\documentclass[11pt]{article}
\usepackage[a4paper,margin=1in]{geometry}
\usepackage[T1]{fontenc}
\usepackage{lmodern}
\usepackage{microtype}
\usepackage{amsmath,amssymb,amsthm,mathtools}
\usepackage{bm}
\usepackage{graphicx}
\usepackage{booktabs}
\usepackage{enumitem}
\usepackage{float}
\usepackage{tabularx}
\usepackage{makecell}
\usepackage{array}
\usepackage{hyperref}
\usepackage[nameinlink,noabbrev]{cleveref}

\hypersetup{
  colorlinks=true,
  linkcolor=blue,
  citecolor=blue,
  urlcolor=blue
}

\newcolumntype{L}{>{\raggedright\arraybackslash}X}

\newtheorem{theorem}{Theorem}
\newtheorem{corollary}[theorem]{Corollary}
\newtheorem{remark}{Remark}

\newcommand{\draftnote}[1]{%
  \par\medskip\noindent\textbf{Drafting note.} #1\par\medskip
}

\title{A Conditional Full Conservative Derivation of the Two-Body $4$PN Sector from the Unified $4$D Toy Model}
\author{Trevor Norris}
\date{\today}

\begin{document}
\maketitle

\begin{abstract}
We present the conditional full conservative two-body $4$PN assembly for the $4$D toy-model program within the same declared closure hierarchy used at $1$PN, $2$PN, and $3$PN. Starting from the exact $4$D reduction backbone, the frozen lower-order conservative ledger, and the $2.5$PN narrowing to the canonically normalized orbital/worldtube STF quadrupole branch, we close the complete \emph{local instantaneous} near-zone ledger through order $c^{-8}$ and isolate the hereditary contribution in the form
\[
L_{4\mathrm{PN}}^{\mathrm{cons}} = L_{4\mathrm{PN}}^{\mathrm{local}} + L_{4\mathrm{PN}}^{\mathrm{tail}},
\qquad
C_{\rm tail}=\frac{GM}{2c^3}\,\gamma_{\rm quad}^{\rm eff}.
\]
The local theorem chain uses the exact Schwarzschild one-body $4$PN gate, the exact quartic Legendre compiler, the Hamiltonian-first generic-frame lift of the fixed local target, the exact canonical residual slice, the generic-frame lift of the aligned-seed ordinary correction, and the exact local referee reconstruction. The only conditional step left for a fully normalized conservative $4$PN theorem is the same passive/outgoing STF quadrupole normalization already isolated by the $2.5$PN program,
\[
\gamma_{\rm quad}^{\rm eff}\stackrel{?}{=}\frac{2G}{5c^5}.
\]
The paper is written as a theorem-status ledger: exact, reduced, and conditional
claims are separated explicitly so that the remaining $4$PN gate is isolated
rather than blurred into the solved local closure.
\end{abstract}

% ---------------------------------------------------------------------------
% Main text
% ---------------------------------------------------------------------------

% Section: Introduction and scope for the conditional 4PN paper

\section{Introduction and scope}
\label{sec:intro}

\subsection{Motivation: from full conservative \texorpdfstring{$3$}{3}PN and the \texorpdfstring{$2.5$}{2.5}PN narrowing to the local-plus-tail \texorpdfstring{$4$}{4}PN problem}
\label{sec:intro-motivation}

The earlier papers in this program already froze a substantial conservative and
reduction-theoretic backbone.  On the parent-theory side, the unified $4$D
framework supplies the exact projection map, the exact projected continuity law
with leakage, the exact longitudinal identity, the controlled quasi-static
Poisson hook, and the coherent-defect / worldtube reduction used to define the
brane-facing point-particle sector.  On the post-Newtonian side, the lower-order
ledger already carries the Newtonian, conservative $1$PN, conservative $2$PN,
and conservative $3$PN sectors within one declared closure hierarchy.  In
particular, the carried viability conditions remain
\[
\kappa_\rho = 1,
\qquad
n = 5,
\qquad
\kappa_{\rm add} = \frac{1}{2},
\qquad
\kappa_{\rm PV} = \frac{3}{2},
\qquad
\beta_{\rm 1PN} = 3.
\]
Just as importantly, the conservative $2$PN and $3$PN steps already sharpened
the particle/throat ontology well beyond a pure scalar monopole picture.  The
carried structure now includes an odd dipole wake sector, a real
$P_0\oplus P_2$ mouth/support layer, and a separate geometry-closure channel,
with the conservative $3$PN middle block already assigned in the fixed ADM chart
inside the same hierarchy.

The $2.5$PN program then made the higher-order bridge much narrower.  Its main
structural lesson was that the surviving universal point-particle retarded route
is the canonically normalized orbital/worldtube STF quadrupole branch.  In the
conservative bookkeeping used by the present series, that branch is the same
higher-order lane that had already begun to appear as the grouped real
\[
P_{20}\oplus P_{21}\oplus P_{22}
\]
bundle.  This narrowing changes the meaning of the next conservative step.
After $3$PN, the right question is no longer whether one can continue adding
local coefficients somehow.  The right question is whether the first
conservative order that must be stated explicitly as a \emph{local instantaneous
sector plus a hereditary/tail sector} can still be closed inside the same
hierarchy, and whether that hereditary sector opens a genuinely new datum or
only reuses the already-known quadrupole-normalization gate.

That is the purpose of the present paper.  At $4$PN the conservative near-zone
ledger through order $c^{-8}$ must be written in the form
\[
L_{\rm cons}
=
L_0+c^{-2}L_1+c^{-4}L_2+c^{-6}L_3+c^{-8}L_4+O(c^{-10}),
\]
with the theorem split
\[
L_{4\mathrm{PN}}^{\mathrm{cons}}
=
L_{4\mathrm{PN}}^{\mathrm{local}}
+
L_{4\mathrm{PN}}^{\mathrm{tail}}.
\]
The local side must close the strict one-body $4$PN Schwarzschild gate, the
quartic ordinary/Hamiltonian compiler, the full local instantaneous comparable-mass
generic-frame sector, and the fixed-chart uniqueness problem.  The hereditary
side must determine whether the $4$PN tail coefficient is a new independent
normalization datum or is instead controlled by the same STF quadrupole branch
that already survived the $2.5$PN audit.

A first structural lesson appears already at the strict one-body level.  The
exact isotropic Schwarzschild gate through $4$PN is \emph{not} reproduced by a
single quartic continuation of the carried $3$PN denominator-style self sector.
The minimal one-body repair ledger is instead
\[
\mu_{\rho,4}=\frac18,
\qquad
 d_4=\frac{205}{16},
\qquad
 s_{34}=-\frac{15}{32},
\qquad
 s_{26}=-\frac1{16},
\]
so the $4$PN one-body problem is already not a one-parameter extension of the
$3$PN closure.  But the one-body sector is still only the front door.  The
deeper algebraic question is whether the full local comparable-mass $c^{-8}$
residual can be assigned exactly once the lower-order ledger is frozen, the
fixed local target is imported, and the generic-frame chart is handled in the
correct order.

The main positive result of the paper is that it can.  The raw exchange-symmetric
local generic-frame scaffold is already COM-complete blockwise; the decisive
chart lesson is that the lift must be done Hamiltonian-first, because the fixed
local Hamiltonian target respects the constant-coefficient scaffold degree
ceilings while the translated ordinary target develops extra $\nu^4$ tails in
the middle and upper interaction blocks.  In the Hamiltonian chart there is no
span obstruction, the canonical comparable-mass local residual translates back
to the ordinary chart by exact sign flip once the lower-order ledger and aligned
seed convention are frozen, the aligned-seed ordinary correction admits an exact
generic-frame lift after the structured seed enlargement recorded in the audit
chain, and the full local instantaneous $4$PN ordinary sector therefore has an
exact generic-frame representative.

The hereditary side is narrower still.  The exact bridge derived in the notes is
\[
C_{\rm tail}=\frac{GM}{2c^3}\,\gamma_{\rm quad}^{\rm eff},
\]
so the $4$PN tail introduces no new independent normalization datum.  Full GR
matching of the hereditary coefficient is equivalent to the already-known
$2.5$PN normalization target
\[
\gamma_{\rm quad}^{\rm eff}\stackrel{?}{=}\frac{2G}{5c^5}.
\]
Thus the remaining full-$4$PN conditional step is not another open-ended
coefficient search.  It is the same passive/outgoing STF quadrupole
normalization theorem that the $2.5$PN program had already isolated.

The scope of the paper is therefore deliberately strong but intermediate.  It is
stronger than a notebook-level or target-matching statement because the entire
local instantaneous $4$PN sector is already assembled exactly and the hereditary
interface has been reduced to one named normalization gate.  It is weaker than a
fully first-principles moving-throat PDE theorem because the passive/outgoing
quadrupole normalization still awaits a completed branch derivation, and the
microscopic origin of the grouped real conservative quadrupole data remains a
deeper dynamical problem.  The correct reading is the same as in the earlier
full conservative $1$PN, $2$PN, and $3$PN papers: a full assembly theorem
\emph{within a declared closure hierarchy}, not an assumption-free theorem of
the fully solved bulk problem.

A second purpose of the paper is bookkeeping clarity.  At $4$PN the argument
mixes carried exact lower-order structure, controlled one-body and COM
reductions, fixed-chart import, protocol-level local seed choices, and one
conditional hereditary interface that is inherited rather than newly opened.
The reader should be able to see, line by line, which ingredients are exact,
which depend on the declared reduction and chart choices, which live only inside
the same response hierarchy that already succeeded below $4$PN, and which are
final theorem consequences.  That separation is part of the result.

\subsection{What this paper claims}
\label{sec:intro-claims}

The main claims of the paper are the following.
\begin{enumerate}
  \item The strict one-body $4$PN Schwarzschild gate is fixed exactly, and its
  minimal repair ledger already shows that $4$PN is not a one-parameter
  continuation of the carried $3$PN one-body closure.

  \item The exact quartic perturbative Legendre-transform identity closes the
  ordinary/Hamiltonian bridge through $4$PN and provides the unique compiler
  needed for any trustworthy local comparable-mass lift.

  \item The raw exchange-symmetric generic-frame local scaffold is already
  COM-complete blockwise, so local $4$PN never had an early span shortage at the
  level of kinematic image size.

  \item The fixed local $4$PN target must be lifted Hamiltonian-first, because
  the Hamiltonian degree ceilings match the constant-coefficient scaffold while
  the translated ordinary target develops extra $\nu^4$ tails.

  \item In the Hamiltonian chart there is no span obstruction: the generic-frame
  local interaction residual is fully reachable block by block, and one exact
  canonical comparable-mass local residual representative exists.

  \item Once the lower-order ledger and aligned-seed convention are frozen, the
  canonical local residual translates back to the ordinary chart by exact sign
  flip,
  \[
  \Delta H_4=-\Delta L_4,
  \]
  and the aligned-seed ordinary correction admits an exact structured
  generic-frame lift.

  \item The entire local instantaneous $4$PN ordinary sector therefore has an
  exact generic-frame representative and can be reconstructed referee-style in
  the fixed local chart.

  \item The hereditary $4$PN coefficient is not a new independent datum; it is
  fixed by the same canonically normalized STF quadrupole branch already
  isolated by the $2.5$PN program, namely
  \[
  C_{\rm tail}=\frac{GM}{2c^3}\,\gamma_{\rm quad}^{\rm eff}.
  \]

  \item Therefore the remaining $4$PN gap is exactly the same
  passive/outgoing quadrupole normalization problem that already survived the
  $2.5$PN audit.  In particular, the local $4$PN bottleneck is no longer an
  algebraic coefficient search.
\end{enumerate}

Stated differently, the point of the paper is not that the model can be tuned
until a $4$PN-looking ledger appears.  The point is that, once the lower-order
stack, the fixed local target, and the Hamiltonian-first chart logic are handled
correctly, the full conservative $c^{-8}$ problem narrows to one closed local
theorem chain plus one already-named hereditary normalization gate.

\subsection{What this paper does not claim}
\label{sec:intro-nonclaims}

The negative side of the scope is equally important.

First, this paper does \emph{not} claim an assumption-free theorem that begins
with a fully solved moving-throat PDE and ends automatically with the complete
conservative two-body $4$PN dynamics.  As in the earlier full conservative
papers, the point-particle limit used here is a controlled reduction inside a
declared hierarchy.  The main text therefore presents a reduced-theory theorem
chain, not a finished theorem of the full moving-throat bulk problem.

Second, this paper does \emph{not} claim that the passive/outgoing quadrupole
normalization has already been derived microscopically from the completed throat
PDE.  The hereditary bridge proved here is that the $4$PN tail coefficient is
controlled by the same canonically normalized STF quadrupole branch as the
$2.5$PN odd channel.  The actual branch normalization on the true moving-throat
solution remains outside the present theorem.

Third, this paper does \emph{not} claim a final uniqueness or interpretation
theorem for every aligned-seed local representative in the ordinary chart.  What
is established here is that the local instantaneous ledger is closed exactly
inside the declared hierarchy and fixed chart once the aligned-seed convention is
handled by the structured generic-frame lift.  The deeper microscopic meaning of
that aligned-seed sector is a separate follow-on problem.

Fourth, the paper does \emph{not} claim any new dissipative or radiative theorem
beyond the already narrowed quadrupole route.  The scalar and dipole dangers had
already been demoted above universal point-particle $2.5$PN on the natural
small-body branch, and the present paper does not reopen those channels.

Fifth, the paper is restricted to the nonspinning conservative two-body sector
through order $c^{-8}$.  It does not attempt a treatment of spin couplings,
strong-field nonperturbative completion, generic many-body response, or higher-PN
sectors beyond the present conservative $4$PN scope.  Nor does it claim to
replace the broader moving-throat script-side program by a single closed PDE
paper.

These non-claims are not qualifications to hide.  They are part of the correct
interpretation of the result.  The paper should be read as establishing the
sharpest honest conservative $4$PN statement currently supported by the program:
a full local-plus-tail assembly theorem within the declared hierarchy, with the
remaining open work shifted from local algebra to microscopic quadrupole
normalization on the true passive/outgoing branch.

\subsection{Claim taxonomy}
\label{sec:intro-taxonomy}

To keep the argument referee-safe, we retain the same bookkeeping discipline
used in the earlier full conservative papers.
\begin{description}
  \item[Exact.] Identities that follow directly from the frozen lower-order
  ledger, the exact one-body $4$PN target, the exact quartic compiler, the exact
  chart-translation theorems, symmetry bookkeeping, or the exact hereditary
  coefficient bridge once the target chart and basis choices have been fixed.

  \item[Controlled reduction.] Steps that require an explicit reduced
  description, such as the carried $4$D-to-brane and worldtube reductions, the
  strict one-body and COM reductions, the fixed local Hamiltonian import, or the
  grouped real quadrupole reduction inherited from the $2.5$PN narrowing.

  \item[Protocol closure.] Statements that depend on the same declared response
  hierarchy used successfully below $4$PN rather than on a fully solved
  moving-throat PDE.  In the present paper the most important examples are the
  carried self/static ledger, the Hamiltonian-aligned seed convention for the
  local ordinary chart, and the use of the already-isolated passive/outgoing STF
  quadrupole normalization as the hereditary interface datum.

  \item[Full assembly.] The final theorem consequence after the previous
  ingredients are inserted.  In the present paper this means that, once the
  lower-order ledger, the fixed local ADM target, the quartic compiler, the
  Hamiltonian-first lift, and the hereditary bridge are frozen, the final
  equality
  \[
  L_{4\mathrm{PN}}^{\mathrm{cons}}
  =
  L_{4\mathrm{PN}}^{\mathrm{local}}
  +
  L_{4\mathrm{PN}}^{\mathrm{tail}}
  \]
  is exact algebra inside the declared hierarchy.
\end{description}

This taxonomy is not cosmetic.  It prevents four logically distinct moves from
being blurred together: carrying forward exact lower-order structure, using the
declared reductions and chart choices, invoking the closure needed to define the
local aligned-seed sector and the hereditary interface, and then stating the
final theorem consequence that survives after those ingredients are frozen.

\subsection{Roadmap}
\label{sec:intro-roadmap}

The next section gives an executive summary of the final theorem ledger.
Section~\ref{sec:dictionary} then records the minimal carry-forward dictionary
from the earlier $4$D, $1$PN, $2$PN, $2.5$PN, and $3$PN papers and states the
exact $4$PN target.  Section~\ref{sec:onebody} begins the proof chain with the
strict one-body Schwarzschild gate, the minimal repair ledger, the natural
self/static seed, and the exact quartic Legendre compiler.

Section~\ref{sec:local-chain} then closes the local instantaneous theorem chain:
the raw local scaffold, the fixed local target import, the Hamiltonian-chart
generic-frame lift, the canonical comparable-mass residual slice, the ordinary
translation, the aligned-seed correction and its generic-frame lift, and the
exact local referee reconstruction.  Section~\ref{sec:tail-bridge} isolates the
hereditary bridge, proves that the tail coefficient is controlled by the same
canonically normalized quadrupole branch already isolated at $2.5$PN, and shows
that $4$PN opens no new independent normalization datum.

Section~\ref{sec:final-conditional-4pn-theorem} states the final conditional conservative $4$PN
theorem in COM and generic-frame form.  Section~\ref{sec:interface-25pn-pde}
then explains exactly how the present result interfaces with the $2.5$PN program
and with the moving-throat PDE: what is fixed now, what remains branch
normalization, and why the correct next theorem gate is the passive/outgoing STF
quadrupole normalization on the true branch.  Section~\ref{sec:fixed-open}
separates frozen inputs, newly fixed quantities, and still-open objects;
Section~\ref{sec:verification} records the role and limits of the referee
archive; and Section~\ref{sec:conclusions} closes the paper.

The appendices collect the detailed one-body algebra, the quartic compiler, the
local scaffold and target-import audits, the Hamiltonian-first lift, the
aligned-seed correction and its structured generic-frame lift, the exact local
referee reconstruction, the hereditary bridge algebra, and the verification
ledger.  They are essential to the theorem chain but would obscure the main
logical structure if left in the body of the paper.

% Section: Executive summary of results for the conditional 4PN paper

\section{Executive summary of results}
\label{sec:executive-summary}

This paper closes the first conservative order of the program that must be
stated explicitly as a sum of a full local instantaneous sector and a separate
hereditary/tail sector.  The earlier papers already froze the exact $4$D
reduction backbone, carried the conservative Newtonian+$1$PN+$2$PN+$3$PN ledger
through its declared closure hierarchy, and narrowed the live higher-order
universal point-particle bridge to the canonically normalized orbital/worldtube
STF quadrupole branch.  The present question is therefore not whether one can
continue adding local $4$PN coefficients somehow.  It is whether the full
nonspinning conservative two-body ledger through order $c^{-8}$ can be assigned
uniquely once the strict one-body gate, the exact quartic
ordinary/Hamiltonian compiler, the fixed local target, the generic-frame lift,
and the hereditary interface are all handled in the correct chart.  The answer
is now sharp: within the same declared closure hierarchy used successfully at
$1$PN, $2$PN, and $3$PN, the entire local instantaneous $4$PN sector is already
assembled exactly, and the only remaining conditional interface for a full
conservative $4$PN theorem is the same passive/outgoing STF quadrupole
normalization already isolated by the $2.5$PN program.

\paragraph{Headline result.}
The final theorem split is
\[
L_{4\mathrm{PN}}^{\mathrm{cons}}
=
\underbrace{L_{4\mathrm{PN}}^{\mathrm{local}}}_{\text{closed exactly}}
+
\underbrace{L_{4\mathrm{PN}}^{\mathrm{tail}}}_{\text{fixed by the STF quadrupole bridge}},
\]
with
\[
C_{\rm tail}=\frac{GM}{2c^3}\,\gamma_{\rm quad}^{\rm eff}.
\]
Thus $4$PN opens no new independent normalization datum beyond the already-known
quadrupole gap.  Full GR recovery of the hereditary coefficient is equivalent to
the single condition
\[
\gamma_{\rm quad}^{\rm eff}=\frac{2G}{5c^5},
\]
in which case
\[
C_{\rm tail}
=
\frac{GM}{2c^3}\,\gamma_{\rm quad}^{\rm eff}
=
\frac{G^2M}{5c^8}
=
C_{\rm tail}^{\rm GR}.
\]
What remains open is therefore no longer a local $4$PN coefficient search.
What remains open is the passive/outgoing STF quadrupole normalization on the
true moving-throat branch.

Table~\ref{tab:executive-ledger-4pn} lists the load-bearing results established
in the present paper and the role each plays in the theorem chain.

\begin{table}[t]
  \centering
  \caption{Load-bearing results of the present $4$PN paper.}
  \label{tab:executive-ledger-4pn}
  \footnotesize
  \setlength{\tabcolsep}{4pt}
  \renewcommand{\arraystretch}{1.12}
  \begin{tabularx}{\linewidth}{@{}p{0.22\linewidth}L p{0.15\linewidth}@{}}
    \toprule
    Item & Result & Status \\
    \midrule
    One-body gate
    & The strict isotropic Schwarzschild target through $c^{-8}$ is fixed
      exactly, and the minimal repair ledger is
      $\mu_{\rho,4}=1/8$, $d_4=205/16$, $s_{34}=-15/32$, $s_{26}=-1/16$.
    & Settled \\

    Quartic compiler
    & The exact quartic perturbative Legendre-transform identity closes the
      ordinary/Hamiltonian bridge through $4$PN and fixes the unique local
      compiler used by the Hamiltonian-first lift.
    & Settled \\

    Raw local scaffold
    & The exchange-symmetric generic-frame local interaction scaffold has block
      sizes $52$, $46$, $29$, $10$, and $2$, and its COM image already spans the
      full $15$ local interaction slots blockwise.
    & Settled \\

    Hamiltonian-chart lift
    & The local Hamiltonian target respects the scaffold degree ceilings
      $(4,3,3,2,2)$, the blockwise coefficient-space ranks are maximal
      $20/20$, $12/12$, $9/9$, $4/4$, and $2/2$, and no Hamiltonian-chart span
      obstruction survives.
    & Exact closure \\

    Canonical quotient slice
    & The $92$-dimensional Hamiltonian-chart null family is pure COM-blind
      freedom.  Fixing the canonical slice gives one sparse exact generic-frame
      Hamiltonian residual representative.
    & Exact quotient \\

    Aligned-seed correction
    & The remaining ordinary-chart obstruction is purely a seed-alignment issue,
      and the correction is exactly generic-frame liftable in the structured
      seed spaces $Q$, $T\oplus(pq)T$, $S\oplus(pq)S$, $U\oplus(p^2q^2)U$, with
      $W=0$.
    & Exact lift \\

    Local referee reconstruction
    & The assembled generic-frame ordinary representative reproduces the full
      fixed-chart local ordinary $4$PN target block by block under COM
      reduction.
    & Exact closure \\

    Tail bridge
    & The GR tail has the exact Newtonian quadrupole shadow in the degree-$2$
      $G^4/r^4$ block, and the hereditary coefficient satisfies
      $C_{\rm tail}=(GM/2c^3)\,\gamma_{\rm quad}^{\rm eff}$.
    & Exact bridge \\

    Remaining gap
    & The only conditional step left for a full conservative $4$PN theorem is
      the passive/outgoing STF quadrupole normalization on the actual
      moving-throat branch.
    & Open \\
    \bottomrule
  \end{tabularx}
\end{table}

\paragraph{How the ledger closes.}
The proof chain has four positive gates and one decisive bookkeeping lesson.
First, the strict one-body Schwarzschild target shows that $4$PN is not a
one-parameter continuation of the carried $3$PN self sector.  The quartic static
gate, quartic denominator datum, and the two new self-sector slots in the
$U^3v^4$ and $U^2v^6$ channels are all forced already at test-mass level.
Second, the exact quartic Legendre compiler turns the local comparable-mass
problem into a controlled ordinary/Hamiltonian bridge rather than a blind
coefficient search.  Third, the raw generic-frame local scaffold is already
blockwise COM-complete, so the real issue is not early image shortage but the
large COM-null sector and the correct choice of quotient chart.  Fourth, once
the local target is imported Hamiltonian-first, the entire comparable-mass local
residual becomes exactly reachable, the canonical residual translates back to
the ordinary chart by exact sign flip, and the aligned-seed ordinary correction
admits its own exact structured generic-frame lift.  The decisive bookkeeping
lesson is therefore that the local $4$PN lift is Hamiltonian-first.  The
translated ordinary target develops extra $\nu^4$ tails in the middle and upper
interaction blocks, whereas the fixed local Hamiltonian target matches the
constant-coefficient scaffold degree ceilings exactly.  Once this is respected,
the local $4$PN bottleneck is no longer algebraic span or chart ambiguity.

\paragraph{One-body gate and quartic compiler.}
The strict isotropic Schwarzschild one-body target is
\[
\frac{L_{4\mathrm{PN},1\text{-body}}}{m}
=
\frac{7}{256}\frac{v^{10}}{c^8}
+\frac{75}{128}\frac{Uv^8}{c^8}
+\frac{59}{16}\frac{U^2v^6}{c^8}
+\frac{203}{32}\frac{U^3v^4}{c^8}
+\frac{31}{32}\frac{U^4v^2}{c^8}
+\frac{1}{16}\frac{U^5}{c^8}.
\]
A direct continuation of the carried $3$PN denominator-style self sector does
not reproduce this answer with only one new quartic datum.  The exact minimal
repair ledger is instead
\[
\mu_{\rho,4}=\frac18,
\qquad
 d_4=\frac{205}{16},
\qquad
 s_{34}=-\frac{15}{32},
\qquad
 s_{26}=-\frac1{16},
\]
so the one-body sector already proves that $4$PN is not merely the next scalar
continuation of the earlier closure.  Applying the exact quartic Legendre
compiler then gives the strict one-body Hamiltonian gate
\[
\frac{7}{256}p^{10},
\qquad
\frac{45}{128}\frac{p^8}{r},
\qquad
\frac{13}{8}\frac{p^6}{r^2},
\qquad
\frac{105}{32}\frac{p^4}{r^3},
\qquad
\frac{105}{32}\frac{p^2}{r^4},
\qquad
-\frac{1}{16}\frac{1}{r^5},
\]
which is the exact self/static Hamiltonian seed data used by the later local
Hamiltonian-first lift.

\paragraph{Local scaffold, degree ceilings, and the Hamiltonian-first lesson.}
Symmetric promotion of the exact one-body target gives the natural local
self/static seed
\[
L_{4,\mathrm{seed}}=
\frac{7}{256}(m_A v_A^{10}+m_B v_B^{10})
+\frac{75Gm_A m_B}{128r}(v_A^8+v_B^8)
+\cdots
+\frac{G^5m_A m_B}{16r^5}(m_B^4+m_A^4),
\]
with the whole expression understood as the coefficient of $c^{-8}$.  The
comparable-mass local interaction scaffold built on top of this seed has raw
block sizes
\[
G/r:52,
\qquad
G^2/r^2:46,
\qquad
G^3/r^3:29,
\qquad
G^4/r^4:10,
\qquad
G^5/r^5:2,
\]
for a total of $139$ exchange-symmetric generic-frame interaction directions.
After COM projection these span exactly the full $15$ local interaction slots,
namely
\[
G/r:5,
\qquad
G^2/r^2:4,
\qquad
G^3/r^3:3,
\qquad
G^4/r^4:2,
\qquad
G^5/r^5:1.
\]
So local $4$PN never had an early kinematic shortage.  The real issue is chart
choice.  In the Hamiltonian chart the imported local interaction blocks have
$\nu$-degree ceilings
\[
(4,3,3,2,2),
\]
whereas after quartic translation the ordinary target becomes
\[
(4,4,4,4,2).
\]
The extra ordinary-chart $\nu^4$ tails in the middle and upper blocks are
therefore a translation artifact rather than evidence of a true local span
failure.  This is why the correct local $4$PN route is
\[
\text{fixed local Hamiltonian target}
\longrightarrow
\text{generic-frame Hamiltonian lift}
\longrightarrow
\text{ordinary translation only afterward}.
\]

\paragraph{Hamiltonian-chart completeness and exact ordinary reconstruction.}
Once the fixed local target is treated in the Hamiltonian chart, the blockwise
coefficient-space ranks are maximal:
\[
G/r:20/20,
\qquad
G^2/r^2:12/12,
\qquad
G^3/r^3:9/9,
\qquad
G^4/r^4:4/4,
\qquad
G^5/r^5:2/2.
\]
So the full interaction coefficient-space dimensions are
\[
\text{rank}=47,
\qquad
\text{nullity}=92,
\]
and there is no Hamiltonian-chart span obstruction.  The null family is purely
COM-blind algebraic freedom.  Fixing the canonical quotient slice by setting all
null coordinates to zero gives one sparse exact generic-frame Hamiltonian
representative.  Translating that representative back to the ordinary chart
exposes the only genuine local obstruction left, namely the seed-alignment
correction
\[
\delta L_{4,\rm seed}
=
L_{4,\rm seed}^{(\rm aligned)}-L_{4,\rm seed}^{(\rm natural)}.
\]
That correction is not a failure of the generic-frame lift.  It is an aligned-seed
issue, and it is exactly generic-frame liftable in the structured seed spaces
\[
Q,
\qquad
T\oplus(pq)T,
\qquad
S\oplus(pq)S,
\qquad
U\oplus(p^2q^2)U,
\qquad
W=0.
\]
Once the natural generic-frame seed, the structured aligned-seed correction, and
the canonical translated local residual are assembled together, the local
referee master verifies that the COM reduction of the resulting generic-frame
ordinary candidate reproduces the full fixed-chart local ordinary $4$PN target
block by block.  The entire local instantaneous $4$PN ordinary sector is
therefore closed exactly.

\paragraph{Tail bridge and the no-new-datum theorem.}
The standard conservative GR tail coefficient is
\[
C_{\rm tail}^{\rm GR}=\frac{G^2M}{5c^8},
\]
and the corresponding time-symmetric nonlocal Hamiltonian has local logarithmic
shadow
\[
F_{\rm tail}(t)=\frac{2}{5}\frac{G^2M}{c^8}\bigl(I_{ij}^{(3)}(t)\bigr)^2.
\]
Using the Newtonian STF quadrupole
\[
I_{ij}=\mu\,\mathrm{STF}(x_i x_j),
\]
one finds the exact order-reduced shadow
\[
\frac{F_{\rm tail}}{\mu}
=
\frac{16}{15}\,\nu\,\frac{U^4}{c^8}
\left(12v^2-11\dot r^2\right),
\qquad
U\equiv\frac{GM}{r},
\]
which occupies only the degree-$2$ $G^4/r^4$ block.  The decisive hereditary
identity is then
\[
C_{\rm tail}=\frac{GM}{2c^3}\,\gamma_{\rm quad}^{\rm eff}.
\]
This is the exact bridge to the already isolated $2.5$PN quadrupole route.  In
particular, no new independent $4$PN normalization datum opens.  The same
canonically normalized STF quadrupole branch that controls the $2.5$PN
Burke--Thorne channel also controls the $4$PN hereditary coefficient.  On the
Burke--Thorne target branch,
\[
\gamma_{\rm quad}^{\rm eff}=\frac{2G}{5c^5},
\]
one recovers
\[
C_{\rm tail}^{\rm toy}=\frac{G^2M}{5c^8}=C_{\rm tail}^{\rm GR}.
\]

\paragraph{What the result fixes for the quadrupole program.}
The main positive interface result is that the local conservative $4$PN problem
is no longer waiting on another comparable-mass algebraic solve.  The exact
one-body gate, quartic compiler, raw local scaffold, Hamiltonian-first lift,
canonical quotient slice, aligned-seed correction, and local referee
reconstruction together remove the local uncertainty completely.  The remaining
open work is therefore much narrower than ``derive $4$PN somehow.''  What the
completed moving-throat PDE still has to supply is the passive/outgoing STF
quadrupole normalization on the true branch, equivalently the canonical
low-frequency pair $(\overline K_0,\overline K_2)$ whose isotropic value fixes
$\gamma_{\rm quad}^{\rm eff}$.  Once that pair is determined on the natural
branch, the normalized $2.5$PN odd quadrupole coefficient closes, the normalized
$4$PN hereditary coefficient closes, and the present conditional $4$PN theorem
becomes unconditional inside the same hierarchy.

% Section: Carry-forward dictionary and theorem target for the conditional 4PN paper

\section{Carry-forward dictionary and theorem target}
\label{sec:dictionary}

This section is intentionally downstream of the parent $4$D action paper, the
full conservative $1$PN, $2$PN, and $3$PN assembly papers, and the $2.5$PN
quadrupole-normalization audit.  We do not repeat those longer derivations
here.  Instead, we record only the imported objects, frozen lower-order values,
chart conventions, and exact target statements that the present $4$PN proof
chain actually uses.  The purpose of the section is fourfold: to freeze
notation, to separate exact imports from controlled reductions, to make clear
what is being carried rather than rederived, and to state the theorem target
that organizes the sections that follow.

The guiding principle is the same one used in the earlier conservative papers.
What is exact in the parent reduction framework should remain visibly exact;
what enters only after a declared near-zone, small-body, COM, or grouped-branch
assumption should remain visibly reduced; and what is fixed only inside the
response hierarchy should remain visibly hierarchical rather than disguised as
an assumption-free theorem of the fully solved moving-throat PDE.  The present
paper is therefore built on an imported dictionary rather than on a fresh
rederivation of the full bulk problem.

\subsection{Minimal imported \texorpdfstring{$4$}{4}D reduction dictionary}
\label{sec:dictionary-4d}

The parent theory lives on a $(4+1)$-dimensional spacetime with coordinates
\[
X^M=(t,x,y,z,w),
\qquad
\mathbf{X}=(x,y,z,w)\in\mathbb R^4,
\qquad
\mathbf{x}=(x,y,z)\in\mathbb R^3.
\]
The fundamental bulk objects carried into the present paper are the matter field
$\psi(\mathbf{X},t)$, its density
\[
\rho(\mathbf{X},t)=|\psi(\mathbf{X},t)|^2,
\]
the optional localized gauge field $A_M(\mathbf{X},t)$ with localization profile
$Z(w)$, and the effective throat geometry variables
\[
a(t),\qquad L(t).
\]
When electric charge is named explicitly in the inherited $4$D notation, it is
always the corrected branch quantity
\[
q_* = \eta_Q e_*,
\qquad
q_{\rm eff}=\frac{q_*}{\sqrt{Z_{\rm int}}},
\qquad
Z_{\rm int}=\int_{-\infty}^{\infty} Z(w)\,dw,
\]
not the gravity-side scalar coefficient $\kappa_\rho$ that appears in the
post-Newtonian ledger below.

A brane observer does not see the full bulk fields directly.  Brane observables
are introduced by a fixed normalized projection kernel $W(w)$,
\[
\int_{-\infty}^{\infty} W(w)\,dw = 1,
\]
through the projection map
\[
\mathcal P_W[Q](t,\mathbf{x})
\equiv
\int_{-\infty}^{\infty} W(w)\,Q(t,\mathbf{x},w)\,dw.
\]
In particular,
\[
\rho_{\rm br}=\mathcal P_W[\rho],
\qquad
\mathbf J_{\rm br}=\mathcal P_W[(j^x,j^y,j^z)],
\qquad
\mathbf v_{\rm br}=\frac{\mathbf J_{\rm br}}{\rho_{\rm br}}
\quad (\rho_{\rm br}>0).
\]
This projection operation is distinct from \emph{reduction}.  Projection defines
what a brane observer measures; reduction refers to the further controlled
elimination of the extra coordinate under an explicit ansatz or scaling regime.

The exact projected continuity law carried forward from the parent theory is
\[
\partial_t\rho_{\rm br}+\nabla_3\!\cdot\mathbf J_{\rm br}=S_{\rm leak},
\]
with leakage source
\[
S_{\rm leak}
=
-\bigl[W(w)j^w\bigr]_{-\infty}^{+\infty}
+
\int_{-\infty}^{\infty} W'(w)j^w\,dw.
\]
After Helmholtz decomposition of the brane velocity,
\[
\mathbf v_{\rm br}=\nabla_3\phi_{\rm br}+\mathbf v_T,
\qquad
\nabla_3\!\cdot\mathbf v_T=0,
\]
one obtains the exact longitudinal identity
\[
\rho_{\rm br}\,\nabla_3^2\phi_{\rm br}
=
S_{\rm leak}
-\partial_t\rho_{\rm br}
-(\nabla_3\rho_{\rm br})\cdot(\nabla_3\phi_{\rm br}+\mathbf v_T).
\]
Under the carried quasi-static near-zone assumptions this reduces to the Poisson
hook
\[
\nabla_3^2\phi_{\rm br}\simeq \frac{S_{\rm eff}}{\rho_0},
\qquad
\Phi_N=\lambda_\Phi\,\phi_{\rm br},
\]
which is the origin of the Newtonian scalar sector used throughout the lower-order
particle reduction.

The second imported reduction step is the controlled worldtube/coherent-defect
limit.  For a compact defect of size $a_{\rm th}\ll r$ in a smooth external field,
the center-of-mass reduction gives the usual point-particle force law
\[
M_A\,\ddot{\mathbf X}_A
=
-M_A\nabla\Phi_{\rm ext}(\mathbf X_A)
+O\!\left(\frac{a_{\rm th}^2}{r^2}\,\frac{\Phi_{\rm ext}}{r}\right),
\]
so that the universal point-particle source at leading order is controlled by
the worldline/worldtube multipoles rather than by arbitrary internal details of
the defect.  That worldtube dictionary is what later allows the higher-order
universal bridge to be packaged in terms of the orbital/worldtube STF
quadrupole rather than a generic internal support excitation.

\subsection{Frozen conservative carry-forward ledger}
\label{sec:dictionary-carry-forward}

The present paper imports the conservative Newtonian+$1$PN+$2$PN+$3$PN ledger as
frozen input, not as open fit parameters.  The surviving lower-order viability
conditions are
\[
\kappa_\rho = 1,
\qquad
n = 5,
\qquad
\kappa_{\rm add}=\frac12,
\qquad
\kappa_{\rm PV}=\frac32,
\qquad
\beta_{\rm 1PN}=3.
\]
The first four values are the scalar mass-dressing, optical/barotropic,
added-mass, and adiabatic breathing-response coefficients frozen by the earlier
program.  The fifth is the resulting carried $1$PN ledger value.  The present
paper is not allowed to repair a $4$PN mismatch by retuning any of these
numbers.

When referenced, the parity-even wake data inherited from the full conservative
$1$PN assembly are also treated as frozen branch data,
\[
\alpha^2=\frac34,
\qquad
a_H=0,
\qquad
K_{\rm vec}=\frac{2}{\pi^2},
\]
not as live higher-order fit parameters.

The conservative $2$PN and $3$PN assemblies sharpen the ontology of the defect
further.  By the start of the present $4$PN paper, the particle can no longer be
treated as a pure scalar monopole with small corrections.  The carried
conservative structure already looks like a small throat-response theory with
three distinct pieces:
\begin{enumerate}
  \item a carried odd dipole wake sector,
  \item a genuine even $P_0\oplus P_2$ mouth/support layer,
  \item and a separate geometry-closure channel.
\end{enumerate}
Moreover, the full conservative $3$PN paper has already fixed the first honest
conservative grouped-real quadrupole front end in the fixed ADM chart.  So the
present $4$PN work is not starting from a blank EFT and is not free to redefine
the higher-order conservative payload.  It is extending a lower-order ledger in
which the grouped real quadrupole sector is already the live conservative
structure, while the remaining universal dissipative route has already been
narrowed at $2.5$PN to the orbital/worldtube STF quadrupole branch.

What remains genuinely open from the lower-order side is not the existence of
these channels but their full dynamical completion.  The unresolved pole,
compliance, and normalization data are not claimed as already derived from the
fully solved moving-throat PDE.  This is one reason the present paper continues
to speak in terms of a declared closure hierarchy rather than an assumption-free
first-principles theorem.

Two structural consequences of this imported ledger should be stated explicitly.
First, the present paper is not allowed to repair a higher-order mismatch by
changing lower-order conservative constants that were already fixed below $4$PN.
Second, the defect already carries enough internal support structure that the
$4$PN analysis must distinguish carefully between universal point-particle
content, self/static carry-forward content, local comparable-mass lift data, and
the separate hereditary quadrupole interface.

\subsection{Why the \texorpdfstring{$2.5$}{2.5}PN notes control \texorpdfstring{$4$}{4}PN}
\label{sec:dictionary-25pn-controls-4pn}

The $2.5$PN theorem audit is what makes the present $4$PN problem sharply
focused rather than diffuse.  By the end of that audit, the scalar and dipole
odd channels had already been demoted above universal point-particle $2.5$PN on
the strict small-body branch, and the surviving retarded quadrupole route had
been isolated in the isotropic form
\[
\mathcal M_{ab}^R(\omega)
=
\bigl[m_0+m_2\omega^2+m_4\omega^4+i\Gamma_5\omega^5+\cdots\bigr]\delta_{ab},
\qquad
m_0\neq 0,
\quad
\Gamma_5>0.
\]
So $4$PN should be read as the next \emph{conservative} use of a quadrupole
route that $2.5$PN had already isolated, not as the opening of a new unrelated
bridge.

In the branch language inherited from the narrowing, the canonically normalized
STF quadrupole odd coefficient is written
\[
\gamma_{\rm quad}^{\rm eff}
=
\mathcal N_Q\frac{a_{\rm th}^5}{27c_s^5}
=
\overline\Gamma_5
=
9\frac{\overline K_2^{5/2}}{\overline K_0^{3/2}}.
\]
The key point for the present paper is that the hereditary $4$PN coefficient is
controlled by \emph{that same} canonically normalized STF quadrupole branch.
Thus the later hereditary calculation is not a fresh search for a new $4$PN tail
data structure.  It is the conservative reuse of the already narrowed $2.5$PN
quadrupole normalization problem.

This is also why the Burke--Thorne target remains the only conditional interface
that matters:
\[
\gamma_{\rm quad}^{\rm eff}=\frac{2G}{5c^5}.
\]
If this holds on the natural passive/outgoing STF quadrupole branch, then the
$2.5$PN odd coefficient and the $4$PN hereditary coefficient close together.
The present paper therefore treats the $2.5$PN narrowing as a carried theorem
gate, not as optional background motivation.

\subsection{PN bookkeeping and invariant variables}
\label{sec:dictionary-pn-bookkeeping}

The conservative expansion convention is
\[
L_{\rm cons}=L_0+c^{-2}L_1+c^{-4}L_2+c^{-6}L_3+c^{-8}L_4+O(c^{-10}).
\]
So throughout this paper, ``$4$PN'' means the coefficient of $c^{-8}$.

In the generic frame, the basic two-body invariants remain
\[
\mathbf r=\mathbf x_A-\mathbf x_B,
\qquad
r=|\mathbf r|,
\qquad
\mathbf n=\mathbf r/r,
\]
\[
v_A^2=\mathbf v_A^2,
\qquad
v_B^2=\mathbf v_B^2,
\qquad
v_{AB}=\mathbf v_A\!\cdot\!\mathbf v_B,
\]
\[
v_{An}=\mathbf v_A\!\cdot\!\mathbf n,
\qquad
v_{Bn}=\mathbf v_B\!\cdot\!\mathbf n.
\]
The usual local pair potentials are
\[
U_A=\frac{Gm_B}{r},
\qquad
U_B=\frac{Gm_A}{r},
\]
and in the strict one-body gate we write
\[
U=\frac{GM}{r}.
\]

In the COM sector, the reduced ordinary-Lagrangian basis uses the usual
velocity invariants together with the symmetric mass ratio
\[
d\equiv \dot r,
\qquad
\nu\equiv\frac{m_A m_B}{(m_A+m_B)^2},
\qquad
v_\perp^2\equiv v^2-d^2.
\]
The matching COM Hamiltonian basis uses the corresponding reduced momentum
invariants.  In the present paper the ordinary COM coefficients are denoted
$l_i$, the COM Hamiltonian coefficients $h_i$, and the residuals beyond the
chosen seed by $\Delta l_i$ and $\Delta h_i$.

For the raw local generic-frame scaffold it is also convenient to use the
invariant shorthand
\[
a=v_A^2,
\qquad
b=v_B^2,
\qquad
c=\mathbf v_A\!\cdot\!\mathbf v_B,
\qquad
d=\mathbf v_A\!\cdot\!\mathbf n,
\qquad
e=\mathbf v_B\!\cdot\!\mathbf n,
\]
\[
p=m_A,
\qquad
q=m_B.
\]
These shorthands are purely local-scaffold variables.  They are not the throat
radius, electric charge, or grouped quadrupole labels.

\subsection{COM versus generic-frame conventions}
\label{sec:dictionary-com-vs-generic}

The settled $4$PN convention is simple: use the COM basis whenever it
simplifies the exact target or the reduced compiler map, but treat the
\emph{Hamiltonian chart} as the authoritative chart for the local generic-frame
lift.

The reason for this convention is now settled.  After the local/nonlocal split,
the imported local $4$PN Hamiltonian target has one pure free slot and fifteen
interaction slots.  Its strict test-mass limit matches the exact one-body $4$PN
Hamiltonian gate.  More importantly, its interaction blocks have
$\nu$-degree ceilings
\[
(4,3,3,2,2)
\]
across the $G/r$, $G^2/r^2$, $G^3/r^3$, $G^4/r^4$, and $G^5/r^5$ blocks, and
those ceilings match the constant-coefficient generic-frame scaffold exactly.

By contrast, after quartic translation the ordinary local target becomes
\[
(4,4,4,4,2),
\]
so the ordinary chart develops extra $\nu^4$ tails in the $G^2/r^2$, $G^3/r^3$,
and $G^4/r^4$ blocks.  The correct $4$PN sequence is therefore
\[
\begin{aligned}
\text{fixed local Hamiltonian target}
&\longrightarrow \text{generic-frame Hamiltonian lift}
\\
&\longrightarrow \text{ordinary translation only afterward}.
\end{aligned}
\]
The local lift is Hamiltonian-first not because the ordinary chart is
illegitimate, but because the ordinary $\nu^4$ tails are a quartic-translation
artifact rather than evidence of a true local span failure.

A second settled convention follows from the canonical-slice result.  In the
Hamiltonian chart the remaining null family is purely COM-blind algebraic
freedom.  Once the canonical slice and aligned-seed convention are fixed, the
later ordinary representative is no longer a live generic ambiguity.  The fixed
Hamiltonian chart is the authoritative local quotient statement.

\subsection{Notation firewall}
\label{sec:dictionary-firewall}

The notation firewall from the corrected $4$D / EM / PN stack must be kept
intact.

First, electric-charge notation is completely separate from the conservative
$4$PN coefficient ledger.  When charge appears at all, it uses the corrected
variables
\[
\eta_Q=\pm1,
\qquad
q_* = \eta_Q e_*,
\qquad
q_{\rm eff}=\frac{q_*}{\sqrt{Z_{\rm int}}},
\]
while the historical gravity-side bare $q=1$ remains renamed as
\[
\kappa_\rho=1.
\]
Quantized circulation belongs to the magnetic/vortical sector, not to the
electric-charge dictionary.

Second, the grouped labels $20,21,22$ always denote grouped real $P_2$
channels or their canonically normalized STF descendants.  They are not tensor
indices.  Similarly, the later branch variables
\[
\overline K_0,
\qquad
\overline K_2,
\qquad
\overline\Gamma_5,
\qquad
\gamma_{\rm quad}^{\rm eff}
\]
are quadrupole-branch data, not COM slot coefficients.

Third, the generic-frame scaffold symbols
\[
a,b,c,d,e,p,q
\]
are only invariant shorthands for the local two-body algebra.  In particular,
the throat-radius scale entering the quadrupole-branch normalization is always
written explicitly as
\[
a_{\rm th},
\]
so that it cannot be confused with the shorthand
\[
a=v_A^2.
\]
If the throat-response source-weight vector
\[
J=\left(\frac{4}{\sqrt5},\frac54,0,0,0,0\right)
\]
appears later when the grouped branch is discussed, it is a throat-response
source-weight vector rather than an electromagnetic current.

Finally, the three layers of notation used by the paper must remain separate:
COM coefficients $(l_i,h_i)$, grouped conservative/outgoing quadrupole data,
and geometric throat variables $(a,L)$.  Most avoidable confusion at this stage
is not physics but notation drift across those three ledgers.

\subsection{Exact theorem target}
\label{sec:dictionary-theorem-target}

The theorem target of the paper is not an unconstrained coefficient fit.  It has
a strict one-body front door, an imported fixed local target, and a hereditary
bridge whose coefficient is already tied to the narrowed $2.5$PN quadrupole
route.

\paragraph{Strict one-body gate.}
The exact isotropic Schwarzschild target through $c^{-8}$ is
\[
\frac{L_{4\mathrm{PN},1\mathrm{body}}}{m}
=
\frac7{256}v^{{10}}
+\frac{75}{128}Uv^8
+\frac{59}{16}U^2v^6
+\frac{203}{32}U^3v^4
+\frac{31}{32}U^4v^2
+\frac1{16}U^5.
\]
Any acceptable $4$PN one-body/self sector must reproduce this answer exactly.
The first task of the paper is therefore to show that the carried denominator-style
package can be repaired minimally by the finite ledger
\[
\mu_{\rho,4}=\frac18,
\qquad
 d_4=\frac{205}{16},
\qquad
 s_{34}=-\frac{15}{32},
\qquad
 s_{26}=-\frac1{16},
\]
rather than by an uncontrolled extension of the lower-order self sector.

\paragraph{Imported fixed local target.}
The local instantaneous $4$PN target is imported only after the standard
local/nonlocal split has been made.  In reduced COM Hamiltonian form it has one
pure free slot and fifteen interaction slots organized as
\[
1\;\oplus\;5\;\oplus\;4\;\oplus\;3\;\oplus\;2\;\oplus\;1
\]
across the free, $G/r$, $G^2/r^2$, $G^3/r^3$, $G^4/r^4$, and $G^5/r^5$ blocks.
Its strict test-mass limit matches the one-body Hamiltonian gate exactly.
The exact local theorem target is to show that, inside the declared hierarchy,
this imported local Hamiltonian target admits a generic-frame representative,
that its canonical comparable-mass residual translates back to the ordinary
chart by exact sign flip,
\[
\Delta H_4=-\Delta L_4,
\]
and that the remaining ordinary-chart seed-alignment correction is generic-frame
liftable in the structured seed spaces
\[
Q,
\qquad
T\oplus(pq)T,
\qquad
S\oplus(pq)S,
\qquad
U\oplus(p^2q^2)U,
\qquad
W=0.
\]
Equivalently, the local target is to prove the existence of one exact
generic-frame ordinary representative of the form
\[
L_{4,\rm loc}^{\rm gen}
=
L_{4,\rm seed}^{(\rm natural,gen)}
+
\delta L_{4,\rm seed}^{(\rm gen)}
+
\Delta L_{4,\rm loc}^{(\rm can,gen)},
\]
whose COM reduction reproduces the full fixed-chart local ordinary target block
by block.

\paragraph{Hereditary bridge target.}
The remaining nonlocal target is the time-symmetric hereditary coefficient.
In standard GR notation,
\[
C_{\rm tail}^{\rm GR}=\frac{G^2M}{5c^8},
\]
and the corresponding local logarithmic shadow occupies only the degree-$2$
$G^4/r^4$ block.  The exact bridge to the already narrowed quadrupole route is
\[
C_{\rm tail}=\frac{GM}{2c^3}\,\gamma_{\rm quad}^{\rm eff}.
\]
This is the decisive theorem target on the hereditary side.  It says that the
$4$PN tail sector is not introducing a new independent normalization constant.
The only conditional interface is the same passive/outgoing STF quadrupole
normalization already isolated by the $2.5$PN program.  On the Burke--Thorne
branch,
\[
\gamma_{\rm quad}^{\rm eff}=\frac{2G}{5c^5}
\qquad\Longleftrightarrow\qquad
\mathcal N_Q^{\rm target}=\frac{54Gc_s^5}{5a_{\rm th}^5c^5},
\]
and then automatically
\[
C_{\rm tail}=\frac{G^2M}{5c^8}=C_{\rm tail}^{\rm GR}.
\]

\paragraph{Fixed-chart theorem target.}
Within the same declared closure hierarchy used successfully at $1$PN, $2$PN,
and $3$PN, the goal is to show all of the following.
\begin{enumerate}
  \item The exact one-body repair ledger closes the strict Schwarzschild gate.

  \item The imported fixed local Hamiltonian target is fully reachable in the
  generic frame once the Hamiltonian chart is used first.

  \item The canonical comparable-mass local residual translates back to the
  ordinary chart by exact sign flip once the lower-order ledger and aligned seed
  convention are fixed.

  \item The seed-alignment correction is exactly liftable in the structured seed
  sector above, so the earlier ordinary-chart obstruction is revealed to be a
  seed issue rather than a local comparable-mass span failure.

  \item The entire local instantaneous $4$PN sector therefore has an exact
  generic-frame ordinary representative.

  \item The hereditary coefficient is fixed by the same canonically normalized
  STF quadrupole branch already isolated by $2.5$PN, so no new independent
  $4$PN normalization datum opens.
\end{enumerate}
Equivalently, the exact conditional full-assembly target to be proved in the
rest of the paper is
\[
\boxed{
L_{4\mathrm{PN}}^{\mathrm{cons}}
=
\underbrace{L_{4\mathrm{PN}}^{\mathrm{local}}}_{\text{exactly closable in generic-frame ordinary form}}
+
\underbrace{L_{4\mathrm{PN}}^{\mathrm{tail}}}_{\text{fixed by }C_{\rm tail}=(GM/2c^3)\,\gamma_{\rm quad}^{\rm eff}}.
}
\]
This is the most compact statement of the theorem target.  It says that once the
strict one-body gate, the fixed local Hamiltonian target, the Hamiltonian-first
generic-frame lift, the aligned-seed ordinary correction, and the hereditary
coefficient bridge are all handled correctly, the remaining gap between the
present hierarchy and a full conservative $4$PN theorem is no longer ``more
$4$PN coefficients.''  It is the already-known passive/outgoing STF quadrupole
normalization on the true moving-throat branch.

The correct theorem-status reading is the same as in the earlier full
conservative papers.  The boxed statement above is a \emph{full assembly target
within a declared closure hierarchy}, not an assumption-free theorem of the
fully solved moving-throat PDE.  What remains open after the present paper is
therefore not generic local $4$PN algebra but the deeper moving-throat
realization of the passive/outgoing quadrupole normalization, together with its
final insertion into the complete referee/master chain.

% Section: One-body 4PN gate, self/static seed, and exact quartic compiler

\section{One-body \texorpdfstring{$4$}{4}PN gate, self/static seed, and exact quartic compiler}
\label{sec:onebody}

The local $4$PN chain begins by freezing the one-body answer and the unique
order-$4$ ordinary/Hamiltonian compiler before any comparable-mass local lift is
attempted.  This is the same methodological move that organized the earlier
$3$PN paper, but the $4$PN content is sharper.  At this order the one-body
continuation is no longer controlled by a single new denominator datum, and the
quartic Legendre bridge is no longer a dispensable bookkeeping convenience.  It
is a theorem-level input.  The purpose of the present section is therefore to
record four things in one place: the exact Schwarzschild test-mass gate through
$c^{-8}$, the minimal direct-continuation repair ledger relative to the carried
$3$PN self-sector packaging, the natural self/static two-body seed obtained by
symmetric promotion of that gate, and the exact quartic perturbative
Legendre-transform identity that governs every later local $4$PN comparison.

\subsection{Exact Schwarzschild test-mass target through \texorpdfstring{$4$}{4}PN}
\label{sec:onebody-schwarzschild}

The strict test-mass gate is read from the exact isotropic Schwarzschild
worldline Lagrangian
\begin{equation}
\frac{L_{\rm exact}}{m}
=
-c^2\sqrt{\left(\frac{1-u/2}{1+u/2}\right)^2-(1+u/2)^4\frac{v^2}{c^2}},
\qquad
u\equiv \frac{U}{c^2}.
\label{eq:4pn-schwarzschild-exact-worldline}
\end{equation}
Expanding \eqref{eq:4pn-schwarzschild-exact-worldline} through order $c^{-8}$
gives the unique admissible one-body $4$PN coefficient
\begin{equation}
\frac{L_{4\mathrm{PN},1\text{-body}}}{m}
=
\frac{7}{256}\frac{v^{10}}{c^8}
+\frac{75}{128}\frac{Uv^8}{c^8}
+\frac{59}{16}\frac{U^2v^6}{c^8}
+\frac{203}{32}\frac{U^3v^4}{c^8}
+\frac{31}{32}\frac{U^4v^2}{c^8}
+\frac{1}{16}\frac{U^5}{c^8}.
\label{eq:4pn-onebody-lagrangian-gate}
\end{equation}
This coefficient is the exact test-mass answer that any acceptable local
$4$PN self sector must reproduce inside the present hierarchy.  In particular,
it freezes simultaneously the free $v^{10}$ slot, the four mixed
self/static channels
\[
Uv^8,
\qquad
U^2v^6,
\qquad
U^3v^4,
\qquad
U^4v^2,
\]
and the pure static $U^5$ gate.  No later comparable-mass construction is
allowed to modify these test-mass coefficients.  If a proposed local $4$PN
self sector fails \eqref{eq:4pn-onebody-lagrangian-gate}, it is excluded
before any generic-frame or fixed-chart comparison even begins.

This is also the correct place to emphasize the methodological status of the
one-body gate.  The present paper does not infer the $4$PN test-mass answer by
continuing a lower-order pattern heuristically.  It imports the exact
Schwarzschild answer first and only afterward asks what continuation of the
carried $3$PN self-sector packaging is actually compatible with it.  The full
algebraic expansion is recorded in Appendix~\ref{app:onebody-expansion}; the
main text keeps only the coefficient that the later local theorem chain must
match exactly.

\subsection{Minimal one-body repair ledger}
\label{sec:onebody-repair}

To test whether the carried $3$PN self-sector packaging continues directly, we
extend the denominator-style one-body ansatz by writing
\begin{equation}
D(u)=1-4u+8u^2+d_3u^3+d_4u^4,
\label{eq:4pn-onebody-denominator}
\end{equation}
with the already frozen lower-order values
\begin{equation}
d_3=-\frac{45}{4},
\qquad
\mu_{\rho,3}=\frac14,
\qquad
s_{24}=-\frac1{16}.
\label{eq:4pn-onebody-carried-values}
\end{equation}
The most direct continuation of the $3$PN one-body package is then
\begin{equation}
\begin{aligned}
L_{\rm cand}
={}&
-c^2(1-u)\sqrt{1-\frac{v^2/c^2}{D(u)}}
-
\frac{U^2}{2c^2}
+
\frac{U^3}{4c^4}
-
\frac{\mu_{\rho,3}U^4}{2c^6}
+
\frac{\mu_{\rho,4}U^5}{2c^8}
\\[0.5ex]
&
+
 s_{24}\frac{U^2v^4}{c^6}
+
 s_{34}\frac{U^3v^4}{c^8}
+
 s_{26}\frac{U^2v^6}{c^8}.
\end{aligned}
\label{eq:4pn-onebody-direct-continuation}
\end{equation}
Comparing \eqref{eq:4pn-onebody-direct-continuation} with the exact gate
\eqref{eq:4pn-onebody-lagrangian-gate} shows that a quartic denominator plus a
quartic static gate is \emph{not} enough.  Two genuinely new self-sector slots
survive at $4$PN, namely the
\[
\frac{U^3v^4}{c^8}
\qquad\text{and}\qquad
\frac{U^2v^6}{c^8}
\]
channels.  The minimal direct-continuation repair ledger is therefore
\begin{equation}
\boxed{
\mu_{\rho,4}=\frac18,
\qquad
 d_4=\frac{205}{16},
\qquad
 s_{34}=-\frac{15}{32},
\qquad
 s_{26}=-\frac1{16}.
}
\label{eq:4pn-onebody-repair-ledger}
\end{equation}
The interpretation is immediate.  Relative to the carried $3$PN one-body
package, $4$PN opens
\begin{enumerate}
  \item a quartic static gate, represented by $\mu_{\rho,4}$,
  \item a quartic denominator datum, represented by $d_4$,
  \item and two genuinely new self-sector slots, represented by $s_{34}$ and
  $s_{26}$.
\end{enumerate}
So $4$PN is already not a one-parameter continuation of the $3$PN one-body
closure.  The admissible one-body front end of the local $4$PN theorem is the
full four-parameter ledger
\[
(\mu_{\rho,4},d_4,s_{34},s_{26}),
\]
not a single extra denominator coefficient.  The detailed repair algebra is
recorded in Appendix~\ref{app:onebody-repair}.

\subsection{Natural \texorpdfstring{$4$}{4}PN self/static seed}
\label{sec:onebody-seed}

Once the exact one-body target is frozen, the natural local self/static seed is
obtained by direct symmetric promotion of the six one-body channels appearing in
\eqref{eq:4pn-onebody-lagrangian-gate}.  The resulting two-body seed is
\begin{equation}
\begin{aligned}
L_{4,\mathrm{seed}}={}&
\frac{7}{256}(m_A v_A^{10}+m_B v_B^{10})
+\frac{75Gm_A m_B}{128r}(v_A^8+v_B^8)
+\frac{59G^2m_A m_B}{16r^2}(m_B v_A^6+m_A v_B^6)
\\[0.5ex]
&
+\frac{203G^3m_A m_B}{32r^3}(m_B^2 v_A^4+m_A^2 v_B^4)
+\frac{31G^4m_A m_B}{32r^4}(m_B^3 v_A^2+m_A^3 v_B^2)
+\frac{G^5m_A m_B}{16r^5}(m_B^4+m_A^4),
\end{aligned}
\label{eq:4pn-natural-self-static-seed}
\end{equation}
understanding the whole block as the coefficient of $c^{-8}$ in the full
conservative expansion.

The role of \eqref{eq:4pn-natural-self-static-seed} is exactly the same one
played by the corresponding natural seed below $4$PN.  It isolates the
universal one-body/self/static content before any genuine comparable-mass local
residual is solved.  No comparable-mass cross information is inserted here.
Each term is simply the unique exchange-symmetric two-body promotion of the
corresponding one-body channel, and the strict test-mass limit of the seed
reduces back to \eqref{eq:4pn-onebody-lagrangian-gate} exactly.

This seed should therefore be read as a theorem-level starting point, not as a
convenient ansatz.  It is the correct local front end because the one-body gate
forces it and because later generic-frame comparable-mass data must be measured
\emph{relative} to it.  In particular, the subsequent local comparable-mass
problem is not allowed to hide one-body repair information inside the residual.
That information has already been frozen here.  The later Hamiltonian-first lift
will modify neither the strict one-body content nor the self/static promotion
principle; it will only determine how the remaining comparable-mass local
content is represented consistently across charts.

\subsection{Exact quartic Legendre-transform identity}
\label{sec:onebody-quartic-compiler}

The second non-negotiable input is the exact quartic perturbative
Legendre-transform identity.  Let
\begin{equation}
L=L_0+\varepsilon L_1+\varepsilon^2L_2+\varepsilon^3L_3+\varepsilon^4L_4,
\label{eq:4pn-legendre-expansion}
\end{equation}
with quadratic $L_0$ and constant Newtonian mass matrix $M$.  Define
\begin{equation}
v_0=M^{-1}p,
\qquad
A_0=(\partial_vL_1)|_{v_0},
\qquad
B_0=(\partial_vL_2)|_{v_0},
\qquad
D_0=(\partial_vL_3)|_{v_0},
\label{eq:4pn-legendre-first-derivatives}
\end{equation}
\begin{equation}
C_0=(\partial_v^2L_1)|_{v_0},
\qquad
E_0=(\partial_v^2L_2)|_{v_0},
\qquad
T_0=(\partial_v^3L_1)|_{v_0}.
\label{eq:4pn-legendre-higher-derivatives}
\end{equation}
Then the exact perturbative Legendre coefficients through quartic order are
\begin{align}
H_1 &= -L_1(v_0),
\label{eq:4pn-legendre-h1}\\
H_2 &= -L_2(v_0)+\frac12A_0^TM^{-1}A_0,
\label{eq:4pn-legendre-h2}\\
H_3 &= -L_3(v_0)+A_0^TM^{-1}B_0-\frac12A_0^TM^{-1}C_0M^{-1}A_0,
\label{eq:4pn-legendre-h3}\\
H_4 &= -L_4(v_0)
+A_0^TM^{-1}D_0
+\frac12B_0^TM^{-1}B_0
-B_0^TM^{-1}C_0M^{-1}A_0
\nonumber\\
&\hspace{2em}
-\frac12A_0^TM^{-1}E_0M^{-1}A_0
+\frac12A_0^TM^{-1}C_0M^{-1}C_0M^{-1}A_0
+\frac16T_0[M^{-1}A_0,M^{-1}A_0,M^{-1}A_0].
\label{eq:4pn-legendre-h4}
\end{align}
Equations~\eqref{eq:4pn-legendre-h1}--\eqref{eq:4pn-legendre-h4} are the exact
ordinary/Hamiltonian bridge used throughout the local $4$PN chain.  Once the
lower-order ledger through $3$PN is frozen, the genuinely new $4$PN block is
not free to enter the Hamiltonian arbitrarily.  It enters through the rigid
quartic compiler above.

This has two immediate consequences.  First, applying
\eqref{eq:4pn-legendre-h1}--\eqref{eq:4pn-legendre-h4} to the strict one-body
Lagrangian gate reproduces the one-body $4$PN Hamiltonian seed used later by the
Hamiltonian-first local lift.  So the Hamiltonian one-body gate is not an
independent heuristic ansatz; it is a direct theorem-level output of the exact
Schwarzschild Lagrangian gate together with the quartic compiler.  Second, no
trustworthy local comparable-mass $4$PN solve can be formulated before this
compiler is frozen, because every later translation between fixed-chart
Hamiltonian targets and ordinary generic-frame representatives depends on it.

The detailed order-by-order inversion and derivation of
\eqref{eq:4pn-legendre-h1}--\eqref{eq:4pn-legendre-h4} are placed in
Appendix~\ref{app:quartic-compiler}.  For the main text, the important lesson is
structural.  The exact one-body gate, the minimal repair ledger, the natural
self/static seed, and the quartic compiler already determine the admissible
self-sector front end of the local $4$PN theorem.  What remains after this
section is therefore not a search for one-body or chart data, but the genuinely
comparable-mass local lift and the separate hereditary quadrupole bridge.

% Section: The local instantaneous 4PN theorem chain

\section{The local instantaneous \texorpdfstring{$4$}{4}PN theorem chain}
\label{sec:local-chain}

This section closes the entire \emph{local instantaneous} $4$PN sector.  Its
role is narrower than the full paper and sharper than the earlier executive
summary.  The hereditary/tail contribution is kept completely separate; the
only question addressed here is whether the strict one-body gate fixed in
Section~\ref{sec:onebody}, together with the exact quartic compiler
\eqref{eq:4pn-legendre-h4}, is enough to determine a full generic-frame local
ordinary representative of the two-body $4$PN sector.  The answer is yes.  The
key lesson is that the local problem never had an early span shortage.  The raw
exchange-symmetric scaffold is already large enough.  What matters instead is
chart choice, quotienting of the COM-blind null sector, and correct treatment of
the aligned ordinary seed.

The logic therefore runs through eight exact steps.  First, one builds the raw
local comparable-mass scaffold relative to the natural self/static seed
\eqref{eq:4pn-natural-self-static-seed} and checks its COM image.  Second, one
imports the fixed local target in the Hamiltonian chart and reads off the true
polynomial degree ceilings.  Third, one proves that the Hamiltonian-chart
scaffold is polynomially complete against the imported local target.  Fourth,
one chooses a canonical representative inside the resulting COM-blind null
family.  Fifth, one translates that representative back to the ordinary chart,
where the comparable-mass residual obeys an exact sign-flip theorem.  Sixth, one
shows that the remaining ordinary obstruction is only a seed-alignment
correction.  Seventh, one lifts that correction in a minimal structured seed
sector.  Finally, one assembles the natural seed, the lifted seed-alignment
correction, and the canonical translated residual into one exact generic-frame
ordinary local representative.  At that point the local theorem is closed.

\subsection{Raw local scaffold and COM completeness}
\label{sec:local-raw-scaffold}

With the strict one-body front end already frozen, the first local question is
whether the exchange-symmetric generic-frame comparable-mass scaffold has enough
room to support a full $4$PN interaction residual.  Using the invariant
variables
\[
a=v_A^2,
\qquad
b=v_B^2,
\qquad
c=\mathbf v_A\!\cdot\!\mathbf v_B,
\qquad
d=\mathbf v_A\!\cdot\!\mathbf n,
\qquad
e=\mathbf v_B\!\cdot\!\mathbf n,
\qquad
p=m_A,
\qquad
q=m_B,
\]
and imposing strict vanishing on the one-body branch, the raw local interaction
families have block sizes
\begin{equation}
\dim \mathcal Q = 52,
\qquad
\dim \mathcal T = 46,
\qquad
\dim \mathcal S = 29,
\qquad
\dim \mathcal U = 10,
\qquad
\dim \mathcal W = 2,
\label{eq:4pn-local-raw-block-sizes}
\end{equation}
where the notation follows the standard block decomposition
\[
\frac{G}{r}\,\mathcal Q,
\qquad
\frac{G^2}{r^2}\,\mathcal T,
\qquad
\frac{G^3}{r^3}\,\mathcal S,
\qquad
\frac{G^4}{r^4}\,\mathcal U,
\qquad
\frac{G^5}{r^5}\,\mathcal W.
\]
So before any quotienting or target import the local generic-frame interaction
scaffold already contains
\begin{equation}
52+46+29+10+2 = 139
\label{eq:4pn-local-raw-total}
\end{equation}
exchange-symmetric directions.

Projecting to the center-of-mass chart does not reveal a shortage.  It reveals a
large null family.  The induced COM interaction-slot structure is
\begin{equation}
\#_{\rm COM}^{\rm int}
=
\bigl(5,4,3,2,1\bigr),
\qquad
5+4+3+2+1 = 15,
\label{eq:4pn-local-com-slot-counts}
\end{equation}
corresponding respectively to the
\[
\frac{G}{r},
\qquad
\frac{G^2}{r^2},
\qquad
\frac{G^3}{r^3},
\qquad
\frac{G^4}{r^4},
\qquad
\frac{G^5}{r^5}
\]
interaction blocks.  The blockwise COM image ranks are already maximal at this
slot-counting level:
\begin{equation}
52\to 5,
\qquad
46\to 4,
\qquad
29\to 3,
\qquad
10\to 2,
\qquad
2\to 1.
\label{eq:4pn-local-com-slot-ranks}
\end{equation}
Equivalently, the raw scaffold spans the full COM interaction image block by
block, with total COM-nullity
\begin{equation}
(52-5)+(46-4)+(29-3)+(10-2)+(2-1)=124.
\label{eq:4pn-local-com-nullity-raw}
\end{equation}

The interpretation is important.  Local $4$PN never suffered from an early image
shortage.  Already at the raw-scaffold stage the question is not whether there
is enough room to span the reduced interaction sector.  The question is which
chart carries the correct polynomial degree ceilings and which quotient slice
extracts a canonical representative from the very large COM-blind family.

\subsection{Fixed local target import and the Hamiltonian-first lesson}
\label{sec:local-target-import}

The next step is to import the fixed-chart \emph{local} $4$PN target.  At this
order the conservative theorem splits into a local instantaneous sector and a
separate hereditary/tail sector, so the target imported here is deliberately the
local ADM Hamiltonian sector only.  This target passes the strict one-body gate:
a direct application of the quartic compiler \eqref{eq:4pn-legendre-h4} to the
one-body Lagrangian gate \eqref{eq:4pn-onebody-lagrangian-gate} yields the exact
reduced one-body Hamiltonian coefficients
\begin{equation}
\frac{7}{256}p^{10},
\qquad
\frac{45}{128}\frac{p^8}{r},
\qquad
\frac{13}{8}\frac{p^6}{r^2},
\qquad
\frac{105}{32}\frac{p^4}{r^3},
\qquad
\frac{105}{32}\frac{p^2}{r^4},
\qquad
-\frac{1}{16}\frac{1}{r^5},
\label{eq:4pn-local-onebody-hamiltonian-gate}
\end{equation}
which the imported local target reproduces exactly in the strict test-mass
limit.

The imported COM Hamiltonian target consists of one free slot plus the same
fifteen interaction slots counted above, so it is already structurally matched
to the scaffold.  The decisive issue is instead the \emph{polynomial degree
ceiling} in the symmetric mass ratio $\nu$.  In the Hamiltonian chart the local
target has blockwise ceilings
\begin{equation}
\bigl(G/r,\,G^2/r^2,\,G^3/r^3,\,G^4/r^4,\,G^5/r^5\bigr)
=
(4,3,3,2,2),
\label{eq:4pn-local-hamiltonian-degree-ceilings}
\end{equation}
whereas translating the same fixed local target back to the ordinary chart by
\eqref{eq:4pn-legendre-h4} produces
\begin{equation}
\bigl(G/r,\,G^2/r^2,\,G^3/r^3,\,G^4/r^4,\,G^5/r^5\bigr)
=
(4,4,4,4,2)
\label{eq:4pn-local-ordinary-degree-ceilings}
\end{equation}
in the same block order.  So the quartic compiler itself creates the extra
ordinary-chart $\nu^4$ tails in the middle and upper interaction blocks
\[
\frac{G^2}{r^2},
\qquad
\frac{G^3}{r^3},
\qquad
\frac{G^4}{r^4},
\]
while leaving the top static block unchanged.

This is the central structural lesson of the whole local chain.  The fixed local
$4$PN lift must be performed \emph{Hamiltonian-first}.  In the Hamiltonian chart
the constant-coefficient generic-frame scaffold matches the imported target
exactly at the level of degree ceilings.  In the translated ordinary chart the
same target develops extra polynomial tails that do not belong to the original
constant-coefficient comparable-mass scaffold.  So the correct order of work is
not
\[
\text{ordinary target} \longrightarrow \text{generic-frame lift},
\]
but rather
\begin{equation}
\boxed{
\text{fixed local Hamiltonian target}
\;\longrightarrow\;
\text{generic-frame Hamiltonian lift}
\;\longrightarrow\;
\text{ordinary translation only afterward}.
}
\label{eq:4pn-local-hamiltonian-first-scheme}
\end{equation}
That statement is no longer a heuristic design preference.  It is the exact
chart lesson forced by the quartic compiler.

\subsection{Hamiltonian-chart generic-frame local span theorem}
\label{sec:local-hamiltonian-lift}

Once the Hamiltonian chart is chosen, the local comparable-mass existence
question becomes exact.  The COM interaction residual is decomposed into
explicit $\nu$-polynomial coefficient space:
\begin{equation}
\begin{aligned}
\frac{G}{r}:&\quad 5\text{ slots}\times(\nu,\nu^2,\nu^3,\nu^4) &&\Rightarrow 20\text{ coefficients},\\
\frac{G^2}{r^2}:&\quad 4\text{ slots}\times(\nu,\nu^2,\nu^3) &&\Rightarrow 12,\\
\frac{G^3}{r^3}:&\quad 3\text{ slots}\times(\nu,\nu^2,\nu^3) &&\Rightarrow 9,\\
\frac{G^4}{r^4}:&\quad 2\text{ slots}\times(\nu,\nu^2) &&\Rightarrow 4,\\
\frac{G^5}{r^5}:&\quad 1\text{ slot}\times(\nu,\nu^2) &&\Rightarrow 2,
\end{aligned}
\label{eq:4pn-local-polynomial-coefficient-space}
\end{equation}
so the full interaction coefficient space has dimension
\begin{equation}
20+12+9+4+2 = 47.
\label{eq:4pn-local-interaction-rank-target}
\end{equation}

The exact blockwise rank test in the Hamiltonian chart is then
\begin{equation}
\begin{aligned}
\frac{G}{r}:&\quad 52\to 20, &\qquad& \text{nullity }32,\\
\frac{G^2}{r^2}:&\quad 46\to 12, && \text{nullity }34,\\
\frac{G^3}{r^3}:&\quad 29\to 9, && \text{nullity }20,\\
\frac{G^4}{r^4}:&\quad 10\to 4, && \text{nullity }6,\\
\frac{G^5}{r^5}:&\quad 2\to 2, && \text{nullity }0.
\end{aligned}
\label{eq:4pn-local-hamiltonian-ranks}
\end{equation}
Therefore the total interaction coefficient-space rank is maximal,
\begin{equation}
\operatorname{rank}=47,
\qquad
\operatorname{nullity}=32+34+20+6+0=92.
\label{eq:4pn-local-hamiltonian-total-rank-nullity}
\end{equation}

This proves the Hamiltonian-chart local span theorem: every interaction
coefficient of the imported fixed-chart COM Hamiltonian residual lies in the
image of the constant-coefficient exchange-symmetric generic-frame scaffold.
There is no Hamiltonian-chart span obstruction, no missing middle block, and no
need for additional seed dressing before the comparable-mass Hamiltonian lift is
performed.  The only ambiguity is the $92$-direction COM-blind null family.
That ambiguity is algebraic.  It is not a failure of local existence.

\subsection{Canonical comparable-mass local residual slice}
\label{sec:local-canonical-slice}

Because the Hamiltonian-chart image is complete, the next question is not
existence but quotienting.  The generic-frame invariants satisfy the usual COM
relations
\begin{equation}
\mathcal C_1 = pa+qc,
\qquad
\mathcal C_2 = pc+qb,
\qquad
\mathcal C_3 = pd+qe,
\label{eq:4pn-local-com-ideal-linear}
\end{equation}
\begin{equation}
\mathcal C_4 = ab-c^2,
\qquad
\mathcal C_5 = ae-cd,
\qquad
\mathcal C_6 = bd-ce.
\label{eq:4pn-local-com-ideal-quadratic}
\end{equation}
The exact ideal-membership audit shows that the $4$PN Hamiltonian-chart nullities
split as
\begin{equation}
Q:32,
\qquad
T:34,
\qquad
S:20,
\qquad
U:6,
\qquad
W:0,
\label{eq:4pn-local-nullities-by-block}
\end{equation}
and moreover that
\begin{enumerate}
  \item all $Q$-block null vectors lie in the full COM ideal
  $\langle \mathcal C_1,\ldots,\mathcal C_6\rangle$,
  \item all $T$-, $S$-, and $U$-block null vectors already lie in the linear
  momentum ideal $\langle \mathcal C_1,\mathcal C_2,\mathcal C_3\rangle$,
  \item and the top static $W$ block has zero nullity.
\end{enumerate}
So the entire $92$-direction ambiguity is purely COM-blind algebraic freedom.
No physical local datum is hiding there once the COM target is fixed.

The canonical quotient slice is therefore defined by the exact Gauss--Jordan lift
with all null coordinates set to zero.  This choice gives one sparse generic-frame
Hamiltonian representative of the comparable-mass local residual,
\begin{equation}
\Delta H_{4,\mathrm{loc}}^{\mathrm{can}}
=
\frac{Gpq}{r}Q_{\mathrm{can}}
+
\frac{G^2pq}{r^2}T_{\mathrm{can}}
+
\frac{G^3pq}{r^3}S_{\mathrm{can}}
+
\frac{G^4}{r^4}U_{\mathrm{can}}
+
\frac{G^5}{r^5}W_{\mathrm{can}},
\label{eq:4pn-canonical-hamiltonian-residual}
\end{equation}
with nonzero scaffold-direction counts
\begin{equation}
Q:11,
\qquad
T:12,
\qquad
S:9,
\qquad
U:4,
\qquad
W:2.
\label{eq:4pn-canonical-hamiltonian-sparsity}
\end{equation}
The explicit block formulas are lengthy and belong in the appendix.  The main
text only needs the structural conclusion: the Hamiltonian-chart local
comparable-mass existence problem is solved by one exact canonical
representative.

\subsection{Ordinary translation and the aligned-seed issue}
\label{sec:local-ordinary-translation}

The exact quartic compiler then transports the canonical comparable-mass
Hamiltonian residual back to the ordinary chart.  Once the lower-order ledger is
held fixed, the translation is rigid and remarkably simple:
\begin{equation}
\boxed{
\Delta L_{4,\mathrm{loc}}^{\mathrm{can}}
=
-\Delta H_{4,\mathrm{loc}}^{\mathrm{can}}.
}
\label{eq:4pn-local-sign-flip}
\end{equation}
This is the exact local $4$PN sign-flip theorem.  It shows that the canonical
comparable-mass residual itself is not the source of the earlier ordinary-chart
obstruction.  That obstruction sits elsewhere.

What remains misaligned in the ordinary chart is the seed.  The correct ordinary
seed is not simply the natural seed obtained by symmetric promotion of the
one-body gate.  It is the \emph{aligned} seed selected by the Hamiltonian-first
translation.  The resulting seed-alignment correction is
\begin{equation}
\delta L_{4,\mathrm{seed}}
=
L_{4,\mathrm{seed}}^{(\mathrm{aligned})}
-
L_{4,\mathrm{seed}}^{(\mathrm{natural})}.
\label{eq:4pn-local-seed-alignment-correction}
\end{equation}
The aligned seed already fixes the full free block of the local ordinary target
exactly, while the translated canonical comparable-mass residual fixes the true
interaction residual by \eqref{eq:4pn-local-sign-flip}.  Therefore the entire
remaining ordinary-chart mismatch is a \emph{seed-alignment} issue.  It is not a
failure of the local generic-frame comparable-mass lift.

This distinction is conceptually important.  The Hamiltonian chart has already
solved the genuine comparable-mass local problem.  The ordinary chart only has
to repair the seed so that the exact translated residual can sit on the correct
free and self/static background.  Once that is recognized, the old ordinary
obstruction becomes narrow and intelligible.

\subsection{Generic-frame lift of the aligned-seed ordinary correction}
\label{sec:local-aligned-seed-lift}

The seed-alignment correction \eqref{eq:4pn-local-seed-alignment-correction}
does not require a new comparable-mass interaction scaffold.  It requires a
minimal structured enlargement of the \emph{seed sector}.  Writing the original
local generic-frame block families as
\[
\mathcal Q,
\qquad
\mathcal T,
\qquad
\mathcal S,
\qquad
\mathcal U,
\qquad
\mathcal W,
\]
the exact lift statement is
\begin{equation}
Q_\delta \in \mathcal Q,
\qquad
T_\delta \in \mathcal T\oplus(pq)\,\mathcal T,
\qquad
S_\delta \in \mathcal S\oplus(pq)\,\mathcal S,
\qquad
U_\delta \in \mathcal U\oplus(p^2q^2)\,\mathcal U,
\qquad
W_\delta = 0.
\label{eq:4pn-local-structured-seed-spaces}
\end{equation}
In these structured seed spaces the coefficient-space ranks are again maximal:
\begin{equation}
Q:20/20,
\qquad
T:16/16,
\qquad
S:12/12,
\qquad
U:8/8.
\label{eq:4pn-local-structured-seed-ranks}
\end{equation}
So the aligned-seed correction is exactly generic-frame liftable once the seed
sector, rather than the residual sector, is dressed in this specific way.

A canonical sparse lift is again obtained by setting all structured null
coordinates to zero.  One such representative uses
\begin{equation}
Q:13,
\qquad
T:14,
\qquad
S:12,
\qquad
U:8,
\qquad
W:0
\label{eq:4pn-local-structured-seed-sparsity}
\end{equation}
nonzero scaffold directions.  The important point is not the particular sparse
choice.  It is the theorem-level structural message: the ordinary seed-alignment
correction is not an inexplicable extra obstruction.  It is an exactly liftable
generic-frame object inside a narrow and intelligible structured seed sector.

\subsection{Exact local referee reconstruction}
\label{sec:local-referee-reconstruction}

At this point all three local ingredients are available in generic-frame
ordinary form.

First, the natural generic-frame local seed is the direct symmetric promotion of
the one-body gate,
\begin{equation}
Q_{\rm nat}=\frac{75}{128}(a^4+b^4),
\qquad
T_{\rm nat}=\frac{59}{16}(qa^3+pb^3),
\qquad
S_{\rm nat}=\frac{203}{32}(q^2a^2+p^2b^2),
\label{eq:4pn-local-natural-generic-seed-1}
\end{equation}
\begin{equation}
U_{\rm nat}=\frac{31}{32}(q^3a+p^3b),
\qquad
W_{\rm nat}=\frac{1}{16}(q^4+p^4).
\label{eq:4pn-local-natural-generic-seed-2}
\end{equation}
Second, the aligned-seed correction has the exact structured lift described in
\eqref{eq:4pn-local-structured-seed-spaces}.  Third, the canonical comparable-mass
local residual has already been translated to the ordinary chart by the sign-flip
identity \eqref{eq:4pn-local-sign-flip}.

These pieces assemble into one exact generic-frame ordinary local candidate,
\begin{equation}
\boxed{
L_{4,\mathrm{loc}}^{(\mathrm{gen})}
=
L_{4,\mathrm{seed}}^{(\mathrm{natural,gen})}
+
\delta L_{4,\mathrm{seed}}^{(\mathrm{gen})}
+
\Delta L_{4,\mathrm{loc}}^{(\mathrm{can,gen})}.
}
\label{eq:4pn-local-generic-ordinary-candidate}
\end{equation}
Let $\Pi_{\rm COM}$ denote exact reduction to the center-of-mass chart.  The end-to-end
referee identity is then
\begin{equation}
\Pi_{\rm COM}\!\left[L_{4,\mathrm{loc}}^{(\mathrm{gen})}\right]
=
L_{4,\mathrm{loc}}^{(\mathrm{target,ord})},
\label{eq:4pn-local-referee-identity}
\end{equation}
with equality holding slot by slot in all five local interaction blocks:
\[
Q:5\text{ slots},
\qquad
T:4\text{ slots},
\qquad
S:3\text{ slots},
\qquad
U:2\text{ slots},
\qquad
W:1\text{ slot}.
\]
In other words, the local referee suite verifies not merely formal existence but
full end-to-end reconstruction of the fixed-chart local ordinary target.

\subsection{Final exact local theorem}
\label{sec:local-final}

The local theorem can now be stated without qualification.  Within the declared
closure hierarchy used throughout the lower conservative program, the complete
local instantaneous two-body $4$PN sector is fixed exactly.  More precisely:
\begin{enumerate}
  \item the strict one-body $4$PN Schwarzschild gate
  \eqref{eq:4pn-onebody-lagrangian-gate} and its natural self/static two-body
  promotion \eqref{eq:4pn-natural-self-static-seed} are frozen exactly;
  \item the raw exchange-symmetric generic-frame local scaffold is already
  blockwise COM-complete, so local $4$PN has no early span shortage;
  \item the correct fixed-chart import is Hamiltonian-first, because only the
  Hamiltonian chart respects the true degree ceilings of the local target;
  \item in that Hamiltonian chart there is no span obstruction, and one exact
  canonical generic-frame comparable-mass residual representative exists;
  \item the canonical residual translates back to the ordinary chart by exact
  sign flip;
  \item the remaining ordinary mismatch is only the aligned-seed correction,
  which is itself exactly generic-frame liftable in the structured seed spaces
  \eqref{eq:4pn-local-structured-seed-spaces};
  \item and the assembled candidate
  \eqref{eq:4pn-local-generic-ordinary-candidate} reproduces the full fixed-chart
  local ordinary target exactly under COM reduction.
\end{enumerate}
Equivalently, the full local instantaneous sector may be written in the exact
assembled form
\begin{equation}
L_{4\mathrm{PN}}^{\mathrm{local}}
=
L_{4,\mathrm{seed}}^{(\mathrm{natural,gen})}
+
\delta L_{4,\mathrm{seed}}^{(\mathrm{gen})}
+
\Delta L_{4,\mathrm{loc}}^{(\mathrm{can,gen})},
\qquad
\Pi_{\rm COM}\!\left[L_{4\mathrm{PN}}^{\mathrm{local}}\right]
=
L_{4,\mathrm{loc}}^{(\mathrm{target,ord})}.
\label{eq:4pn-local-final-theorem}
\end{equation}
So the local instantaneous $4$PN bottleneck is no longer algebraic span, chart
ambiguity, or generic-frame existence.  All of that is already closed.  What
remains for the full conservative $4$PN theorem is not more local algebra.  It
is the separate hereditary/tail bridge, whose unique compatible coefficient is
controlled by the same passive/outgoing STF quadrupole normalization already
isolated by the $2.5$PN program.

% Section: Tail / hereditary 4PN bridge

\section{Tail / hereditary \texorpdfstring{$4$}{4}PN bridge}
\label{sec:tail-bridge}

With the local instantaneous sector already fixed exactly, the remaining $4$PN
question is not another generic-frame existence solve.  The $2.5$PN audit has
already reduced the universal higher-order retarded route to the canonically
normalized orbital/worldtube STF quadrupole branch, while the local referee
chain has already closed the full instantaneous $4$PN sector.  The hereditary
side is therefore sharply narrower than the local side ever was.  The purpose of
this section is to make that narrowing explicit.  First, we freeze the standard
conservative GR tail functional and its universal local logarithmic shadow.
Second, we show that the hereditary coefficient is tied exactly to the same
canonically normalized STF quadrupole coefficient that already controls the
$2.5$PN Burke--Thorne channel.  Third, we isolate the only remaining model-side
ambiguity as a single scalar transport factor rather than a new tensor channel,
generic-frame lift, or normalization family.  The resulting conclusion is the
main structural message of the section: $4$PN opens no new independent
normalization datum beyond the already-known passive/outgoing STF quadrupole
gate.

\subsection{Exact GR tail structure}
\label{sec:tail-gr-structure}

The conservative $4$PN dynamics in GR splits into a local instantaneous sector
plus a time-symmetric nonlocal tail sector.  On the action side the hereditary
piece may be written as
\begin{equation}
S_{\mathrm{tail}}
=
\frac{G^2 M}{5c^8}
\operatorname{Pf}_{2s/c}
\iint \frac{dt\,dt'}{|t-t'|}
I_{ij}^{(3)}(t)
I_{ij}^{(3)}(t'),
\label{eq:4pn-tail-action}
\end{equation}
while on the Hamiltonian side the corresponding time-symmetric nonlocal term is
\begin{equation}
H_{4\mathrm{PN}}^{\mathrm{tail,sym}}(t)
=
-
\frac{G^2 M}{5c^8}
I_{ij}^{(3)}(t)
\operatorname{Pf}_{2s/c}
\int_{-\infty}^{+\infty}
\frac{dv}{|v|}
I_{ij}^{(3)}(t+v).
\label{eq:4pn-tail-hamiltonian}
\end{equation}
So the universal GR hereditary coefficient is
\begin{equation}
C_{\mathrm{tail}}^{\mathrm{GR}}
=
\frac{G^2 M}{5c^8}.
\label{eq:4pn-tail-gr-coefficient}
\end{equation}
Equations~\eqref{eq:4pn-tail-action}--\eqref{eq:4pn-tail-gr-coefficient}
show that the conservative hereditary $4$PN sector is a single STF-quadrupole
functional with one scalar coefficient.  In particular, the tail side is not a
second local-generic-frame lift problem.  Its representation is already fixed at
the level of the canonical STF quadrupole.

The same nonlocal functional has a local logarithmic shadow.  In the present
normalization that shadow is controlled by
\begin{equation}
F_{\mathrm{tail}}(t)
=
\frac{2}{5}\frac{G^2 M}{c^8}
\bigl(I_{ij}^{(3)}(t)\bigr)^2.
\label{eq:4pn-tail-local-shadow-functional}
\end{equation}
This is the local object that later appears when one compares the hereditary
sector with the local instantaneous $4$PN blocks.

\subsection{Universal Newtonian quadrupole shadow}
\label{sec:tail-newtonian-shadow}

Because the hereditary coefficient first enters at $4$PN, its local logarithmic
shadow is fixed already by the Newtonian order-reduced STF quadrupole.  Writing
\begin{equation}
I_{ij}
=
\mu\,\mathrm{STF}(x_i x_j),
\qquad
M \equiv m_A+m_B,
\qquad
\mu \equiv \frac{m_A m_B}{M},
\qquad
\nu \equiv \frac{\mu}{M},
\label{eq:4pn-tail-newtonian-quadrupole}
\end{equation}
one finds the exact order-reduced third derivative
\begin{equation}
I_{ij}^{(3)}
=
-
\frac{2GM\mu}{r^3}
\left(
4x_{\langle i}v_{j\rangle}
-
3\frac{\dot r}{r}x_{\langle i}x_{j\rangle}
\right),
\label{eq:4pn-tail-third-derivative}
\end{equation}
and therefore
\begin{equation}
\bigl(I_{ij}^{(3)}\bigr)^2
=
\frac{8}{3}\frac{G^2 M^2 \mu^2}{r^4}
\left(12v^2-11\dot r^2\right).
\label{eq:4pn-tail-third-derivative-square}
\end{equation}
Substituting \eqref{eq:4pn-tail-third-derivative-square} into
\eqref{eq:4pn-tail-local-shadow-functional} gives the exact local logarithmic
tail shadow
\begin{equation}
\boxed{
\frac{F_{\mathrm{tail}}}{\mu}
=
\frac{16}{15}\,\nu\,\frac{U^4}{c^8}
\left(12v^2-11\dot r^2\right),
\qquad
U\equiv \frac{GM}{r}.
}
\label{eq:4pn-tail-local-log-shadow}
\end{equation}
Equation~\eqref{eq:4pn-tail-local-log-shadow} makes the structural point
transparent.  The hereditary shadow lives only in the degree-$2$ local
\(G^4/r^4\) block, i.e. in the $U$ block of the local $4$PN bookkeeping.  That
is the exact reason the hereditary side is much smaller than the full local lift
problem.  The local theorem had to reconstruct an entire generic-frame
interaction sector.  The tail theorem only has to attach one scalar coefficient
to one already-fixed STF quadrupole representation.

\subsection{Exact bridge to the \texorpdfstring{$2.5$}{2.5}PN quadrupole coefficient}
\label{sec:tail-exact-bridge}

The $2.5$PN narrowing already fixed the correct representation and normalization
language for the surviving radiative branch.  On that canonical passive/outgoing
STF quadrupole branch the carried coefficient is written
\begin{equation}
\gamma_{\rm quad}^{\rm eff}
=
\mathcal N_Q\frac{a_{\rm th}^5}{27c_s^5}
=
\overline\Gamma_5
=
9\frac{\overline K_2^{5/2}}{\overline K_0^{3/2}},
\label{eq:4pn-tail-carried-quadrupole-coefficient}
\end{equation}
where $a_{\rm th}$ is the throat-radius scale carried by the quadrupole branch.
In GR the corresponding Burke--Thorne coefficient is
\begin{equation}
\gamma_{\rm GR}
=
\frac{2G}{5c^5}.
\label{eq:4pn-tail-burke-thorne-gr}
\end{equation}
The crucial arithmetic observation is that the $4$PN GR hereditary coefficient
is related exactly to the $2.5$PN Burke--Thorne coefficient by
\begin{equation}
\boxed{
C_{\mathrm{tail}}^{\mathrm{GR}}
=
\frac{GM}{2c^3}\,\gamma_{\rm GR}.
}
\label{eq:4pn-tail-gr-bridge}
\end{equation}
Indeed,
\[
\frac{GM}{2c^3}\,\frac{2G}{5c^5}
=
\frac{G^2 M}{5c^8}
=
C_{\mathrm{tail}}^{\mathrm{GR}}.
\]
So the hereditary coefficient is not merely analogous to the already-isolated
quadrupole odd coefficient; it is fixed by the same coefficient through the
exact bridge \eqref{eq:4pn-tail-gr-bridge}.

This observation upgrades directly from the GR reference coefficient to the
canonically normalized branch carried by the toy-model program.  The unique
compatible conservative hereditary coefficient is
\begin{equation}
\boxed{
C_{\mathrm{tail}}
=
\frac{GM}{2c^3}\,\gamma_{\rm quad}^{\rm eff}.
}
\label{eq:4pn-tail-exact-bridge}
\end{equation}
Equation~\eqref{eq:4pn-tail-exact-bridge} is the cleanest structural statement
of the whole hereditary bridge.  Once the STF quadrupole branch is canonically
normalized, the conservative $4$PN tail coefficient follows automatically.
There is no second normalization problem hiding on the hereditary side.

\subsection{Toy-model branch form and the remaining transport factor}
\label{sec:tail-toy-branch}

The same bridge can be written in a form that isolates the only remaining
model-side ambiguity.  Let
\begin{equation}
\gamma_{\rm toy}
\equiv
\hat m_0^2\Gamma_5
\label{eq:4pn-tail-toy-gamma}
\end{equation}
be the canonically normalized STF quadrupole coefficient inherited from the
$2.5$PN audit.  Then the hereditary side may be parameterized by one scalar
transport factor $\Theta_{\rm tail}$ through
\begin{equation}
\boxed{
\alpha_{\rm tail}^{\rm toy}
=
\Theta_{\rm tail}
\frac{GM}{2c_s^3}
\gamma_{\rm toy}.
}
\label{eq:4pn-tail-toy-transport}
\end{equation}
On the branch with $c_s=c$ and
\begin{equation}
\gamma_{\rm toy}=
\gamma_{\rm quad}^{\rm eff},
\label{eq:4pn-tail-toy-canonical-branch}
\end{equation}
Eq.~\eqref{eq:4pn-tail-toy-transport} reduces to the exact bridge
\eqref{eq:4pn-tail-exact-bridge} when
\begin{equation}
\Theta_{\rm tail}=1.
\label{eq:4pn-tail-theta-one}
\end{equation}
This is the sharpest way to say what is and is not still open.  The
representation/source problem has already been solved by the $2.5$PN quadrupole
audit.  The local $4$PN existence problem has already been solved by the local
referee chain.  The only possible remaining hereditary mismatch is one scalar
monopole-scattering transport factor.

\subsection{Why no new independent \texorpdfstring{$4$}{4}PN normalization datum opens}
\label{sec:tail-no-new-datum}

The hereditary theorem may now be stated without further search language.
Inside the same closure hierarchy used throughout the lower-order conservative
program, the following statements hold:
\begin{enumerate}
  \item the unique compatible conservative hereditary coefficient is the bridge
  coefficient \eqref{eq:4pn-tail-exact-bridge},
  \item therefore the $4$PN tail sector introduces no new independent
  quadrupole-normalization datum,
  \item and full GR matching of the hereditary coefficient is equivalent to the
  already-known $2.5$PN normalization target
  \begin{equation}
  \gamma_{\rm quad}^{\rm eff}
  =
  \frac{2G}{5c^5}.
  \label{eq:4pn-tail-gr-recovery-condition}
  \end{equation}
\end{enumerate}
The same recovery condition may be written in the normalization language carried
by the moving-throat program as
\begin{equation}
\mathcal N_Q^{\rm target}
=
\frac{54Gc_s^5}{5a_{\rm th}^5 c^5}.
\label{eq:4pn-tail-normalization-target}
\end{equation}
Substituting \eqref{eq:4pn-tail-gr-recovery-condition} into
\eqref{eq:4pn-tail-exact-bridge} gives
\begin{equation}
C_{\mathrm{tail}}
=
\frac{GM}{2c^3}\,\frac{2G}{5c^5}
=
\frac{G^2 M}{5c^8}
=
C_{\mathrm{tail}}^{\mathrm{GR}}.
\label{eq:4pn-tail-gr-recovery}
\end{equation}
So the $4$PN hereditary bridge closes automatically once the already-isolated
passive/outgoing STF quadrupole normalization is fixed on the natural branch.

What remains open is correspondingly narrow.  The present section does not claim
a first-principles moving-throat PDE derivation of the passive/outgoing
quadrupole normalization, nor does it claim the final insertion of that
normalized tail module into a single end-to-end referee/master script.  But the
bottleneck is now fully exposed: the local $4$PN sector is already closed, the
hereditary coefficient is fixed by the same invariant quadrupole branch as
$2.5$PN, and the remaining full-$4$PN gap is therefore the same passive/outgoing
quadrupole-normalization theorem already identified by the $2.5$PN program.

% Section: Final conditional conservative 4PN theorem

\section{Final conditional conservative \texorpdfstring{$4$}{4}PN theorem}
\label{sec:final-conditional-4pn-theorem}

The preceding sections have separated the $4$PN problem into two logically distinct
pieces.  The first is the full local instantaneous sector, whose one-body gate,
quartic compiler, Hamiltonian-chart lift, canonical quotient slice, aligned-seed
repair, and referee reconstruction have already been closed exactly inside the
declared hierarchy.  The second is the hereditary/tail sector, whose coefficient
is not an independent new tensor datum but a scalar bridge to the same
canonically normalized STF quadrupole branch that was already isolated by the
$2.5$PN audit.  The purpose of the present section is therefore not to introduce
new machinery, but to state the final theorem in the sharpest possible form and
record exactly what remains conditional.

\begin{theorem}[Final conditional conservative \texorpdfstring{$4$}{4}PN theorem]
\label{thm:final-conditional-4pn}
Work within the same declared closure hierarchy used in the Newtonian, $1$PN,
$2$PN, and $3$PN papers, together with the carried $2.5$PN narrowing to the
canonically normalized orbital/worldtube STF quadrupole branch.  Then the full
conservative two-body $4$PN coefficient admits the exact split
\begin{equation}
L_{4\mathrm{PN}}^{\mathrm{cons}}
=
L_{4\mathrm{PN}}^{\mathrm{local}}
+
L_{4\mathrm{PN}}^{\mathrm{tail}},
\label{eq:final-4pn-split}
\end{equation}
with the following properties.
\begin{enumerate}[label=(\roman*)]
  \item The local instantaneous sector \(L_{4\mathrm{PN}}^{\mathrm{local}}\)
  is already fixed exactly inside the declared hierarchy.  In particular, there
  exists an exact generic-frame ordinary representative whose COM reduction
  reproduces the full fixed-chart local ordinary $4$PN target slot by slot.

  \item The hereditary coefficient is fixed by the exact bridge
  \begin{equation}
  C_{\rm tail}=
  \frac{GM}{2c^3}\,\gamma_{\rm quad}^{\rm eff},
  \label{eq:final-4pn-tail-bridge}
  \end{equation}
  where \(\gamma_{\rm quad}^{\rm eff}\) is the canonically normalized passive/
  outgoing STF quadrupole coefficient already isolated by the $2.5$PN program.

  \item Consequently, $4$PN opens no new independent normalization datum beyond
  that already-known quadrupole coefficient.  Full conservative GR recovery at
  $4$PN is equivalent to the single scalar condition
  \begin{equation}
  \gamma_{\rm quad}^{\rm eff}=\frac{2G}{5c^5},
  \label{eq:final-4pn-gr-recovery-condition}
  \end{equation}
  in which case
  \begin{equation}
  C_{\rm tail}
  =
  \frac{GM}{2c^3}\,\gamma_{\rm quad}^{\rm eff}
  =
  \frac{G^2M}{5c^8}
  =
  C_{\rm tail}^{\rm GR}.
  \label{eq:final-4pn-gr-tail-recovery}
  \end{equation}
\end{enumerate}
Therefore the only remaining conditional interface for a full conservative $4$PN
GR theorem is the passive/outgoing STF quadrupole normalization on the actual
moving-throat branch.
\end{theorem}

\begin{proof}
The proof is the synthesis of the earlier sections.

First, the strict one-body Schwarzschild gate, the exact quartic perturbative
Legendre compiler, the raw exchange-symmetric generic-frame scaffold, the fixed
local Hamiltonian target, and the Hamiltonian-first generic-frame lift together
show that the local instantaneous $4$PN problem has no remaining span
obstruction.  The canonical Hamiltonian residual is exactly reachable in the
Hamiltonian chart, its ordinary translation obeys the exact sign-flip theorem,
and the only ordinary-chart obstruction is an aligned-seed issue rather than a
failure of the local generic-frame lift itself.  That aligned-seed correction is
then exactly liftable in the structured seed spaces
\[
Q,
\qquad
T\oplus(pq)T,
\qquad
S\oplus(pq)S,
\qquad
U\oplus(p^2q^2)U,
\qquad
W=0,
\]
so the full fixed-chart local ordinary $4$PN target is reproduced exactly after
COM reduction.  This establishes item~(i).

Second, the hereditary side is already much narrower than the local side.  The
conservative GR tail is a single STF-quadrupole functional with coefficient
\(C_{\rm tail}^{\rm GR}=G^2M/(5c^8)\), and its Newtonian local shadow lives only
in the degree-$2$ \(G^4/r^4\) block.  The $2.5$PN narrowing had already shown
that the only surviving universal higher-order branch is the canonically
normalized passive/outgoing STF quadrupole route.  Matching the universal GR
hereditary coefficient to that carried quadrupole coefficient gives the exact
bridge
\[
C_{\rm tail}=\frac{GM}{2c^3}\,\gamma_{\rm quad}^{\rm eff},
\]
which proves item~(ii).

Finally, because the hereditary coefficient depends only on the same scalar
quadrupole normalization already isolated by the $2.5$PN program, no second
independent $4$PN normalization datum can open.  GR recovery of the hereditary
sector is therefore equivalent to the single condition
\(
\gamma_{\rm quad}^{\rm eff}=2G/(5c^5)
\), which immediately yields
\(
C_{\rm tail}=G^2M/(5c^8)
\).  This proves item~(iii) and the final claim.
\end{proof}

\paragraph{Corollary (no new \texorpdfstring{$4$}{4}PN normalization datum).}
Within the present hierarchy, the $4$PN conservative sector introduces no new
independent normalization problem beyond the already-known passive/outgoing STF
quadrupole normalization.  Equivalently, once the moving-throat PDE fixes
\(\gamma_{\rm quad}^{\rm eff}\), both the $2.5$PN Burke--Thorne coefficient and
the conservative $4$PN tail coefficient are fixed simultaneously.  This follows
immediately from \eqref{eq:final-4pn-tail-bridge}: the same scalar branch
normalization that determines the odd quadrupole coefficient at $2.5$PN also
determines the conservative hereditary coefficient at $4$PN.

\paragraph{Interpretation.}
The theorem above is the correct endpoint of the present paper.  The local
instantaneous $4$PN sector is no longer provisional, symbolic, or span-limited:
it is already closed exactly inside the declared hierarchy.  The hereditary side
is likewise no longer diffuse: it is controlled by one scalar bridge to the
canonically normalized STF quadrupole branch.  Thus the remaining open problem
is not a generic $4$PN algebraic search, but the completed moving-throat PDE
realization of the passive/outgoing quadrupole normalization.

\paragraph{Practical reading rule.}
The correct conservative $4$PN verdict is therefore:
\begin{center}
\fbox{%
\parbox{0.9\linewidth}{%
The local instantaneous sector is already exact.  The full conservative $4$PN
sector is conditionally complete, and the sole remaining condition is
\(
\gamma_{\rm quad}^{\rm eff}=2G/(5c^5)
\), i.e. the same passive/outgoing STF
quadrupole normalization already isolated by the $2.5$PN program.}}
\end{center}

% Section: Interface to the 2.5PN program and the moving-throat PDE

\section{Interface to the \texorpdfstring{$2.5$}{2.5}PN program and the moving-throat PDE}
\label{sec:interface-25pn-pde}

From the viewpoint of the present paper, the $2.5$PN audit and the conditional
$4$PN theorem are not two unrelated higher-order stories.  They are the
retarded and conservative faces of one and the same quadrupole program.  The
$2.5$PN work already narrowed the universal higher-order bridge to the
canonically normalized orbital/worldtube STF quadrupole branch; the present
$4$PN paper shows that the full local instantaneous algebra is already closed
and that the hereditary coefficient depends on that same branch through one
scalar bridge.  The purpose of this section is to state that interface cleanly,
record what the present $4$PN result contributes to it, and identify the exact
piece that the completed moving-throat PDE must still supply.

The strategic message is simple.  The next theorem gate is no longer another
local $4$PN generic-frame solve, another seed search, or another free
normalization family.  It is the passive/outgoing STF quadrupole normalization
on the natural moving-throat branch.  Once that branch is fixed, the odd
$2.5$PN Burke--Thorne coefficient and the conservative $4$PN hereditary
coefficient close together.

\subsection{Why the same STF quadrupole branch controls \texorpdfstring{$2.5$}{2.5}PN and \texorpdfstring{$4$}{4}PN}
\label{sec:interface-same-branch}

The $2.5$PN audit already showed that the frozen conservative hierarchy cannot
produce a GR-like local radiation-reaction theorem by itself and that, on the
strict small-body branch, the scalar and dipole odd channels are pushed above
universal point-particle $2.5$PN.  The surviving universal route is the
canonically normalized orbital/worldtube STF quadrupole branch packaged in the
isotropic grouped-real $P_2$ language.  On the passive/outgoing compact branch,
the normalized outgoing $\ell=2$ fingerprint is
\begin{equation}
\widehat Y_2^{\rm(out)}(\omega)
=
1+\frac{a_{\rm th}^2\omega^2}{9c_s^2}
+\frac{4a_{\rm th}^4\omega^4}{81c_s^4}
+i\,\frac{a_{\rm th}^5\omega^5}{27c_s^5}
+O(\omega^6),
\label{eq:interface-outgoing-l2-fingerprint}
\end{equation}
where $a_{\rm th}$ is the throat-radius scale on the moving-throat branch.  On
the isotropic grouped branch the conservative and outgoing low-frequency data
collapse to common coefficients,
\begin{equation}
D_{20,n}=D_{21,n}=D_{22,n}=D_n,
\qquad
N_{20,n}=N_{21,n}=N_{22,n}=N_n,
\label{eq:interface-grouped-collapse}
\end{equation}
so the static outgoing prefactor is
\begin{equation}
P_0=\frac{N_0}{D_0}.
\label{eq:interface-P0}
\end{equation}
The leading odd quadrupole coefficient is then controlled by the same static
prefactor,
\begin{equation}
\Gamma_5
=
P_0\,\frac{a_{\rm th}^5}{27c_s^5},
\qquad
\gamma_{\rm quad}^{\rm eff}
=
\hat m_0^{\,2}\Gamma_5
=
\hat m_0^{\,2}P_0\,\frac{a_{\rm th}^5}{27c_s^5},
\label{eq:interface-gamma-from-P0}
\end{equation}
with $\hat m_0$ the source-map normalization factor on the natural branch.
Equivalently, in the invariant-pair language already isolated by the quadrupole
normalization package,
\begin{equation}
\gamma_{\rm quad}^{\rm eff}
=
9\,\frac{\overline K_2^{5/2}}{\overline K_0^{3/2}}.
\label{eq:interface-gamma-invariant-pair}
\end{equation}

This is exactly the same scalar coefficient that enters the conservative
hereditary bridge at $4$PN.  The present paper has shown that the hereditary
coefficient is not an independent new tensor datum, but rather
\begin{equation}
C_{\rm tail}
=
\frac{GM}{2c^3}\,\gamma_{\rm quad}^{\rm eff}.
\label{eq:interface-tail-bridge}
\end{equation}
So the $2.5$PN odd channel and the $4$PN conservative tail are controlled by
one and the same canonically normalized STF quadrupole branch.  The difference
between the two orders is not a difference of branch, but only a difference in
which coefficient of that branch is being read: the odd $\omega^5$ coefficient
for $2.5$PN and the conservative hereditary coefficient for $4$PN.

This is the main conceptual simplification produced by the combined $2.5$PN and
$4$PN work.  The higher-order bridge is no longer split into an apparently
``dissipative'' quadrupole problem and a separate ``conservative tail''
problem.  It is one normalization problem viewed in two different PN sectors.

\subsection{What the present \texorpdfstring{$4$}{4}PN result fixes for the quadrupole program}
\label{sec:interface-what-4pn-fixes}

The positive interface result of the present paper is not that it computes a new
branch parameter.  Its actual contribution is sharper: it proves that no such
new $4$PN-specific parameter is needed.  The local instantaneous $4$PN sector is
already closed exactly inside the declared hierarchy, and the hereditary
coefficient has been reduced to the scalar bridge
\eqref{eq:interface-tail-bridge}.  Therefore the quadrupole program no longer
has two conceptually independent open fronts.  The only live front is the
normalization of the passive/outgoing STF quadrupole branch itself.

The implication chain may be written in the compact form
\begin{equation}
\hat m_0^{\,2}P_0
=
\frac{54Gc_s^5}{5a_{\rm th}^5c^5}
\Longleftrightarrow
\gamma_{\rm quad}^{\rm eff}=\frac{2G}{5c^5}
\Longleftrightarrow
C_{\rm tail}=\frac{G^2M}{5c^8}.
\label{eq:interface-implication-chain}
\end{equation}
The first equality is the invariant moving-throat normalization target inherited
from the reduced $2.5$PN quadrupole package; the middle equality is the
canonical quadrupole-normalization statement; and the final equality is the GR
value of the conservative $4$PN hereditary coefficient.  Thus the present $4$PN
paper fixes the conservative end of the chain completely: once the moving-throat
side lands on the target value of $\hat m_0^{\,2}P_0$, the hereditary $4$PN
sector is no longer open.

In particular, the role of the present paper in the broader quadrupole program
is now very specific.  The $2.5$PN audit had already shown that, on the
isotropic real-$P_2$ branch, the grouped conservative coefficients through
$O(\omega^4)$ form the front end from which the odd Burke--Thorne coefficient is
read.  The present $4$PN paper proves that the same normalized branch also fixes
the conservative hereditary coefficient and that $4$PN opens no second
normalization datum beyond that already-known branch.  So the old question
``what extra $4$PN information is still needed after the $2.5$PN audit?'' now
has a very short answer: none beyond the same branch normalization.

This also clarifies what the present $4$PN theorem does \emph{not} leave open.
It does not leave a local-algebra bottleneck.  It does not leave an unfixed
local generic-frame existence question.  And it does not leave a separate tail
normalization family.  The remaining problem is genuinely a moving-throat branch
problem, not a missing post-Newtonian bookkeeping problem.

\subsection{What the full moving-throat PDE still has to supply}
\label{sec:interface-pde-still-open}

What remains open is therefore sharp and narrow.  The completed moving-throat
PDE must still supply the actual passive/outgoing quadrupole normalization on
its natural branch.  In the grouped-real $P_2$ language, this means that the PDE
must determine the conservative operator moments
\begin{equation}
D_A^{\rm(cons)}(\omega)
=
D_{A,0}+D_{A,2}\omega^2+D_{A,4}\omega^4+O(\omega^6),
\label{eq:interface-grouped-operator}
\end{equation}
together with the outgoing-transfer moments
\begin{equation}
N_A(\omega)
=
N_{A,0}+N_{A,2}\omega^2+N_{A,4}\omega^4+O(\omega^6),
\label{eq:interface-grouped-transfer}
\end{equation}
for the grouped lanes $A\in\{20,21,22\}$.  On the natural isotropic branch,
these data must satisfy the grouped isotropy gate
\begin{equation}
a_2=b_2=a_4=b_4=0,
\label{eq:interface-isotropy-gate}
\end{equation}
so that the outgoing prefactor collapses to the scalar object $P_0=N_0/D_0$ of
\eqref{eq:interface-P0}.  Equivalently, in the invariant-pair language, the PDE
must determine the canonical pair $(\overline K_0,\overline K_2)$ on the
passive/outgoing STF quadrupole branch.

The point is that the remaining PDE task is no longer diffuse.  It is not
``derive all of higher-order radiation reaction from scratch.''  It is to compute
one isotropic branch and evaluate one invariant normalization product.  In the
most compact form, the outstanding theorem gate is
\begin{equation}
\boxed{
\hat m_0^{\,2}P_0
=
\frac{54Gc_s^5}{5a_{\rm th}^5c^5}
}
\label{eq:interface-normalization-target}
\end{equation}
or, equivalently,
\begin{equation}
\boxed{
\gamma_{\rm quad}^{\rm eff}=\frac{2G}{5c^5}.
}
\label{eq:interface-gamma-target}
\end{equation}
Once this is true on the actual moving-throat branch, both the normalized
$2.5$PN odd quadrupole coefficient and the normalized $4$PN conservative tail
coefficient close simultaneously.

The current reduced moving-throat program sharpens this still further.  On the
natural isotropic grouped-$P_2$ one-pole branch, if the geometry lane is static
through $O(\omega^4)$, then the conservative precursor is forced into the
minimal form
\begin{equation}
\widehat Y_Q^{\rm cons}(\omega)
=
\frac34+\frac14\frac{1}{1-\omega^2/\Omega_Q^2},
\label{eq:interface-minimal-conservative-module}
\end{equation}
and the explicit Family-1 support/source side is no longer the active bottleneck.
A convenient reduced defect packaging is then
\begin{equation}
N_Q:=\frac{\overline K_0}{\overline K_0^{\rm target}},
\qquad
\frac{\overline\Gamma_5}{\overline\Gamma_5^{\rm target}}=\chi_Q N_Q,
\qquad
\hat m_0^{\,2}\chi_Q N_Q=1,
\label{eq:interface-NQ-chiQ}
\end{equation}
where $N_Q$ measures the conservative grouped-quadrupole normalization and
$\chi_Q$ measures the outgoing $\ell=2$ DtN normalization.  On the canonical
compact outgoing branch one finds the reduced closure $\chi_Q=1$, so in the
strict point-particle limit the condition becomes $N_Q=1$.  This is a sharpened
reduced guide to the PDE problem, not yet a theorem that the completed
moving-throat PDE must realize the canonical branch rather than a deformed one.

Two additional tasks remain attached to that PDE-side normalization problem, but
both are now secondary rather than diffuse.  First, the moving-throat PDE should
still explain the deeper microscopic origin of the grouped conservative
quadrupole data that appear on the isotropic branch.  Second, once the
normalization is fixed, the normalized hereditary module should be inserted into
a full release-grade referee/master package so that the conditional $4$PN
statement becomes a fully assembled theorem package rather than only a theorem
ledger.  Neither of these is a new $4$PN-specific parameter search; both are
assembly problems downstream of the same quadrupole branch.

\paragraph{Correct next theorem gate.}
The correct next step after the present paper is therefore not more local $4$PN
algebra.  It is to derive the passive/outgoing STF quadrupole normalization on
the moving-throat branch, or equivalently to determine $(\overline K_0,
\overline K_2)$ --- or the invariant product $\hat m_0^{\,2}P_0$ --- on the
natural isotropic branch.  Once that gate is closed, the normalized $2.5$PN odd
quadrupole coefficient and the normalized $4$PN hereditary coefficient become
simultaneously fixed inside the same hierarchy.

% Section: Fixed versus open quantities

\section{Fixed versus open quantities}
\label{sec:fixed-open}

The conditional theorem established above closes the full local instantaneous
$4$PN ledger and fixes the hereditary bridge, but it does not erase the status
boundary that runs through the whole series.  As in the earlier conservative
papers, one must keep separate what is exact in the carried reduction backbone,
what is frozen as lower-order input, what is fixed by the present $4$PN audit,
and what still belongs to the unsolved moving-throat completion.  The purpose of
this section is to record that boundary in one place.  Its role is not to weaken
the theorem, but to state its true strength: the conservative $4$PN coefficient
problem is no longer open-ended, and the remaining open work has collapsed to
one already-known quadrupole-normalization interface.

\subsection{Frozen carry-forward inputs}
\label{sec:fixed-open-frozen}

The present paper does not begin from a blank effective theory.  It imports the
exact $4$D parent reduction backbone already frozen by the earlier papers: the
projection map, projected continuity with leakage, the exact longitudinal
identity, the quasi-static Poisson hook, and the coherent-defect / worldtube
particle reduction.  These formulas remain the structural dictionary that turns
the bulk ontology into the near-zone two-body language used throughout the
post-Newtonian series.

The lower-order conservative ledger is carried in exactly the same way.  The
frozen Newtonian and $1$PN constants remain
\begin{equation}
\kappa_\rho=1,
\qquad
n=5,
\qquad
\kappa_{\rm add}=\frac12,
\qquad
\kappa_{\rm PV}=\frac32,
\qquad
\beta_{\rm 1PN}=3,
\label{eq:fixed-open-lower-order-ledger}
\end{equation}
and the full conservative Newtonian+$1$PN, $2$PN, and $3$PN assemblies are
already treated as fixed lower-order output inside the same declared closure
hierarchy.  In particular, the exact $1$PN EIH match, the exact $2$PN ADM
match, and the full conservative $3$PN chart-level closure are imported here as
frozen theorem input rather than reopened as free data.

The $2.5$PN narrowing is likewise a frozen input.  The surviving universal
higher-order route had already been reduced to the orbital/worldtube STF
quadrupole branch, packaged in the grouped real language as
\begin{equation}
P_{20}\oplus P_{21}\oplus P_{22}.
\label{eq:fixed-open-grouped-bundle}
\end{equation}
That means the present $4$PN paper is not searching for a new higher-order
carrier.  It is built from the start on the statement that the only universal
retarded interface already lives on the canonically normalized STF quadrupole
branch isolated by the $2.5$PN audit.

\subsection{Quantities fixed by the present \texorpdfstring{$4$}{4}PN audit}
\label{sec:fixed-open-fixed-now}

The first new fixed data are the strict one-body $4$PN repair parameters.  The
exact Schwarzschild gate freezes the four-parameter continuation ledger
\begin{equation}
\mu_{\rho,4}=\frac18,
\qquad
 d_4=\frac{205}{16},
\qquad
 s_{34}=-\frac{15}{32},
\qquad
 s_{26}=-\frac1{16}.
\label{eq:fixed-open-onebody-ledger}
\end{equation}
So the carried $3$PN self-sector packaging does \emph{not} extend to $4$PN with
one new scalar knob.  The exact one-body front end already fixes a quartic
static datum, a quartic denominator datum, and two genuinely new self-sector
slots.

The second new fixed object is the exact order-$4$ ordinary/Hamiltonian bridge.
Once the quartic compiler is frozen, the local $4$PN chart problem is no longer
a matter of heuristic translation.  Combined with the imported local ADM target,
it yields the Hamiltonian-first local lift theorem: the raw exchange-symmetric
local scaffold is already complete blockwise in the Hamiltonian chart, the
canonical comparable-mass local residual exists there exactly, and the ordinary
translation obeys the exact sign-flip rule
\begin{equation}
\Delta H_4=-\Delta L_4.
\label{eq:fixed-open-signflip}
\end{equation}
So the ordinary-chart obstruction is not a failure of the local generic-frame
lift itself.  It is an aligned-seed issue, and that issue is also fixed by the
present audit.  In particular, the aligned ordinary local seed is generically
liftable once the seed sector is enlarged in the exact structured way
\begin{equation}
Q,
\qquad
T\oplus(pq)T,
\qquad
S\oplus(pq)S,
\qquad
U\oplus(p^2q^2)U,
\qquad
W=0.
\label{eq:fixed-open-structured-seed-space}
\end{equation}
Therefore the whole fixed-chart local ordinary $4$PN target now has an exact
generic-frame representative.  In theorem language, the entire local
instantaneous sector is already fixed inside the declared hierarchy.

The third new fixed result is the hereditary bridge itself.  The tail side does
not open a new tensor family or a second normalization ladder.  Once the
canonically normalized STF quadrupole coefficient is denoted by
$\gamma_{\rm quad}^{\rm eff}$, the exact hereditary coefficient is
\begin{equation}
C_{\rm tail}
=
\frac{GM}{2c^3}\,\gamma_{\rm quad}^{\rm eff}.
\label{eq:fixed-open-tail-bridge}
\end{equation}
So the full conservative $4$PN theorem split may now be written as
\begin{equation}
L_{4\mathrm{PN}}^{\mathrm{cons}}
=
L_{4\mathrm{PN}}^{\mathrm{local}}
+
L_{4\mathrm{PN}}^{\mathrm{tail}},
\qquad
C_{\rm tail}
=
\frac{GM}{2c^3}\,\gamma_{\rm quad}^{\rm eff},
\label{eq:fixed-open-final-split}
\end{equation}
and the meaning of \eqref{eq:fixed-open-final-split} is now sharp: the entire
local instantaneous sector is already exact, and the hereditary side is already
reduced to the same quadrupole branch that controlled the $2.5$PN
Burke--Thorne channel.  In particular, \eqref{eq:fixed-open-tail-bridge} proves
that $4$PN opens no new independent normalization datum beyond that branch.

\begin{center}
\renewcommand{\arraystretch}{1.15}
\begin{tabular}{@{}>{\raggedright\arraybackslash}p{0.31\linewidth}>{\raggedright\arraybackslash}p{0.14\linewidth}>{\raggedright\arraybackslash}p{0.47\linewidth}@{}}
\textbf{Object} & \textbf{Status} & \textbf{Role at $4$PN} \\
\hline
Parent reduction backbone & frozen & Exact/reduced dictionary inherited from the $4$D parent theory \\
Lower-order Newtonian+$1$PN+$2$PN+$3$PN ledger & frozen & Fixed conservative input and chart conventions \\
$2.5$PN STF quadrupole narrowing & frozen & Fixes the only surviving universal higher-order branch \\
$(\mu_{\rho,4},d_4,s_{34},s_{26})$ & fixed here & Exact one-body $4$PN repair ledger \\
Quartic compiler & fixed here & Exact ordinary/Hamiltonian bridge through $4$PN \\
Hamiltonian-first local lift and \(\Delta H_4=-\Delta L_4\) & fixed here & Exact comparable-mass local chart theorem \\
Structured aligned-seed lift & fixed here & Exact ordinary-chart seed correction repair \\
$L_{4\mathrm{PN}}^{\mathrm{local}}$ & fixed here & Complete local instantaneous $4$PN sector \\
$C_{\rm tail}=\tfrac{GM}{2c^3}\gamma_{\rm quad}^{\rm eff}$ & fixed here & Hereditary bridge; no new $4$PN datum \\
\(\gamma_{\rm quad}^{\rm eff}=2G/(5c^5)\) & open & Sole remaining normalization condition \\
\end{tabular}
\end{center}

\subsection{Quantities that remain symbolic or open}
\label{sec:fixed-open-still-open}

What remains open is no longer algebraic coefficient matching.  The remaining
open problem is the deeper moving-throat realization of already isolated branch
data.  The only live conditional theorem gate for a full conservative $4$PN
statement is
\begin{equation}
\gamma_{\rm quad}^{\rm eff}\stackrel{?}{=}\frac{2G}{5c^5}.
\label{eq:fixed-open-gamma-target}
\end{equation}
Equivalently, in the reduced moving-throat normalization language inherited from
the $2.5$PN package, the open condition is
\begin{equation}
\hat m_0^{\,2}P_0
\stackrel{?}{=}
\frac{54Gc_s^5}{5a_{\rm th}^5c^5},
\label{eq:fixed-open-P0-target}
\end{equation}
with $a_{\rm th}$ the throat-radius scale on the natural branch.  If
\eqref{eq:fixed-open-gamma-target} or, equivalently,
\eqref{eq:fixed-open-P0-target} holds, then the GR hereditary coefficient
follows automatically.  If it fails, the present theorem remains conditional.

This is the main open theorem quantity.  Two further open layers should be kept
conceptually separate from it.  First, the present paper does not yet derive
from a completed moving-throat PDE the actual grouped-lane conservative and
outgoing data on the true branch.  In the grouped-real language, those still
open branch quantities are the conservative operator moments
\begin{equation}
D_A^{\rm(cons)}(\omega)
=
D_{A,0}+D_{A,2}\omega^2+D_{A,4}\omega^4+O(\omega^6),
\label{eq:fixed-open-grouped-operator}
\end{equation}
together with the outgoing-transfer moments
\begin{equation}
N_A(\omega)
=
N_{A,0}+N_{A,2}\omega^2+N_{A,4}\omega^4+O(\omega^6),
\label{eq:fixed-open-grouped-transfer}
\end{equation}
for $A\in\{20,21,22\}$.  On the natural isotropic branch, these data must
collapse to the scalar prefactor $P_0=N_0/D_0$ that appears in
\eqref{eq:fixed-open-P0-target}.  So the remaining PDE task is not to search
for a new $4$PN coefficient family.  It is to compute one already-isolated
quadrupole branch and evaluate one already-isolated invariant normalization
product.

Second, a few objects remain open only at the level of deeper microscopic
interpretation rather than at the level of coefficient matching.  The present
paper does not yet provide a first-principles moving-throat derivation of the
conservative grouped real quadrupole data beneath
\eqref{eq:fixed-open-grouped-operator}; it does not yet prove a final
uniqueness/interpretation theorem for every aligned-seed local representative;
and it does not attempt any dissipative or radiative completion beyond the
already narrowed quadrupole route.  These are genuine open problems, but they
are no longer algebraic obstacles to the local $4$PN theorem chain.

Finally, there is a practical assembly distinction that should remain visible.
Once the normalization target is fixed on the actual branch, the normalized
hereditary module still has to be inserted into a fully assembled release-grade
referee/master package.  That remaining step is an end-to-end assembly task, not
a new coefficient-discovery problem.

\subsection{Interpretation}
\label{sec:fixed-open-interpretation}

The correct reading rule is now short.  What should be treated as fixed
carry-forward $4$PN data are the one-body repair ledger
\eqref{eq:fixed-open-onebody-ledger}, the quartic compiler, the Hamiltonian
chart lift and sign-flip theorem \eqref{eq:fixed-open-signflip}, the structured
aligned-seed lift \eqref{eq:fixed-open-structured-seed-space}, the existence of
an exact generic-frame representative for the full local instantaneous sector,
and the hereditary bridge \eqref{eq:fixed-open-tail-bridge}.  What should \emph{not}
be treated as fixed yet is the final passive/outgoing quadrupole normalization on
the actual moving-throat branch.

So the main status verdict of the whole paper may be summarized in one sentence:
\begin{center}
\fbox{%
\parbox{0.9\linewidth}{%
The only live conditional interface left at $4$PN is the already-known
passive/outgoing STF quadrupole-normalization gate,
\(\gamma_{\rm quad}^{\rm eff}=2G/(5c^5)\), equivalently
\(\hat m_0^{\,2}P_0=54Gc_s^5/(5a_{\rm th}^5c^5)\).}}
\end{center}

That is the practical meaning of the present theorem package.  The local
instantaneous $4$PN problem is no longer open.  The hereditary side no longer
opens any new datum.  The remaining open work is concentrated in one already
identified branch-normalization theorem that is inherited from the $2.5$PN
program rather than created afresh by $4$PN.

% Section: Verification and reproducibility ledger

\section{Verification and reproducibility ledger}
\label{sec:verification}

The derivation in this paper is analytic rather than computational.  The
accompanying symbolic archive is therefore not presented as a black-box generator
of the theorem.  Its role is narrower and more auditable: it serves as a referee
ledger for the exact one-body gates, compiler identities, local target imports,
generic-frame lift statements, aligned-seed corrections, referee reconstructions,
and hereditary bridge formulas used in the main text.  The archive verifies the
paper as written \emph{within the declared closure hierarchy}; it does not
replace any analytic step, and it does not by itself upgrade the present work
into an assumption-free theorem of a fully solved moving-throat bulk PDE.

At $4$PN the verification story is naturally split into three layers.  The
first layer is the inherited lower-order archive carried forward from the $4$D,
$1$PN, $2$PN, $2.5$PN, and $3$PN packages, which freezes the reduction
backbone, notation firewall, and already solved coefficient ledger used
unchanged here.  The second layer is the dedicated dual-CAS $4$PN referee
suite, consisting of a stagewise SymPy ledger together with a grouped
Mathematica mirror.  Its job is to check the new theorem chain specific to the
present paper: the strict one-body gate, the quartic Legendre compiler, the
local scaffold and fixed-chart import, the Hamiltonian-first generic-frame
lift, the aligned-seed correction, the local referee reconstruction, and the
hereditary/tail bridge.  The third layer is the conditional end-to-end replay,
which ties the already-frozen $2.5$PN quadrupole narrowing to the exact local
$4$PN closure and the hereditary bridge, and thereby verifies that the only
remaining $4$PN interface is the same passive/outgoing STF quadrupole
normalization already isolated on the $2.5$PN side.

\subsection{Role of the \texorpdfstring{$4$}{4}PN referee suite}
\label{sec:verification-role}

The present referee suite should be read as a symbolic reproducibility archive,
not as an independent source of physical assumptions.  Its function is to make
each load-bearing algebraic step re-runnable in a transparent way.  In
particular, the suite is designed to verify the statements that matter most to
a referee reading the paper as a theorem chain:
\begin{itemize}
  \item whether the exact Schwarzschild one-body gate really forces the $4$PN
  repair ledger
  \(
  \mu_{\rho,4}=1/8,
  \ d_4=205/16,
  \ s_{34}=-15/32,
  \ s_{26}=-1/16
  \);
  \item whether the quartic perturbative Legendre-transform identity used to
  move between ordinary and Hamiltonian descriptions is exact;
  \item whether the imported local ADM target, the Hamiltonian-first lift, and
  the canonical comparable-mass slice really close the full local instantaneous
  sector;
  \item whether the aligned ordinary seed correction is generically liftable in
  the structured seed spaces claimed in the main text; and
  \item whether the hereditary coefficient really reduces to the single bridge
  relation
  \begin{equation}
  C_{\rm tail}=\frac{GM}{2c^3}\,\gamma_{\rm quad}^{\rm eff},
  \label{eq:verification-tail-bridge}
  \end{equation}
  so that the only remaining $4$PN gate is the already-known quadrupole
  normalization.
\end{itemize}

So the archive is not being asked to prove a different theorem from the one
stated in the paper.  It is being asked to verify that the paper's own exact,
reduced, and closure-level statements close where they are claimed to close.

\subsection{Core artifacts in the present package}
\label{sec:verification-artifacts}

In the published supplementary archive, the dedicated $4$PN package is organized
around a stagewise SymPy ledger together with a grouped Mathematica mirror.  The
SymPy side carries the full fine-grained theorem chain, while the Mathematica
side independently rechecks the same load-bearing gates in a smaller set of
bundle-level scripts.

\begin{table}[H]
  \centering
  \caption{Core $4$PN referee artifacts and their intended roles in the theorem chain.}
  \label{tab:verification-4pn-artifacts}
  \small
  \renewcommand{\arraystretch}{1.15}
  \begin{tabularx}{\linewidth}{@{}L X@{}}
    \toprule
    Artifact & Intended role in the $4$PN chain \\
    \midrule
    \nolinkurl{4pn_onebody_audit.py} & Exact one-body $4$PN Schwarzschild gate and repair ledger. \\
    \nolinkurl{4pn_quartic_legendre_audit.py} & Exact quartic Legendre compiler through order $c^{-8}$. \\
    \nolinkurl{4pn_local_scaffold_audit.py} & Raw local generic-frame scaffold, COM completeness, and pivot-spanning local basis. \\
    \nolinkurl{4pn_local_target_import_audit.py} & Fixed-chart local ADM target import and the blockwise local degree-ceiling audit. \\
    \nolinkurl{4pn_local_hamiltonian_to_ordinary_audit.py} & Reduced COM Hamiltonian-to-ordinary translation of the fixed local target. \\
    \nolinkurl{4pn_hamiltonian_chart_generic_frame_lift_audit.py} & Hamiltonian-chart generic-frame lift of the local comparable-mass residual. \\
    \nolinkurl{4pn_hamiltonian_chart_canonical_slice_audit.py} & Canonical comparable-mass residual slice in the Hamiltonian chart. \\
    \nolinkurl{4pn_generic_frame_ordinary_translation_audit.py} & Ordinary-chart residual translation and aligned-seed setup. \\
    \nolinkurl{4pn_generic_frame_aligned_seed_lift_audit.py} & Generic-frame lift of the aligned-seed ordinary correction. \\
    \nolinkurl{4pn_local_referee_master_sympy_audit.py} & End-to-end local referee reconstruction of the full fixed-chart ordinary target. \\
    \nolinkurl{4pn_tail_bridge_audit.py} & Tail-side bridge setup and universal quadrupole local shadow. \\
    \nolinkurl{4pn_tail_hereditary_bridge_audit.py} & Exact hereditary coefficient bridge to the canonically normalized STF quadrupole branch. \\
    \nolinkurl{4pn_conditional_referee_master_audit.py} & Conditional full conservative $4$PN replay tying the local closure to the inherited $2.5$PN quadrupole-normalization gate. \\
    \bottomrule
  \end{tabularx}
\end{table}

The package should be read together with the carried lower-order referee archive.
In particular, the final conditional replay imports the already-frozen $2.5$PN
master audit rather than rederiving the quadrupole narrowing from scratch,
because the point of the $4$PN package is precisely to show that no new
independent normalization datum opens beyond that carried interface.

The grouped Mathematica mirror lives beside the stagewise SymPy archive and is
organized around the seven scripts
\nolinkurl{4pn_onebody_mathematica_audit.wl},
\nolinkurl{4pn_quartic_legendre_mathematica_audit.wl},
\nolinkurl{4pn_local_scaffold_target_mathematica_audit.wl},
\nolinkurl{4pn_hamiltonian_chain_mathematica_audit.wl},
\nolinkurl{4pn_ordinary_chain_mathematica_audit.wl},
\nolinkurl{4pn_tail_bridge_mathematica_audit.wl}, and
\nolinkurl{4pn_conditional_referee_master_mathematica_audit.wl}.  These files do
not replace the finer stage structure on the SymPy side; their role is to
provide an independent second-engine replay of the same one-body, compiler,
local-chain, hereditary-bridge, and conditional-closure statements.

\subsection{Stage audits replayed by the master runner}
\label{sec:verification-stage-audits}

The stage structure mirrors the theorem chain in the main text and is organized
so that failure in one lane cannot be hidden by compensating changes in another.
At a practical level the archive is best read in four tranches.

\paragraph{One-body and compiler front end.}
The first tranche consists of the one-body and compiler audits.  These scripts
expand the exact Schwarzschild target, solve the minimal $4$PN repair ledger,
and freeze the quartic ordinary/Hamiltonian compiler that every later local
comparison depends on.

\paragraph{Local scaffold and fixed-chart import.}
The second tranche consists of the raw local scaffold, target-import, and
reduced translation audits.  Their job is to verify that the raw
exchange-symmetric local scaffold is already COM-complete block by block, that
the fixed local ADM target is imported in the correct chart, and that the
Hamiltonian chart rather than the translated ordinary chart carries the correct
local degree ceilings.

\paragraph{Generic-frame local reconstruction.}
The third tranche consists of the Hamiltonian-chart lift, canonical slice,
ordinary translation, aligned-seed lift, and local referee master audits.
Together they verify the full Hamiltonian-first local theorem chain: existence
of a generic-frame Hamiltonian representative, exact ordinary sign flip,
structured liftability of the aligned-seed correction, and final reconstruction
of one exact generic-frame ordinary representative whose COM reduction
reproduces the complete fixed-chart local ordinary target.

\paragraph{Hereditary and conditional interface lane.}
The final tranche consists of the tail bridge, hereditary bridge, and
conditional master replay.  This lane verifies the universal Newtonian
quadrupole shadow of the GR tail, the exact coefficient bridge
\eqref{eq:verification-tail-bridge}, and the interface theorem stating that the
full conservative $4$PN theorem is conditional only on the same
passive/outgoing STF quadrupole normalization already isolated by the $2.5$PN
program.

This modular staging is intentional.  It keeps the one-body gate, the local
Hamiltonian-first lift, the aligned-seed repair, and the hereditary bridge
separate, so that a failure in one sector cannot be disguised as a different
kind of local or tail compensation.

\subsection{What the suite verifies}
\label{sec:verification-verifies}

The suite verifies the load-bearing symbolic statements of the paper sector by
sector.

First, the early audits verify the exact one-body and compiler backbone.  They
check the strict Schwarzschild $4$PN target, confirm that the carried $3$PN
self-sector packaging is insufficient by itself, solve the minimal repair
ledger, and then prove the quartic perturbative Legendre-transform identity
needed to move cleanly between ordinary and Hamiltonian descriptions at $4$PN.

Second, the local middle audits verify the entire comparable-mass chart-level
closure chain.  They check the exact local scaffold and COM completeness, the
fixed local target import, the Hamiltonian-first degree-ceiling lesson, the
Hamiltonian-chart generic-frame lift, the canonical residual slice, the exact
ordinary sign-flip rule
\begin{equation}
\Delta H_4=-\Delta L_4,
\label{eq:verification-signflip-4pn}
\end{equation}
and the structured lift of the aligned-seed ordinary correction.  In the
language of the main text, this is the stage chain that verifies that the
entire local instantaneous $4$PN sector has an exact generic-frame ordinary
representative inside the declared hierarchy.

Third, the local referee master verifies the end-to-end local assembly itself.
It rebuilds the exact reduced COM local ordinary target from the imported local
ADM Hamiltonian, verifies the one-body gate and natural seed, replays the
canonical residual translation and aligned-seed correction, and confirms that
the assembled generic-frame candidate reduces back to the full fixed-chart local
ordinary target slot by slot.  This is the direct executable check of the main
local theorem claim.

Fourth, the tail-side audits verify the narrow hereditary bridge.  They confirm
that the GR tail contributes a universal Newtonian quadrupole shadow, and that
on the canonically normalized STF quadrupole branch the coefficient is exactly
\begin{equation}
C_{\rm tail}=\frac{GM}{2c^3}\,\gamma_{\rm quad}^{\rm eff}.
\label{eq:verification-tail-bridge-repeat}
\end{equation}
So the hereditary side is not an independent $4$PN coefficient hunt.  It is an
interface theorem tying the conservative hereditary coefficient to the same
quadrupole normalization that already controls the $2.5$PN odd sector.

Finally, the conditional master replay verifies the paper's sharpest global
status statement: once the inherited $2.5$PN narrowing, the exact local $4$PN
closure, and the hereditary bridge are replayed together, the only remaining
full-$4$PN gate is the already-known passive/outgoing quadrupole normalization
condition,
\begin{equation}
\gamma_{\rm quad}^{\rm eff}\stackrel{?}{=}\frac{2G}{5c^5},
\label{eq:verification-gamma-target}
\end{equation}
or, equivalently on the moving-throat side,
\begin{equation}
\hat m_0^{\,2}P_0\stackrel{?}{=}\frac{54Gc_s^5}{5a_{\rm th}^5c^5}.
\label{eq:verification-P0-target}
\end{equation}
The executable archive therefore verifies the same conclusion as the analytic
paper: no new independent normalization datum opens at $4$PN.

Two reproducibility conventions are worth recording explicitly.  First, the
archive is stagewise rather than monolithic, so each load-bearing symbolic claim
is attached to a named audit rather than hidden inside one opaque master file.
Second, the local and hereditary lanes are kept separate all the way to the
conditional replay, which makes the theorem split
\(
L_{4\mathrm{PN}}^{\mathrm{cons}}=L_{4\mathrm{PN}}^{\mathrm{local}}+L_{4\mathrm{PN}}^{\mathrm{tail}}
\)
visible in the executable package itself rather than only in the prose of the
paper.

\subsection{What the suite does not verify}
\label{sec:verification-does-not-verify}

The suite does \emph{not} replace the analytic argument in the main text, and it
does not solve the full moving-throat PDE.  In particular, the archive does not
by itself derive the passive/outgoing STF quadrupole normalization on the actual
moving-throat branch; it does not prove from first principles that
\eqref{eq:verification-gamma-target} or \eqref{eq:verification-P0-target}
holds; and it does not upgrade the present work into an assumption-free theorem
of the full bulk dynamics.

It also does not attempt to settle every deeper microscopic interpretation issue
that remains outside the present $4$PN scope.  The archive is not a proof of the
full moving-throat origin of all grouped conservative data, it is not a final
uniqueness theorem for every aligned-seed local representative, and it is not a
general radiative or higher-PN completion package.  Those questions remain
outside the stated theorem target of this paper.

So the correct reading rule is simple.  The referee archive verifies that the
paper's exact local theorem chain and hereditary bridge close where claimed, and
that the whole conservative $4$PN problem has been reduced to one already-known
quadrupole-normalization gate.  What it does not verify is the final PDE-side
realization of that gate.  That distinction is exactly the one on which the
conditional status of the present $4$PN theorem rests.

% Section: Conclusions

\section{Conclusions}
\label{sec:conclusions}

This paper closes the conservative two-body $4$PN problem in the same precise
sense that the earlier full papers closed the conservative $1$PN, $2$PN, and
$3$PN sectors: not as an assumption-free theorem of a fully solved
moving-throat PDE, but as an exact algebraic assembly within a declared closure
hierarchy.  The gain is larger than one more coefficient list.  At $4$PN the
entire local instantaneous sector is now fixed exactly, while the hereditary
sector has collapsed to a single scalar bridge to the same canonically
normalized passive/outgoing STF quadrupole branch that had already been
isolated by the $2.5$PN program.  In that sharpened sense, the conservative
$4$PN bottleneck is no longer an unresolved coefficient search.  It is one
already-known normalization interface.

\subsection{What has been established}
\label{sec:conclusions-established}

The proof line established in the paper may be stated compactly.

\begin{enumerate}
  \item The strict one-body Schwarzschild gate through $c^{-8}$ is fixed
  exactly, and its minimal repair ledger is not one-dimensional but
  \begin{equation}
  \mu_{\rho,4}=\frac18,
  \qquad
  d_4=\frac{205}{16},
  \qquad
  s_{34}=-\frac{15}{32},
  \qquad
  s_{26}=-\frac1{16}.
  \label{eq:conclusions-onebody-ledger-4pn}
  \end{equation}
  So the carried $3$PN self-sector package does not extend to $4$PN by a single
  denominator coefficient or a single static gate alone.

  \item The exact quartic ordinary/Hamiltonian compiler is rigid, and once the
  lower-order ledger is frozen the correct comparable-mass lift is
  Hamiltonian-first.  In the local chart this gives the exact sign-flip rule
  \begin{equation}
  \Delta H_4=-\Delta L_4,
  \label{eq:conclusions-signflip-4pn}
  \end{equation}
  so the comparable-mass local residual is most naturally fixed on the
  Hamiltonian side and only translated back to the ordinary chart afterward.

  \item In that Hamiltonian chart there is no local span obstruction.  The raw
  exchange-symmetric generic-frame scaffold is already complete against the full
  fixed-chart local interaction target, the canonical comparable-mass residual
  is exactly reachable, and the only ordinary-chart obstruction is an
  aligned-seed issue rather than a failure of the local generic-frame lift
  itself.

  \item That aligned-seed obstruction is not a new free datum.  It is exactly
  liftable in the structured seed spaces
  \begin{equation}
  Q,
  \qquad
  T\oplus(pq)T,
  \qquad
  S\oplus(pq)S,
  \qquad
  U\oplus(p^2q^2)U,
  \qquad
  W=0,
  \label{eq:conclusions-structured-seed-spaces-4pn}
  \end{equation}
  so one exact generic-frame ordinary representative of the full local
  instantaneous sector exists inside the declared hierarchy.

  \item Consequently, the local instantaneous $4$PN theorem is already exact in
  the form
  \begin{equation}
  L_{4\mathrm{PN}}^{\mathrm{local}}
  =
  L_{4,\mathrm{seed}}^{(\mathrm{natural,gen})}
  +
  \delta L_{4,\mathrm{seed}}^{(\mathrm{gen})}
  +
  \Delta L_{4,\mathrm{loc}}^{(\mathrm{can,gen})},
  \qquad
  \Pi_{\rm COM}\!\left[L_{4\mathrm{PN}}^{\mathrm{local}}\right]
  =
  L_{4,\mathrm{loc}}^{(\mathrm{target,ord})}.
  \label{eq:conclusions-local-theorem-4pn}
  \end{equation}
  So the local instantaneous $4$PN bottleneck is no longer algebraic span,
  chart ambiguity, or generic-frame existence.

  \item The hereditary side is now equally sharp.  Its unique compatible
  coefficient is
  \begin{equation}
  C_{\rm tail}
  =
  \frac{GM}{2c^3}\,\gamma_{\rm quad}^{\rm eff},
  \label{eq:conclusions-tail-bridge-4pn}
  \end{equation}
  where $\gamma_{\rm quad}^{\rm eff}$ is the canonically normalized
  passive/outgoing STF quadrupole coefficient already isolated by the $2.5$PN
  program.  Therefore $4$PN opens no new independent normalization datum.
\end{enumerate}

When these facts are assembled, the full conservative $4$PN theorem takes the
exact split
\begin{equation}
\boxed{
L_{4\mathrm{PN}}^{\mathrm{cons}}
=
L_{4\mathrm{PN}}^{\mathrm{local}}
+
L_{4\mathrm{PN}}^{\mathrm{tail}},
\qquad
C_{\rm tail}
=
\frac{GM}{2c^3}\,\gamma_{\rm quad}^{\rm eff}.
}
\label{eq:conclusions-final-split-4pn}
\end{equation}
Full conservative GR recovery is therefore equivalent to the single scalar
condition
\begin{equation}
\gamma_{\rm quad}^{\rm eff}=\frac{2G}{5c^5},
\label{eq:conclusions-gamma-target-4pn}
\end{equation}
which immediately yields
\begin{equation}
C_{\rm tail}=\frac{G^2M}{5c^8}=C_{\rm tail}^{\rm GR}.
\label{eq:conclusions-tail-gr-4pn}
\end{equation}
Conceptually, the paper does three things at once.  First, it closes the entire
local instantaneous $4$PN sector exactly inside the declared hierarchy.
Second, it proves that the hereditary side is not a second independent
higher-order search, but a scalar bridge to the same canonically normalized STF
quadrupole branch already identified by the $2.5$PN audit.  Third, it sharpens
the bookkeeping of what is fixed, what is closure-level, and what still has to
be supplied by the completed moving-throat PDE.  The accompanying referee
archive supports this reading in an auditable way: the one-body, compiler,
local-lift, aligned-seed, referee-master, and tail-bridge audits collectively
verify that the exact theorem chain closes where claimed inside the same
hierarchy.

\subsection{What remains open}
\label{sec:conclusions-open}

The negative statement is equally important.

First, the paper does not solve the fully dynamical moving-throat PDE and then
recover the $4$PN theorem as an automatic bulk consequence.  The present result
remains a full conservative derivation only within the declared reduction and
closure hierarchy already used by the lower-order papers.  That status is not a
weakness of the theorem; it is the correct statement of its scope.

Second, the paper does not yet derive from first principles the passive/
outgoing STF quadrupole normalization on the actual moving-throat branch.  The
remaining universal target may be stated in either of the equivalent forms
\begin{equation}
\boxed{
\gamma_{\rm quad}^{\rm eff}=\frac{2G}{5c^5},
\qquad
\hat m_0^{\,2}P_0=\frac{54Gc_s^5}{5a_{\rm th}^5c^5}.
}
\label{eq:conclusions-normalization-target-4pn}
\end{equation}
The present $4$PN theorem removes the conservative algebraic uncertainty
\emph{beneath} this bridge, but it does not by itself prove that the physical
moving-throat branch realizes \eqref{eq:conclusions-normalization-target-4pn}.

Third, the deeper microscopic derivation of the grouped conservative quadrupole
data also remains open.  On the moving-throat side, the relevant grouped bundle
still has to be computed as an actual branch output,
\begin{equation}
D_A(\omega)=D_{A0}+D_{A2}\omega^2+D_{A4}\omega^4+O(\omega^6),
\qquad
N_A(\omega)=N_{A0}+N_{A2}\omega^2+N_{A4}\omega^4+O(\omega^6),
\label{eq:conclusions-grouped-bundle-open}
\end{equation}
with $A\in\{20,21,22\}$ and with the isotropic branch determined rather than
assumed.  These quantities are not leftover algebraic residues in the present
$4$PN theorem.  They are dynamical response data of the microscopic
mouth/support/mixed-sector system.

Fourth, the paper does not yet prove a final uniqueness or interpretation theorem
for every aligned-seed local representative.  What has been proved is the exact
existence of a full local generic-frame representative and the structured
liftability of the aligned-seed correction.  That is sufficient for the present
conservative theorem, but it is not the last possible microscopic word on the
ordinary-chart seed organization.

Finally, larger sectors remain outside the honest scope of the present paper:
dissipative leakage and radiation reaction beyond the already narrowed
quadrupole route, spin couplings, strong-field completion, and higher-PN work
beyond the solved conservative $4$PN scope.  What remains open is therefore no
longer a local coefficient-matching problem.  It is the deeper moving-throat
realization of data that are already algebraically assigned.

\subsection{Interface to the next calculation}
\label{sec:conclusions-next}

This gives the project a much cleaner handoff than before.  The next job is no
longer to reopen the local $4$PN generic-frame algebra, to search for another
hidden seed parameter, or to treat the hereditary sector as a fresh independent
normalization family.  That work is substantially finished in the present
chart.  The next job is the moving-throat grouped-quadrupole calculation
itself.

On the conservative side, the physical branch must first determine whether the
natural grouped real $P_2$ isotropy gate actually holds,
\begin{equation}
\boxed{a_2=0,
\qquad
b_2=0,}
\label{eq:conclusions-isotropy-gate-4pn}
\end{equation}
and, equivalently on the outgoing side,
\begin{equation}
\boxed{a(P_0)=0,
\qquad
b(P_0)=0.}
\label{eq:conclusions-prefactor-isotropy-gate-4pn}
\end{equation}
If these first gates fail, then the minimal isotropic passive/outgoing
quadrupole branch is already dead before the final normalization test is even
attempted.  If they pass, then the next task is simply to evaluate the single
invariant normalization product in
\eqref{eq:conclusions-normalization-target-4pn}.

Equivalently, the completed moving-throat PDE must determine the grouped-bundle
coefficients in \eqref{eq:conclusions-grouped-bundle-open} together with the
source-map factor $\hat m_0$, collapse them to the isotropic branch if that
branch is realized, and then test whether
\begin{equation}
\hat m_0^{\,2}P_0=\frac{54Gc_s^5}{5a_{\rm th}^5c^5}.
\label{eq:conclusions-final-pde-gate-4pn}
\end{equation}
At that point the odd $2.5$PN Burke--Thorne coefficient and the conservative
$4$PN hereditary coefficient close together; before that point, both remain
conditional on the same microscopic normalization theorem.

So the correct next calculation is sharply defined.  One must compute the
canonically normalized grouped quadrupole data on the physical moving-throat
branch, decide whether the isotropic passive/outgoing branch is actually
realized there, and then evaluate the single invariant normalization product in
\eqref{eq:conclusions-final-pde-gate-4pn}.  That is the only place where the
present theorem still needs to be sharpened from a closure-level conservative
result into a fully microscopic one.

The strongest honest endpoint of the paper is therefore also the strongest
honest starting point for the next one: within the declared hierarchy, the
local instantaneous $4$PN bottleneck is gone, the hereditary bridge is exact,
and the only serious universal question left open is the passive/outgoing
normalization of the already identified STF quadrupole branch.

\section*{Acknowledgments}
This work was carried out by an independent researcher.
Generative-AI tools were used during the development of the manuscript to help
translate conceptual ideas into mathematical form, assist with symbolic
manipulations, prototype and revise companion Mathematica and SymPy
verification scripts, and improve the exposition.
The conceptual direction, hypotheses, and overall structure of the work were
directed by the author.
The resulting arguments and calculations were checked for internal consistency
using the accompanying symbolic scripts, but the work should be regarded as
exploratory and would benefit from review by subject-matter experts.
Any errors, omissions, or misunderstandings remain the author's.


% ---------------------------------------------------------------------------
% Appendices
% ---------------------------------------------------------------------------

\appendix

% Appendix: Carry-forward formulas from the earlier papers

\section{Carry-forward formulas from the earlier papers}
\label{app:carry-forward}

This appendix records the specific formulas and status statements imported from
 the parent $4$D paper and the full conservative $1$PN, $2$PN, $2.5$PN, and
 $3$PN packages and used without rederivation in the present $4$PN work.  Its
 role is technical rather than historical.  We do not repeat the full parent
 action, the full lower-order proofs, or the longer appendix-side derivations
 that fixed the wake, breathing, support, grouped-quadrupole, and
 geometry-completion data.  Instead, we collect only the load-bearing formulas
 that the present $4$PN paper actually uses: the exact parent reduction
 identities, the frozen Newtonian+$1$PN+$2$PN ledger, the $2.5$PN narrowing to
 the canonically normalized STF quadrupole branch, and the $3$PN front end that
 the $4$PN one-body and local theorem chains continue.

The claim-status boundary is the same one used throughout the earlier papers.
Exact projection identities from the parent $4$D theory are kept visibly exact.
The near-zone Poisson hook, the zero-mode Maxwell law, and the worldtube
 point-particle limit are carried only as controlled reductions.  The
 conservative Newtonian+$1$PN+$2$PN+$3$PN ledgers are treated as frozen outputs
 of the earlier program within its declared closure hierarchy, not as live
 parameters to be retuned in the present paper.  The $2.5$PN package is used in
 the same way: not as a finished moving-throat PDE theorem, but as the theorem
 audit that narrowed the surviving higher-order universal route to the
 orbital/worldtube STF quadrupole branch.

\subsection{Parent \texorpdfstring{$4$}{4}D reduction dictionary and Newtonian+\texorpdfstring{$1$}{1}PN carried data}
\label{app:carry-forward-1pn}

The parent theory lives on a $(4+1)$-dimensional spacetime with coordinates
\[
X^M=(t,x,y,z,w),
\qquad
\mathbf{X}=(x,y,z,w)\in\mathbb R^4,
\qquad
\mathbf{x}=(x,y,z)\in\mathbb R^3.
\]
The bulk variables carried into the present paper are the matter field
$\psi(\mathbf X,t)$, its density
\[
\rho(\mathbf X,t)=|\psi(\mathbf X,t)|^2,
\]
the optional localized gauge field $A_M(\mathbf X,t)$ with localization profile
$Z(w)$, and the effective throat geometry variables
\[
a(t),\qquad L(t).
\]
When electric charge is named explicitly in this downstream notation, it is
always the corrected branch quantity
\[
q_* = \eta_Q e_*,
\qquad
q_{\rm eff}=\frac{q_*}{\sqrt{Z_{\rm int}}},
\qquad
Z_{\rm int}=\int_{-\infty}^{\infty} Z(w)\,dw,
\]
not the gravity-side scalar coefficient $\kappa_\rho$ that appears in the
post-Newtonian ledger below.

Brane observables are defined by a fixed normalized projection kernel $W(w)$,
\[
\int_{-\infty}^{\infty} W(w)\,dw = 1,
\qquad
\mathcal P_W[Q](t,\mathbf x)
\equiv
\int_{-\infty}^{\infty} W(w)\,Q(t,\mathbf x,w)\,dw.
\]
In particular,
\[
\rho_{\rm br}=\mathcal P_W[\rho],
\qquad
\mathbf J_{\rm br}=\mathcal P_W[(j^x,j^y,j^z)],
\qquad
\mathbf v_{\rm br}=\frac{\mathbf J_{\rm br}}{\rho_{\rm br}}
\quad (\rho_{\rm br}>0).
\]
Projection is not the same thing as reduction.  Projection defines what a brane
 observer measures; reduction refers to the further elimination of the extra
 coordinate under an explicit ansatz or scaling regime.

From exact bulk continuity,
\[
\partial_t\rho+\partial_i j^i=0,
\qquad i\in\{x,y,z,w\},
\]
one obtains the exact open-system brane continuity law
\[
\partial_t\rho_{\rm br}+\nabla_3\!\cdot\mathbf J_{\rm br}=S_{\rm leak},
\]
with leakage source
\[
S_{\rm leak}
=
-\bigl[W(w)j^w\bigr]_{-\infty}^{+\infty}
+
\int_{-\infty}^{\infty} W'(w)j^w\,dw.
\]
After Helmholtz decomposition of the brane velocity,
\[
\mathbf v_{\rm br}=\nabla_3\phi_{\rm br}+\mathbf v_T,
\qquad
\nabla_3\!\cdot\mathbf v_T=0,
\]
projection continuity becomes the exact longitudinal identity
\[
\rho_{\rm br}\,\nabla_3^2\phi_{\rm br}
=
S_{\rm leak}
-\partial_t\rho_{\rm br}
-(\nabla_3\rho_{\rm br})\cdot(\nabla_3\phi_{\rm br}+\mathbf v_T).
\]
Under the usual quasi-static near-zone simplifications this becomes the carried
 Poisson hook
\[
\nabla_3^2\phi_{\rm br}\simeq \frac{S_{\rm eff}}{\rho_0},
\qquad
\Phi_N=\lambda_\Phi\,\phi_{\rm br},
\]
which is the origin of the Newtonian scalar sector used in the lower-order
 particle reduction.  The corresponding coherent-defect / small-body worldtube
 limit gives the point-particle force law
\[
M_A\,\ddot{\mathbf X}_A
=
-M_A\nabla\Phi_{\rm ext}(\mathbf X_A)
+O\!\left(\frac{a_{\rm th}^2}{r^2}\,\frac{\Phi_{\rm ext}}{r}\right),
\]
with the first finite-size correction quadrupolar and therefore suppressed in
 the controlled small-worldtube regime.  For completeness, the same parent paper
 also carries a controlled zero-mode Maxwell reduction.  Under the far-field
 brane assumptions
\[
A_w\approx 0,
\qquad
\partial_w A_\mu\approx 0,
\qquad
J^w\approx 0,
\qquad
F_{\mu w}\approx 0,
\]
integration over $w$ gives the effective brane Maxwell law
\[
\partial_\mu F^{\mu\nu}=\mu_0^{\rm eff}J_{\rm eff}^\nu,
\qquad
\mu_0^{\rm eff}=\frac{\mu_0}{Z_{\rm int}}.
\]
This reduction suppresses the mixed channels only in the controlled far-field
 limit.  It does not erase $(A_w,J^w,F_{\mu w})$ from the microscopic ontology.

The Newtonian point-particle block carried into every later paper is therefore
\[
L_0
=
\frac12 m_A v_A^2
+\frac12 m_B v_B^2
+\frac{Gm_A m_B}{r}.
\]
At $1$PN the scalar/optical one-body reduction fixes the lower-order viability
 conditions
\[
\kappa_\rho=1,
\qquad
n=5,
\qquad
\kappa_{\rm add}=\frac12,
\qquad
\kappa_{\rm PV}=\frac32,
\qquad
\beta_{\rm 1PN}=3.
\]
In the parity-even wake sector, the constructive isotropic Fourier-space closure
 freezes
\[
\alpha^2=\frac34,
\qquad
a_H=0,
\qquad
K_{\rm vec}=\frac{2}{\pi^2},
\]
which in turn give the EIH cross coefficients
\[
C_\parallel=-\frac72,
\qquad
C_L=-\frac12.
\]
With these values frozen, the full carried Newtonian+$1$PN two-body block is
\[
\begin{aligned}
L_1=
&\frac18(m_Av_A^4+m_Bv_B^4)
+\frac{Gm_A m_B}{r}
\left[
\frac32(v_A^2+v_B^2)
-\frac72 v_{AB}
-\frac12 v_{An}v_{Bn}
\right]
\\
&-\frac{G^2m_A m_B(m_A+m_B)}{2r^2}.
\end{aligned}
\]
Nothing in the present $4$PN paper is allowed to retune these lower-order
 values.  They are carry-forward anchors rather than open coefficients.

\subsection{\texorpdfstring{$2$}{2}PN carried data and the \texorpdfstring{$P_0\oplus P_2$}{P0 + P2} support hierarchy}
\label{app:carry-forward-2pn}

The full conservative $2$PN paper closes the comparable-mass $c^{-4}$ ledger
 against the standard generic-frame ADM target within the same declared closure
 hierarchy, but the present $4$PN paper imports from it only the load-bearing
 formulas and the response-side structure that sharpened the ontology of the
 defect.  At the strict one-body level, the exact isotropic Schwarzschild gate
 through $2$PN is packaged by
\[
D_{\rm target}(u)=1-4u+8u^2+O(u^3),
\qquad
u\equiv\frac{U}{c^2},
\]
and the pure-static $2$PN gate fixes
\[
\lambda_\rho=\frac12,
\qquad
L_{2,\rm static}=\frac{G^3m_A m_B(m_A^2+m_B^2)}{4r^3}.
\]
The full carried $2$PN coefficient inserted into the Newtonian+$1$PN ledger is
\[
\begin{aligned}
L_{2,\rm full}
={}&\frac{m_Av_A^6+m_Bv_B^6}{16}
\\
&+\frac{Gm_A m_B}{r}
\Big[
\frac78(v_A^4+v_B^4)
-\frac74 v_{AB}(v_A^2+v_B^2)
-\frac14 v_{An}v_{Bn}(v_A^2+v_B^2)
\\
&\hspace{4em}
+\frac{11}{8}v_A^2v_B^2
+\frac14 v_{AB}^2
-\frac58(v_A^2v_{Bn}^2+v_B^2v_{An}^2)
+\frac32 v_{AB}v_{An}v_{Bn}
+\frac38 v_{An}^2v_{Bn}^2
\Big]
\\
&+\frac{G^2m_A m_B}{r^2}
\Big[
\Big(2m_B+\frac{11}{8}m_A\Big)v_A^2
+\Big(2m_A+\frac{11}{8}m_B\Big)v_B^2
\\
&\hspace{4em}
-\frac{15}{4}(m_A+m_B)v_{AB}
+\frac{15}{8}(m_Av_{An}^2+m_Bv_{Bn}^2)
\Big]
\\
&+\frac{G^3m_A m_B}{4r^3}(m_A^2+5m_A m_B+m_B^2).
\end{aligned}
\]
So by the start of the present paper, the full conservative
 Newtonian+$1$PN+$2$PN ledger is already frozen in the declared hierarchy.

Conceptually, the decisive ontology change at $2$PN is that the defect is no
 longer described adequately as a pure scalar monopole.  The conservative cross
 sector already splits into three response lanes:
\begin{enumerate}
  \item a carried odd dipole wake sector,
  \item a genuine even $P_0\oplus P_2$ mouth/support layer,
  \item and a separate geometry-closure channel.
\end{enumerate}
The appendix-side zero-frequency throat data make this explicit.  The odd dipole
 residues are frozen to
\[
R_{1\perp}=\frac72,
\qquad
R_{10}=4,
\]
while the canonical even support residues are
\[
R_0=R_{20}=R_{21}=R_{22}=1.
\]
The same solved static even operator fixes the source-weight vector
\[
J=\left(\frac{4}{\sqrt5},\frac54,0,0,0,0\right),
\qquad
J\cdot J=\frac{381}{80},
\]
and the pure-potential geometry-closure deficit
\[
\Delta_{\rm geom}=\frac{281}{80}.
\]
The monopole wall add-on
\[
\Delta K_{00}=\frac{109}{280}
\]
is kept distinct from $\Delta_{\rm geom}$; they are not the same object and
 should not be merged in later bookkeeping.  What remains open after the
 conservative $2$PN theorem is not these zero-frequency residues but the fully
 dynamical completion.  In particular, the pole or compliance scales
\[
\Omega_{1\perp},\ \Omega_{10},\ \Omega_0,
\ \Omega_{20},\ \Omega_{21},\ \Omega_{22},\ \Omega_g
\]
are not claimed as already derived from the completed moving-throat PDE.

\subsection{\texorpdfstring{$2.5$}{2.5}PN narrowing to the grouped real \texorpdfstring{$P_2$}{P2} quadrupole branch}
\label{app:carry-forward-25pn}

The $2.5$PN theorem audit is what makes the present $4$PN problem sharply
 focused rather than diffuse.  Its decisive benchmark shows that the frozen
 conservative hierarchy by itself cannot generate a nonzero local $2.5$PN
 reaction theorem, while a minimal retarded quadrupole odd term lands on the
 standard Burke--Thorne / Iyer--Will family member
\[
\alpha=4,
\qquad
\beta=5.
\]
For the present conservative $4$PN paper, however, the important carry-forward
 message is narrower: after the scalar, dipole, and quadrupole channel audits,
 the live higher-order universal point-particle route collapses to the
 orbital/worldtube STF quadrupole branch
\[
\boxed{I_{ij}^{\rm orb}=\mu x_{\langle i}x_{j\rangle}.}
\]
The odd-channel hierarchy on the narrow branch is
\[
\text{scalar }\sim i\omega,
\qquad
\text{dipole }\sim i\omega^3,
\qquad
\text{quadrupole }\sim i\omega^5,
\]
and the quadrupole normalization target is
\[
\boxed{\hat m_0^2\Gamma_5=\frac{2G}{5c^5}.}
\]

In grouped real language, the conservative front end is organized as
\[
Y_{20}(\omega)=1+u_2^{(20)}\omega^2+u_4^{(20)}\omega^4+O(\omega^6),
\]
\[
Y_{21}(\omega)=1+u_2^{(21)}\omega^2+u_4^{(21)}\omega^4+O(\omega^6),
\]
\[
Y_{22}(\omega)=1+u_2^{(22)}\omega^2+u_4^{(22)}\omega^4+O(\omega^6).
\]
It is convenient to repackage these grouped triples into trace/anomaly
 variables.  At $O(\omega^2)$ one defines
\[
\bar u_2=\frac{u_2^{(20)}+2u_2^{(21)}+2u_2^{(22)}}{5},
\qquad
a_2=\frac{2u_2^{(20)}-u_2^{(21)}-u_2^{(22)}}{10},
\qquad
b_2=\frac{u_2^{(21)}-u_2^{(22)}}{2},
\]
and at $O(\omega^4)$
\[
\bar u_4=\frac{u_4^{(20)}+2u_4^{(21)}+2u_4^{(22)}}{5},
\qquad
a_4=\frac{2u_4^{(20)}-u_4^{(21)}-u_4^{(22)}}{10},
\qquad
b_4=\frac{u_4^{(21)}-u_4^{(22)}}{2}.
\]
The exact inverse map is
\[
u_2^{(20)}=\bar u_2+4a_2,
\qquad
u_2^{(21)}=\bar u_2-a_2+b_2,
\qquad
u_2^{(22)}=\bar u_2-a_2-b_2,
\]
\[
u_4^{(20)}=\bar u_4+4a_4,
\qquad
u_4^{(21)}=\bar u_4-a_4+b_4,
\qquad
u_4^{(22)}=\bar u_4-a_4-b_4.
\]
So the grouped-real isotropy gate on the minimal branch is simply
\[
a_2=b_2=a_4=b_4=0.
\]
On that branch,
\[
\bar u_4=4\bar u_2^2,
\qquad
\Omega_Q^2=\frac{1}{4\bar u_2},
\qquad
\overline K_0^{\rm target}=\frac{2G}{45c^5\bar u_2^{5/2}}.
\]
For later use, the normalized outgoing $\ell=2$ moment package can be written as
\[
\widehat Y_2(k)=1+\frac{a_{\rm th}^2k^2}{9}+\frac{4a_{\rm th}^4k^4}{81}+i\frac{a_{\rm th}^5k^5}{27}+\cdots,
\]
with the moment identities
\[
m_4=\frac{4m_2^2}{m_0},
\qquad
\Gamma_5^{\rm PDE}=9\frac{m_2^{5/2}}{m_0^{3/2}}.
\]
These formulas are not the present paper's theorem.  Their role here is to make
 explicit why the hereditary $4$PN bridge can be controlled by the same narrow
 quadrupole-normalization problem that already survived the $2.5$PN audit.

\subsection{\texorpdfstring{$3$}{3}PN carried data and the frozen local front end}
\label{app:carry-forward-3pn}

The full conservative $3$PN paper is the immediate lower-order predecessor of
 the present work.  What the $4$PN paper imports from it is not just the fact of
 a lower-order match, but the exact one-body repair ledger, the fixed-chart
 local front end, and the sharpened statement of what remained open after the
 conservative $3$PN bottleneck was removed.

At the strict one-body level, the exact isotropic Schwarzschild target through
 $c^{-6}$ is
\[
\frac{L_{3\mathrm{PN},1\text{-body}}}{m}
=
\frac5{128}v^8
+\frac{11}{16}Uv^6
+\frac{47}{16}U^2v^4
+\frac{13}{8}U^3v^2
-\frac18U^4.
\]
Trying to continue the lower-order denominator strategy with only one new cubic
 denominator datum fails, so the exact one-body repair ledger opened at $3$PN is
\[
\mu_{\rho,3}=\frac14,
\qquad
 d_3=-\frac{45}{4},
\qquad
 s_{24}=-\frac1{16}.
\]
These are exactly the lower-order self-sector values that the present $4$PN
 one-body continuation carries into the repair ansatz.

At the full conservative $3$PN level, the comparable-mass residual in COM form
 is fixed by the split
\[
\Delta \hat L_3^{\rm GR}
=
\underbrace{\Delta l_1v^8}_{\text{compiler image of the free COM Hamiltonian}}
+
\underbrace{L_{P_2}^{\rm mid}}_{\text{exact grouped real }P_2\text{ closure}}
+
\underbrace{\Delta l_{15}^{(g)}U^4}_{\text{unique geometry completion}}.
\]
Within the same declared closure hierarchy as the $1$PN and $2$PN papers, the
 earlier paper also fixed the exact sign and determinant statements
\[
\Delta H_3=-\Delta L_3,
\qquad
\det M_{\rm mid}=-\frac4{27},
\]
so that the grouped-real quadrupole middle block and the unique geometry
 completion are already frozen inputs by the time the present $4$PN work begins.

The carry-forward lesson for the present paper is therefore very specific.
After the full conservative $3$PN assembly, what remained open was no longer
 algebraic $3$PN bookkeeping.  The remaining bridge problem had become: derive
 the grouped real $P_2$ constitutive law and the unique geometry completion
 microscopically from the moving-throat dynamics, then insert that data into the
 already narrowed $2.5$PN outgoing quadrupole-normalization problem.  The
 present $4$PN paper starts exactly from that status boundary.  It carries the
 frozen $3$PN front end, extends the local conservative ledger through $c^{-8}$,
 and then meets the same narrowed quadrupole-normalization interface from the
 conservative side.

Putting the carried formulas together, the lower-order conservative ladder used
 without rederivation in the present paper is
\[
L_{\rm cons}^{(<4\mathrm{PN})}
=
L_0+c^{-2}L_1+c^{-4}L_{2,\rm full}+c^{-6}L_3+O(c^{-8}),
\]
where $L_0$, $L_1$, and $L_{2,\rm full}$ are the explicit blocks recorded
 above, while $L_3$ is the fully assigned fixed-chart conservative $3$PN block
 of the earlier paper.  The present $4$PN theorem chain is not allowed to alter
 any part of this carried ledger.  It begins only after these lower-order data,
 identities, and grouped-branch status statements have already been frozen.

% Appendix: Exact 4PN one-body Schwarzschild expansion

\section{Exact \texorpdfstring{$4$}{4}PN one-body Schwarzschild expansion}
\label{app:onebody-expansion}

This appendix records the exact test-mass Schwarzschild expansion used in
Section~\ref{sec:onebody-schwarzschild}.  Its role is purely foundational.  We
keep the full isotropic worldline series in one place so that the carried
$1$PN, $2$PN, and $3$PN one-body coefficients and the new $4$PN coefficient can
all be read off from one exact formula before any self-sector continuation or
comparable-mass local lift is attempted.

\subsection{Exact isotropic worldline form}
\label{app:onebody-expansion-exact}

On the strict test-mass branch, with
\[
U=\frac{GM}{r},
\qquad
u_{\rm sym}\to 0,
\]
and in isotropic Schwarzschild coordinates, the exact line element is
\begin{equation}
 ds^2
 =
 -\left(\frac{1-u/2}{1+u/2}\right)^2 c^2 dt^2
 +(1+u/2)^4 d\mathbf{x}^2,
 \qquad
 u\equiv \frac{U}{c^2}.
\label{eq:app-4pn-isotropic-schwarzschild-metric}
\end{equation}
Therefore the exact worldline Lagrangian is
\begin{equation}
\frac{L_{\rm exact}}{m}
=
-c^2\sqrt{\left(\frac{1-u/2}{1+u/2}\right)^2-(1+u/2)^4\frac{v^2}{c^2}},
\qquad
v^2\equiv \mathbf v^2.
\label{eq:app-4pn-isotropic-schwarzschild-worldline}
\end{equation}
This is the unique one-body source for the local conservative post-Newtonian
series used throughout the paper.  No self-sector continuation is allowed to
modify the coefficients obtained by expanding
\eqref{eq:app-4pn-isotropic-schwarzschild-worldline}.

\subsection{Expansion through \texorpdfstring{$4$}{4}PN}
\label{app:onebody-expansion-series}

Expanding
\eqref{eq:app-4pn-isotropic-schwarzschild-worldline} in powers of $c^{-2}$ at
fixed $U$ and $v^2$ gives
\begin{equation}
\begin{aligned}
\frac{L_{\rm exact}}{m}
={}&
-c^2
+\left(\frac12 v^2+U\right)
+\frac{1}{c^2}\left[
\frac18 v^4
+\frac32 Uv^2
-\frac12 U^2
\right]
+\frac{1}{c^4}\left[
\frac1{16}v^6
+\frac78 Uv^4
+2U^2v^2
+\frac14 U^3
\right]
\\[0.75ex]
&
+\frac{1}{c^6}\left[
\frac5{128}v^8
+\frac{11}{16}Uv^6
+\frac{47}{16}U^2v^4
+\frac{13}{8}U^3v^2
-\frac18 U^4
\right]
+\frac{1}{c^8}\left[
\frac7{256}v^{10}
+\frac{75}{128}Uv^8
+\frac{59}{16}U^2v^6
+\frac{203}{32}U^3v^4
+\frac{31}{32}U^4v^2
+\frac1{16}U^5
\right]
+O(c^{-10}).
\end{aligned}
\label{eq:app-4pn-exact-schwarzschild-series}
\end{equation}
Equation~\eqref{eq:app-4pn-exact-schwarzschild-series} is the complete exact
one-body expansion through $4$PN.  In particular, the $4$PN coefficient is read
off without ambiguity as
\begin{equation}
\frac{L_{4\mathrm{PN},1\text{-body}}}{m}
=
\frac7{256}v^{10}
+\frac{75}{128}Uv^8
+\frac{59}{16}U^2v^6
+\frac{203}{32}U^3v^4
+\frac{31}{32}U^4v^2
+\frac1{16}U^5.
\label{eq:app-4pn-onebody-coefficient}
\end{equation}
This is the exact gate used in
Section~\ref{sec:onebody-schwarzschild} and nowhere in the later local chain is
there any freedom to alter it.

\subsection{Order-by-order one-body coefficients carried into the \texorpdfstring{$4$}{4}PN chain}
\label{app:onebody-expansion-orders}

If one drops the irrelevant additive rest-mass term $-mc^2$, then the strict
one-body coefficients imported into the conservative hierarchy are simply the
successive coefficients of
\eqref{eq:app-4pn-exact-schwarzschild-series}.  Writing
\[
L_{\rm exact}
=
-mc^2
+L_{\rm N,1\text{-body}}
+\frac{1}{c^2}L_{1\mathrm{PN},1\text{-body}}
+\frac{1}{c^4}L_{2\mathrm{PN},1\text{-body}}
+\frac{1}{c^6}L_{3\mathrm{PN},1\text{-body}}
+\frac{1}{c^8}L_{4\mathrm{PN},1\text{-body}}
+O(c^{-10}),
\]
one has
\begin{equation}
\frac{L_{\rm N,1\text{-body}}}{m}
=
\frac12 v^2+U,
\label{eq:app-4pn-newtonian-onebody}
\end{equation}
\begin{equation}
\frac{L_{1\mathrm{PN},1\text{-body}}}{m}
=
\frac18 v^4
+\frac32 Uv^2
-\frac12 U^2,
\label{eq:app-4pn-1pn-onebody}
\end{equation}
\begin{equation}
\frac{L_{2\mathrm{PN},1\text{-body}}}{m}
=
\frac1{16}v^6
+\frac78 Uv^4
+2U^2v^2
+\frac14 U^3,
\label{eq:app-4pn-2pn-onebody}
\end{equation}
\begin{equation}
\frac{L_{3\mathrm{PN},1\text{-body}}}{m}
=
\frac5{128}v^8
+\frac{11}{16}Uv^6
+\frac{47}{16}U^2v^4
+\frac{13}{8}U^3v^2
-\frac18 U^4,
\label{eq:app-4pn-3pn-onebody}
\end{equation}
and finally the $4$PN coefficient
\eqref{eq:app-4pn-onebody-coefficient}.  The first three of these are the
carried exact test-mass anchors from the earlier papers; the last one is the new
order-$c^{-8}$ input that the present paper adds.

The practical lesson is the one emphasized in the main text.  The exact
Schwarzschild series fixes the one-body front end before any denominator-style
continuation of the self sector is attempted.  The later repair ledger
\((\mu_{\rho,4},d_4,s_{34},s_{26})\) is therefore not a guess about the
Schwarzschild answer; it is the minimal continuation required to reproduce the
already frozen coefficient
\eqref{eq:app-4pn-onebody-coefficient} exactly.

% Appendix: One-body repair algebra

\section{One-body repair algebra}
\label{app:onebody-repair}

Appendix~\ref{app:onebody-expansion} freezes the exact isotropic
Schwarzschild one-body target through total order $c^{-8}$.  The natural next
question is whether the carried $3$PN denominator/self/static packaging can
simply be pushed one order higher, or whether the strict test-mass gate already
forces genuinely new $4$PN one-body data.  The purpose of the present appendix
is to answer that question algebraically and slot by slot.

The outcome is sharper than at $3$PN.  A quartic denominator datum is needed to
repair the $U^4v^2$ slot, a quartic static datum is needed to repair the pure
$U^5$ slot, and two additional explicit self-sector invariants are needed
because the $U^3v^4$ and $U^2v^6$ slots remain wrong even after the first two
repairs are imposed.  So the full one-body $4$PN repair ledger is already a
four-parameter ledger at the strict Schwarzschild gate.

\subsection{Quartic denominator datum}
\label{app:onebody-repair-denominator}

The first test is the minimal quartic continuation of the carried
 denominator-side self sector,
\begin{equation}
L_{\rm red}
=-mc^2(1-u)\sqrt{1-\frac{v^2/c^2}{D(u)}},
\qquad
u\equiv\frac{U}{c^2},
\qquad
D(u)=1-4u+8u^2+d_3u^3+d_4u^4+O(u^5).
\label{eq:app-4pn-onebody-repair-Lred}
\end{equation}
Expanding \eqref{eq:app-4pn-onebody-repair-Lred} through total order $c^{-8}$
gives
\begin{align}
\frac{L_{\rm red}}{m}
={}&-c^2+\frac12v^2+U
+\frac{1}{c^2}\left(\frac18v^4+\frac32Uv^2\right)
+\frac{1}{c^4}\left(\frac1{16}v^6+\frac78Uv^4+2U^2v^2\right)
\nonumber\\
&
+\frac{1}{c^6}\left(
\frac5{128}v^8
+\frac{11}{16}Uv^6
+3U^2v^4
+\left(-4-\frac{d_3}{2}\right)U^3v^2
\right)
\nonumber\\
&
+\frac{1}{c^8}\left(
\frac7{256}v^{10}
+\frac{75}{128}Uv^8
+\frac{15}{4}U^2v^6
+\left(4-\frac{d_3}{4}\right)U^3v^4
+\left(-32-\frac{7d_3}{2}-\frac{d_4}{2}\right)U^4v^2
\right)
+O(c^{-10}).
\label{eq:app-4pn-onebody-repair-Lred-series}
\end{align}
Two features of \eqref{eq:app-4pn-onebody-repair-Lred-series} are immediate.
First, the carried denominator logic already reproduces the universal $4$PN
coefficients of $v^{10}$ and $Uv^8$ automatically.  Second, once the carried
$3$PN denominator datum
\begin{equation}
\boxed{d_3=-\frac{45}{4}}
\label{eq:app-4pn-onebody-repair-d3-carried}
\end{equation}
is imposed, the quartic denominator coefficient $d_4$ can influence only the
$U^4v^2$ slot.  The $4$PN denominator-side coefficient then becomes
\begin{equation}
\left(\frac{L_{\rm red}}{m}\right)_{c^{-8},\,d_3=-45/4}
=
\frac7{256}v^{10}
+\frac{75}{128}Uv^8
+\frac{15}{4}U^2v^6
+\frac{109}{16}U^3v^4
+\left(\frac{59}{8}-\frac{d_4}{2}\right)U^4v^2.
\label{eq:app-4pn-onebody-repair-Lred-4pn}
\end{equation}
Matching the exact Schwarzschild target
\eqref{eq:app-4pn-onebody-coefficient} in the $U^4v^2$ channel gives
\begin{equation}
\frac{59}{8}-\frac{d_4}{2}=\frac{31}{32}
\qquad\Longrightarrow\qquad
\boxed{d_4=\frac{205}{16}.}
\label{eq:app-4pn-onebody-repair-d4}
\end{equation}
So the quartic denominator datum is fixed exactly by the $U^4v^2$ slot.

This also shows why the denominator-side continuation is not enough by itself.
Even before the static gate is discussed, the carried denominator sector still
predicts
\begin{equation}
\frac{15}{4}U^2v^6,
\qquad
\frac{109}{16}U^3v^4,
\label{eq:app-4pn-onebody-repair-denominator-mispredictions}
\end{equation}
whereas the exact Schwarzschild target requires
\begin{equation}
\frac{59}{16}U^2v^6,
\qquad
\frac{203}{32}U^3v^4.
\label{eq:app-4pn-onebody-repair-denominator-targets}
\end{equation}
So the quartic denominator correction repairs only the $U^4v^2$ gate.  It does
not repair the new $U^3v^4$ or $U^2v^6$ channels.

\subsection{Quartic static gate}
\label{app:onebody-repair-static}

The pure static slot gives an independent exact gate.  Extend the carried local
mass-scaling sector by one quartic static coefficient $\mu_{\rho,4}$,
\begin{equation}
\frac{L_{\rm static,1body}}{m}
=
U-\frac12\frac{U^2}{c^2}+\frac14\frac{U^3}{c^4}
-\frac{\mu_{\rho,3}}{2}\frac{U^4}{c^6}
+\frac{\mu_{\rho,4}}{2}\frac{U^5}{c^8}
+O(c^{-10}),
\label{eq:app-4pn-onebody-repair-static-sector}
\end{equation}
with the carried $3$PN value
\begin{equation}
\boxed{\mu_{\rho,3}=\frac14.}
\label{eq:app-4pn-onebody-repair-murho3-carried}
\end{equation}
The exact Schwarzschild target requires the pure static $4$PN term
\begin{equation}
\frac1{16}\frac{U^5}{c^8}.
\label{eq:app-4pn-onebody-repair-static-target}
\end{equation}
So the quartic static gate is fixed immediately by
\begin{equation}
\frac{\mu_{\rho,4}}{2}=\frac1{16}
\qquad\Longrightarrow\qquad
\boxed{\mu_{\rho,4}=\frac18.}
\label{eq:app-4pn-onebody-repair-murho4}
\end{equation}
As at lower order, the important point is structural rather than cosmetic: the
pure static slot remains a separate exact gate of its own.  It is not absorbed
into the denominator datum and it is not a bookkeeping convenience.

\subsection{Two new self-sector slots and why a quartic denominator plus a quartic static gate is not enough}
\label{app:onebody-repair-self}

At this point the minimal direct-continuation one-body candidate is
\begin{equation}
\begin{aligned}
\frac{L_{\rm cand}}{m}
={}&
\frac{L_{\rm red}}{m}
-\frac12\frac{U^2}{c^2}
+\frac14\frac{U^3}{c^4}
-\frac{\mu_{\rho,3}}{2}\frac{U^4}{c^6}
+\frac{\mu_{\rho,4}}{2}\frac{U^5}{c^8}
\\[0.5ex]
&
+s_{24}\frac{U^2v^4}{c^6}
+s_{34}\frac{U^3v^4}{c^8}
+s_{26}\frac{U^2v^6}{c^8},
\end{aligned}
\label{eq:app-4pn-onebody-repair-candidate}
\end{equation}
where the first five displayed corrections restore the carried lower-order
static ledger together with the new quartic static gate, and the last two terms
are the minimal new explicit $4$PN self-sector slots.  Subtracting
\eqref{eq:app-4pn-onebody-repair-candidate} from the exact target and imposing
the already frozen lower-order values
\begin{equation}
\boxed{
 d_3=-\frac{45}{4},
 \qquad
 \mu_{\rho,3}=\frac14,
 \qquad
 s_{24}=-\frac1{16}
 }
\label{eq:app-4pn-onebody-repair-carried-lower-order-ledger}
\end{equation}
gives the exact residual
\begin{equation}
\left(\frac{L_{4\mathrm{PN},1\text{-body}}}{m}-\frac{L_{\rm cand}}{m}\right)_{c^{-8}}
=
\left(\frac1{16}-\frac{\mu_{\rho,4}}{2}\right)U^5
+\left(-\frac{205}{32}+\frac{d_4}{2}\right)U^4v^2
+\left(-\frac{15}{32}-s_{34}\right)U^3v^4
+\left(-\frac1{16}-s_{26}\right)U^2v^6.
\label{eq:app-4pn-onebody-repair-residual}
\end{equation}
So the one-body repair algebra again separates cleanly by slot.

The first two coefficients in
\eqref{eq:app-4pn-onebody-repair-residual} vanish once
\eqref{eq:app-4pn-onebody-repair-d4} and
\eqref{eq:app-4pn-onebody-repair-murho4} are imposed.  But even after those two
exact repairs, the candidate with no new $4$PN self slots still leaves the
uncancelled discrepancy
\begin{equation}
\left(\frac{L_{4\mathrm{PN},1\text{-body}}}{m}-\frac{L_{\rm cand}}{m}\right)_{c^{-8},\,d_4=205/16,\,\mu_{\rho,4}=1/8,\,s_{34}=s_{26}=0}
=
-\frac{15}{32}U^3v^4-\frac1{16}U^2v^6.
\label{eq:app-4pn-onebody-repair-two-slot-mismatch}
\end{equation}
So a quartic denominator plus a quartic static gate is still not enough.  Two
new self-sector slots survive exactly at $4$PN.  They are repaired by the
conditions
\begin{equation}
-\frac{15}{32}-s_{34}=0
\qquad\Longrightarrow\qquad
\boxed{s_{34}=-\frac{15}{32},}
\label{eq:app-4pn-onebody-repair-s34}
\end{equation}

after which the remaining $U^2v^6$ discrepancy is fixed by
\begin{equation}
-\frac1{16}-s_{26}=0
\qquad\Longrightarrow\qquad
\boxed{s_{26}=-\frac1{16}.}
\label{eq:app-4pn-onebody-repair-s26}
\end{equation}
Collecting everything, the exact $4$PN one-body repair ledger is
\begin{equation}
\boxed{
\mu_{\rho,4}=\frac18,
\qquad
 d_4=\frac{205}{16},
\qquad
 s_{34}=-\frac{15}{32},
\qquad
 s_{26}=-\frac1{16}.
}
\label{eq:app-4pn-onebody-repair-ledger}
\end{equation}
This is the cleanest strict one-body lesson of the whole $4$PN analysis.  The
$4$PN gate does \emph{not} say ``continue the $3$PN denominator idea by one new
coefficient.''  It says that four distinct data are already needed at the
strict test-mass level:
\begin{enumerate}
    \item a quartic static gate,
    \item a quartic denominator datum,
    \item and two genuinely new self-sector slots.
\end{enumerate}
That is exactly why the full $4$PN program had to begin with a one-body audit
before the comparable-mass local lift was even well posed.  Once the exact
target and its minimal repair ledger are frozen, the next step is to promote the
strict one-body structure into the natural two-body self/static seed used in the
local theorem chain.

% Appendix: Exact quartic Legendre-transform identity at 4PN

\section{Exact quartic Legendre-transform identity at \texorpdfstring{$4$}{4}PN}
\label{app:quartic-compiler}

Section~\ref{sec:onebody-quartic-compiler} recorded the compact quartic
ordinary-Lagrangian $\to$ Hamiltonian compiler used throughout the local $4$PN
construction.  The purpose of the present appendix is to freeze that identity in
its full perturbative form, in the same spirit in which the earlier $3$PN paper
froze the exact cubic compiler.  At $4$PN the first genuinely new compiler
content is the \emph{quartic-order} extension.  Once that identity is
available, the later reduced COM translations and generic-frame
Hamiltonian-first lifts are controlled algebraically rather than guessed.

The role of this appendix is therefore very specific.  It does not solve the
full comparable-mass problem by itself.  What it does is fix the exact ordinary
Lagrangian $\to$ Hamiltonian compiler through order $\varepsilon^4$ with
\(\varepsilon=c^{-2}\).  Once that compiler is frozen, the remaining local
$4$PN comparison against the fixed-chart Hamiltonian target becomes a finite
algebraic residual problem.

\subsection{Perturbative setup}
\label{app:quartic-legendre-setup}

Write the conservative ordinary Lagrangian as
\begin{equation}
L(v)=L_0(v)+\varepsilon L_1(v)+\varepsilon^2L_2(v)+\varepsilon^3L_3(v)+\varepsilon^4L_4(v),
\qquad
\varepsilon=c^{-2},
\label{eq:app-4pn-quartic-legendre-L-expansion}
\end{equation}
with quadratic zeroth-order block
\begin{equation}
L_0(v)=\frac12 v^T M v,
\label{eq:app-4pn-quartic-legendre-L0}
\end{equation}
where the Newtonian mass matrix \(M\) is constant at the order relevant for the
perturbative inversion.  Throughout this appendix, gradients with respect to
\(v\) are taken as column vectors and Hessians are taken as matrices.

The canonical momentum is
\begin{equation}
p\equiv \frac{\partial L}{\partial v}
= Mv
+\varepsilon\frac{\partial L_1}{\partial v}
+\varepsilon^2\frac{\partial L_2}{\partial v}
+\varepsilon^3\frac{\partial L_3}{\partial v}
+\varepsilon^4\frac{\partial L_4}{\partial v}.
\label{eq:app-4pn-quartic-legendre-momentum}
\end{equation}
Define the Newtonian momentum-space velocity
\begin{equation}
v_0=M^{-1}p,
\label{eq:app-4pn-quartic-legendre-v0}
\end{equation}
and the derivative data evaluated at \(v_0\),
\begin{equation}
A_0\equiv \left.\frac{\partial L_1}{\partial v}\right|_{v_0},
\qquad
B_0\equiv \left.\frac{\partial L_2}{\partial v}\right|_{v_0},
\qquad
D_0\equiv \left.\frac{\partial L_3}{\partial v}\right|_{v_0},
\label{eq:app-4pn-quartic-legendre-ABD}
\end{equation}
\begin{equation}
C_0\equiv \left.\frac{\partial^2L_1}{\partial v^2}\right|_{v_0},
\qquad
E_0\equiv \left.\frac{\partial^2L_2}{\partial v^2}\right|_{v_0},
\qquad
T_0\equiv \left.\frac{\partial^3L_1}{\partial v^3}\right|_{v_0}.
\label{eq:app-4pn-quartic-legendre-CET}
\end{equation}
Here \(T_0\) is the symmetric trilinear form defined by the third derivatives of
\(L_1\).  For vectors \(x,y,z\), the notation \(T_0[x,y,z]\) denotes full
contraction, while \(T_0[x,y,\cdot]\) denotes the vector obtained by leaving
one slot uncontracted.

The perturbative momentum inversion is written as
\begin{equation}
v=v_0+\varepsilon v_1+\varepsilon^2v_2+\varepsilon^3v_3+O(\varepsilon^4).
\label{eq:app-4pn-quartic-legendre-v-series}
\end{equation}
For the quartic Hamiltonian coefficient \(H_4\), only \(v_1\), \(v_2\), and
\(v_3\) are needed explicitly.

\subsection{Order-by-order momentum inversion}
\label{app:quartic-legendre-inversion}

Insert \eqref{eq:app-4pn-quartic-legendre-v-series} into
\eqref{eq:app-4pn-quartic-legendre-momentum} and expand around \(v_0\).  To the
order needed for \(H_4\), the Taylor expansions of the derivative blocks are
\begin{equation}
\frac{\partial L_1}{\partial v}(v)
=
A_0
+\varepsilon C_0v_1
+\varepsilon^2\left(C_0v_2+\frac12T_0[v_1,v_1,\cdot]\right)
+O(\varepsilon^3),
\label{eq:app-4pn-quartic-legendre-dL1}
\end{equation}
\begin{equation}
\frac{\partial L_2}{\partial v}(v)
=
B_0+\varepsilon E_0v_1+O(\varepsilon^2),
\qquad
\frac{\partial L_3}{\partial v}(v)
=
D_0+O(\varepsilon).
\label{eq:app-4pn-quartic-legendre-dL23}
\end{equation}
Using \(Mv_0=p\), the momentum equation becomes
\begin{equation}
\begin{aligned}
0
={}&
\varepsilon\bigl(Mv_1+A_0\bigr)
+\varepsilon^2\bigl(Mv_2+B_0+C_0v_1\bigr)
\\[0.5ex]
&
+\varepsilon^3\bigl(Mv_3+D_0+E_0v_1+C_0v_2+\tfrac12T_0[v_1,v_1,\cdot]\bigr)
+O(\varepsilon^4).
\end{aligned}
\label{eq:app-4pn-quartic-legendre-momentum-orders}
\end{equation}
So the coefficients are fixed order by order:
\begin{align}
Mv_1+A_0 &= 0,
\label{eq:app-4pn-quartic-legendre-order-eqs-a}\\
Mv_2+B_0+C_0v_1 &= 0,
\label{eq:app-4pn-quartic-legendre-order-eqs-b}\\
Mv_3+D_0+E_0v_1+C_0v_2+\frac12T_0[v_1,v_1,\cdot] &= 0.
\label{eq:app-4pn-quartic-legendre-order-eqs-c}
\end{align}
Equations~\eqref{eq:app-4pn-quartic-legendre-order-eqs-a}--\eqref{eq:app-4pn-quartic-legendre-order-eqs-c}
fix the inversion coefficients order by order.
Therefore
\begin{equation}
\boxed{
\begin{aligned}
v_1={}&-M^{-1}A_0,
\\[0.5ex]
v_2={}&M^{-1}C_0M^{-1}A_0-M^{-1}B_0,
\\[0.5ex]
v_3={}&-M^{-1}D_0
+M^{-1}E_0M^{-1}A_0
+M^{-1}C_0M^{-1}B_0
\\[0.5ex]
&
-M^{-1}C_0M^{-1}C_0M^{-1}A_0
-\frac12M^{-1}T_0[M^{-1}A_0,M^{-1}A_0,\cdot].
\end{aligned}
}
\label{eq:app-4pn-quartic-legendre-v123}
\end{equation}
Equation~\eqref{eq:app-4pn-quartic-legendre-v123} is the exact perturbative
inversion needed through quartic order.

\subsection{Exact quartic Hamiltonian identity}
\label{app:quartic-legendre-identity}

Now substitute \eqref{eq:app-4pn-quartic-legendre-v-series} and
\eqref{eq:app-4pn-quartic-legendre-v123} into the Legendre definition
\begin{equation}
H(p)=p\cdot v-L(v).
\label{eq:app-4pn-quartic-legendre-H-def}
\end{equation}
The Hamiltonian expansion takes the form
\begin{equation}
H=H_0+\varepsilon H_1+\varepsilon^2H_2+\varepsilon^3H_3+\varepsilon^4H_4+O(\varepsilon^5),
\label{eq:app-4pn-quartic-legendre-H-series}
\end{equation}
with zeroth-order piece
\begin{equation}
H_0=\frac12 p^TM^{-1}p.
\label{eq:app-4pn-quartic-legendre-H0}
\end{equation}
Carrying the quartic expansion through exactly gives
\begin{equation}
H_1=-L_1(v_0),
\label{eq:app-4pn-quartic-legendre-H1}
\end{equation}
\begin{equation}
H_2=-L_2(v_0)+\frac12A_0^TM^{-1}A_0,
\label{eq:app-4pn-quartic-legendre-H2}
\end{equation}
\begin{equation}
H_3=-L_3(v_0)+A_0^TM^{-1}B_0-\frac12A_0^TM^{-1}C_0M^{-1}A_0,
\label{eq:app-4pn-quartic-legendre-H3}
\end{equation}
and the exact $4$PN compiler identity
\begin{equation}
\boxed{
\begin{aligned}
H_4={}&-L_4(v_0)
+A_0^TM^{-1}D_0
+\frac12B_0^TM^{-1}B_0
-B_0^TM^{-1}C_0M^{-1}A_0
\\[0.5ex]
&
-\frac12A_0^TM^{-1}E_0M^{-1}A_0
+\frac12A_0^TM^{-1}C_0M^{-1}C_0M^{-1}A_0
+\frac16T_0[M^{-1}A_0,M^{-1}A_0,M^{-1}A_0].
\end{aligned}
}
\label{eq:app-4pn-quartic-legendre-H4}
\end{equation}
Equation~\eqref{eq:app-4pn-quartic-legendre-H4} is the exact quartic analogue
of the quadratic and cubic perturbative compiler identities used at lower
order.  It is the formula actually used to carry the ordinary $4$PN data into
the Hamiltonian chart and back again.

It is also useful to record the reduced COM specialization.  In the normalized
COM variables used later, the Newtonian quadratic block is the identity, so
\(M=I\) and therefore
\begin{equation}
v_0=p,
\qquad
\begin{aligned}
H_4={}&-L_4(p)
+A_0\cdot D_0
+\frac12B_0\cdot B_0
-B_0^TC_0A_0
\\[0.5ex]
&
-\frac12A_0^TE_0A_0
+\frac12A_0^TC_0C_0A_0
+\frac16T_0[A_0,A_0,A_0].
\end{aligned}
\label{eq:app-4pn-quartic-legendre-H4-com}
\end{equation}
This is the compact form checked by the local referee scripts before the later
slot-by-slot COM compiler and fixed-chart local target comparison are built.

\subsection{Practical meaning for the \texorpdfstring{$4$}{4}PN lift}
\label{app:quartic-legendre-meaning}

The practical meaning of \eqref{eq:app-4pn-quartic-legendre-H4} is exactly the
one that makes the local $4$PN chain manageable.  Once the lower-order
$1$PN/$2$PN/$3$PN ledger is frozen, the feedback tensors \(A_0\), \(B_0\),
\(C_0\), \(D_0\), \(E_0\), and \(T_0\) are already known.  So the genuinely
new $4$PN Hamiltonian content enters only through the explicit
\(-L_4(v_0)\) term, while the remaining quartic feedback terms are completely
fixed by carried lower-order data.

That observation is what turns the local comparable-mass $4$PN problem into a
finite algebraic residual rather than a blind full derivation campaign.  It is
also why the compiler has to be frozen before the local fixed-chart target is
translated: without \eqref{eq:app-4pn-quartic-legendre-H4}, the ordinary-chart
and Hamiltonian-chart comparisons are not rigidly related.  With it, the
Hamiltonian-first local lift becomes a theorem-level statement rather than an
ansatz.

Finally, the same compiler applied to the exact one-body Schwarzschild
Lagrangian gate reproduces the strict one-body $4$PN Hamiltonian seed used in
the main local theorem chain.  So the one-body Hamiltonian gate is not an
independent imported object.  It is a direct consequence of the exact
Schwarzschild Lagrangian gate together with the quartic compiler frozen here.

% Appendix: Local generic-frame scaffold and COM projection

\section{Local generic-frame scaffold and COM projection}
\label{app:local-generic-frame-scaffold-com}

Section~\ref{sec:local-chain} used the raw local scaffold only through its
headline counts.  The purpose of the present appendix is to freeze that scaffold
itself, together with the exact center-of-mass (COM) projection map that first
reveals its structure.  At this stage no fixed-chart target has yet been
imported and no Hamiltonian-first quotient has yet been taken.  The question is
therefore purely kinematic: once the strict one-body content has been absorbed
into the natural self/static seed \eqref{eq:4pn-natural-self-static-seed}, how
large is the remaining exchange-symmetric generic-frame local interaction space,
and what does COM projection determine about it exactly?

The answer is already very sharp.  The raw local comparable-mass scaffold is
blockwise COM-complete from the start.  There is no early span shortage.  What
COM projection leaves behind is not a missing image direction but a very large
COM-blind null family.  This is exactly why the later local theorem chain must
be completed by importing the fixed-chart Hamiltonian target and quotienting the
null sector rather than by enlarging the raw scaffold itself.

\subsection{Exchange-symmetric local interaction scaffold}
\label{app:local-generic-frame-scaffold-basis}

Use the standard generic-frame invariants
\begin{equation}
 a=v_A^2,
 \qquad
 b=v_B^2,
 \qquad
 c=\mathbf v_A\!\cdot\!\mathbf v_B,
 \qquad
 d=\mathbf v_A\!\cdot\!\mathbf n,
 \qquad
 e=\mathbf v_B\!\cdot\!\mathbf n,
 \label{eq:app-4pn-local-scaffold-abcde}
\end{equation}
with mass variables
\begin{equation}
 p=m_A,
 \qquad
 q=m_B.
 \label{eq:app-4pn-local-scaffold-pq}
\end{equation}
The raw local comparable-mass scaffold is defined relative to the natural
self/static seed by
\begin{equation}
L_{4,\mathrm{loc}}^{\mathrm{raw}}
=
L_{4,\mathrm{seed}}
+
L_{4,\mathrm{res}}^{\mathrm{raw}},
\label{eq:app-4pn-local-scaffold-split}
\end{equation}
where
\begin{equation}
\begin{aligned}
L_{4,\mathrm{res}}^{\mathrm{raw}}
={}&
\frac{Gm_A m_B}{r}\sum_{i=1}^{52}\alpha_i Q_i
+\frac{G^2m_A m_B}{r^2}\sum_{j=1}^{46}\beta_j T_j
+\frac{G^3m_A m_B}{r^3}\sum_{k=1}^{29}\gamma_k S_k
\\[0.5ex]
&
+\frac{G^4m_A m_B}{r^4}\sum_{\ell=1}^{10}\delta_{\ell} U_{\ell}
+\frac{G^5m_A m_B}{r^5}\sum_{m=1}^{2}\epsilon_m W_m.
\end{aligned}
\label{eq:app-4pn-local-scaffold-residual}
\end{equation}
The basis families are generated algorithmically as exchange-symmetric monomials
in \eqref{eq:app-4pn-local-scaffold-abcde}--\eqref{eq:app-4pn-local-scaffold-pq}
with blockwise velocity degrees
\begin{equation}
\deg_v(Q_i,T_j,S_k,U_{\ell},W_m)=(8,6,4,2,0)
\label{eq:app-4pn-local-scaffold-velocity-degrees}
\end{equation}
and additional mass degrees
\begin{equation}
\deg_m(Q_i,T_j,S_k,U_{\ell},W_m)=(0,1,2,3,4).
\label{eq:app-4pn-local-scaffold-mass-degrees}
\end{equation}
Throughout this appendix, ``raw'' means that no contact/gauge quotient has yet
been imposed.

The strict one-body gate is kept clean by requiring each residual basis element
to vanish on the test-mass branch in which body $A$ becomes the test particle
and body $B$ the fixed source.  In the invariant variables this means
\begin{equation}
 b=0,
 \qquad
 c=0,
 \qquad
 e=0,
 \qquad
 p=0,
 \qquad
 q=1.
 \label{eq:app-4pn-local-scaffold-testmass-branch}
\end{equation}
So every monomial retained in
\eqref{eq:app-4pn-local-scaffold-residual} vanishes on the strict one-body
branch, and the one-body content stays entirely inside
\(L_{4,\mathrm{seed}}\).

The exact block dimensions obtained by this construction are
\begin{equation}
\dim\mathcal Q=52,
\qquad
\dim\mathcal T=46,
\qquad
\dim\mathcal S=29,
\qquad
\dim\mathcal U=10,
\qquad
\dim\mathcal W=2,
\label{eq:app-4pn-local-scaffold-block-dims}
\end{equation}
so the total raw local interaction scaffold contains
\begin{equation}
52+46+29+10+2=139
\label{eq:app-4pn-local-scaffold-total-dim}
\end{equation}
exchange-symmetric directions before any fixed-chart conditions are imposed.

\subsection{Exact COM reduction map and interaction-slot extraction}
\label{app:local-generic-frame-scaffold-com-map}

To read off the COM image, write the standard reduced COM variables as
\begin{equation}
\mathbf v_A=X_B\mathbf v,
\qquad
\mathbf v_B=-X_A\mathbf v,
\qquad
X_A+X_B=1,
\qquad
\nu=X_AX_B,
\label{eq:app-4pn-local-scaffold-com-kinematics}
\end{equation}
with the usual mass-fraction parameterization
\begin{equation}
X_A=\frac{1+\Delta}{2},
\qquad
X_B=\frac{1-\Delta}{2},
\qquad
\Delta^2=1-4\nu.
\label{eq:app-4pn-local-scaffold-XA-XB}
\end{equation}
Therefore the generic-frame invariants reduce to
\begin{equation}
 a=X_B^2 v^2,
 \qquad
 b=X_A^2 v^2,
 \qquad
 c=-X_A X_B v^2=-\nu v^2,
 \label{eq:app-4pn-local-scaffold-com-abc}
\end{equation}
\begin{equation}
 d=X_B\dot r,
 \qquad
 e=-X_A\dot r.
 \label{eq:app-4pn-local-scaffold-com-de}
\end{equation}
Because the raw basis is exchange-symmetric, the COM image is obtained by first
substituting \eqref{eq:app-4pn-local-scaffold-com-abc}--\eqref{eq:app-4pn-local-scaffold-com-de},
then symmetrizing under \(\Delta\mapsto-\Delta\), and finally eliminating all
remaining even powers of \(\Delta\) using \eqref{eq:app-4pn-local-scaffold-XA-XB}.
Every basis element then becomes a polynomial in \(\nu\) multiplying one of the
standard reduced COM monomials in \(v^2\) and \(\dot r^2\).

The induced COM interaction slots are exactly the fifteen slots already used in
Section~\ref{sec:local-raw-scaffold}.  Define the blockwise coefficient
extraction maps by
\begin{equation}
\Pi_Q(F)
\equiv
\bigl(
[v^8]F_{\rm COM},
[v^6\dot r^2]F_{\rm COM},
[v^4\dot r^4]F_{\rm COM},
[v^2\dot r^6]F_{\rm COM},
[\dot r^8]F_{\rm COM}
\bigr)^T,
\label{eq:app-4pn-local-scaffold-PiQ}
\end{equation}
\begin{equation}
\Pi_T(F)
\equiv
\bigl(
[v^6]F_{\rm COM},
[v^4\dot r^2]F_{\rm COM},
[v^2\dot r^4]F_{\rm COM},
[\dot r^6]F_{\rm COM}
\bigr)^T,
\label{eq:app-4pn-local-scaffold-PiT}
\end{equation}
\begin{equation}
\Pi_S(F)
\equiv
\bigl(
[v^4]F_{\rm COM},
[v^2\dot r^2]F_{\rm COM},
[\dot r^4]F_{\rm COM}
\bigr)^T,
\label{eq:app-4pn-local-scaffold-PiS}
\end{equation}
\begin{equation}
\Pi_U(F)
\equiv
\bigl(
[v^2]F_{\rm COM},
[\dot r^2]F_{\rm COM}
\bigr)^T,
\qquad
\Pi_W(F)
\equiv
\bigl([1]F_{\rm COM}\bigr).
\label{eq:app-4pn-local-scaffold-PiUW}
\end{equation}
So, in a chosen raw basis ordering, the exact COM projection matrices are
\begin{equation}
M_Q\in\mathrm{Mat}_{5\times 52}(\mathbb Q[\nu]),
\qquad
M_T\in\mathrm{Mat}_{4\times 46}(\mathbb Q[\nu]),
\qquad
M_S\in\mathrm{Mat}_{3\times 29}(\mathbb Q[\nu]),
\label{eq:app-4pn-local-scaffold-proj-mats-a}
\end{equation}
\begin{equation}
M_U\in\mathrm{Mat}_{2\times 10}(\mathbb Q[\nu]),
\qquad
M_W\in\mathrm{Mat}_{1\times 2}(\mathbb Q[\nu]).
\label{eq:app-4pn-local-scaffold-proj-mats-b}
\end{equation}
The audit establishes the exact blockwise ranks
\begin{equation}
\operatorname{rank}M_Q=5,
\qquad
\operatorname{rank}M_T=4,
\qquad
\operatorname{rank}M_S=3,
\qquad
\operatorname{rank}M_U=2,
\qquad
\operatorname{rank}M_W=1,
\label{eq:app-4pn-local-scaffold-ranks}
\end{equation}
with nullities
\begin{equation}
\dim\ker M_Q=47,
\qquad
\dim\ker M_T=42,
\qquad
\dim\ker M_S=26,
\qquad
\dim\ker M_U=8,
\qquad
\dim\ker M_W=1.
\label{eq:app-4pn-local-scaffold-nullities}
\end{equation}
Equivalently, the exact COM interaction-slot structure is
\begin{equation}
\#_{\rm COM}^{\rm int}=(5,4,3,2,1),
\qquad
5+4+3+2+1=15,
\label{eq:app-4pn-local-scaffold-slot-counts}
\end{equation}
so the total COM interaction rank is
\begin{equation}
5+4+3+2+1=15,
\label{eq:app-4pn-local-scaffold-total-rank}
\end{equation}
and the total COM-blind nullity is
\begin{equation}
47+42+26+8+1=124.
\label{eq:app-4pn-local-scaffold-total-nullity}
\end{equation}
The interaction-slot count in
\eqref{eq:app-4pn-local-scaffold-slot-counts} is the interaction analogue of
the full $21$-slot local $4$PN reduced ordinary basis, whose remaining six
slots are pure kinetic.

So the first exact conclusion is already visible before any target import:
\begin{equation}
\boxed{
\text{the raw local generic-frame scaffold is blockwise COM-complete.}
}
\label{eq:app-4pn-local-scaffold-complete}
\end{equation}
At this stage the only large ambiguity is the $124$-dimensional COM-blind null
family.

\subsection{A minimal pivot-spanning subset}
\label{app:local-generic-frame-scaffold-pivots}

A convenient blockwise pivot-spanning subset is the following.  For the
\(G/r\) block one may take
\begin{equation}
\begin{aligned}
\mathcal P_Q=
\bigl\{
& a^2b^2,
\quad
a^2bd^2+ab^2e^2,
\quad
a^2d^2e^2+b^2d^2e^2,
\\[0.5ex]
& ad^2e^4+bd^4e^2,
\quad
d^4e^4
\bigr\}.
\end{aligned}
\label{eq:app-4pn-local-scaffold-pivots-Q}
\end{equation}
Their COM images are
\begin{equation}
\begin{aligned}
\Pi_Q(\mathcal P_Q)=
\bigl\{
& \nu^4 v^8,
\quad
(\nu^2-4\nu^3+2\nu^4) v^6\dot r^2,
\quad
(\nu^2-4\nu^3+2\nu^4) v^4\dot r^4,
\\[0.5ex]
& 2\nu^4 v^2\dot r^6,
\quad
\nu^4 \dot r^8
\bigr\},
\end{aligned}
\label{eq:app-4pn-local-scaffold-pivot-images-Q}
\end{equation}
which already span the full five-dimensional COM image of the \(G/r\) block.

For the \(G^2/r^2\) block one may take
\begin{equation}
\begin{aligned}
\mathcal P_T=
\bigl\{
& a^2bp+ab^2q,
\quad
a^2d^2p+b^2e^2q,
\quad
ad^2e^2p+bd^2e^2q,
\quad
d^3e^3p+d^3e^3q
\bigr\},
\end{aligned}
\label{eq:app-4pn-local-scaffold-pivots-T}
\end{equation}
with COM images
\begin{equation}
\begin{aligned}
\Pi_T(\mathcal P_T)=
\bigl\{
& \nu^3 v^6,
\quad
(\nu-5\nu^2+5\nu^3)v^4\dot r^2,
\quad
\nu^3 v^2\dot r^4,
\quad
-\nu^3\dot r^6
\bigr\}.
\end{aligned}
\label{eq:app-4pn-local-scaffold-pivot-images-T}
\end{equation}

For the \(G^3/r^3\) block one may take
\begin{equation}
\mathcal P_S=
\bigl\{
 a^2p^2+b^2q^2,
 \qquad
 ad^2p^2+be^2q^2,
 \qquad
 d^2e^2p^2+d^2e^2q^2
\bigr\},
\label{eq:app-4pn-local-scaffold-pivots-S}
\end{equation}
with COM images
\begin{equation}
\Pi_S(\mathcal P_S)=
\bigl\{
(\nu^2-2\nu^3)v^4,
\qquad
(\nu^2-2\nu^3)v^2\dot r^2,
\qquad
(\nu^2-2\nu^3)\dot r^4
\bigr\}.
\label{eq:app-4pn-local-scaffold-pivot-images-S}
\end{equation}

For the \(G^4/r^4\) block one may take
\begin{equation}
\mathcal P_U=
\bigl\{
 ap^2q+bpq^2,
 \qquad
 d^2p^2q+e^2pq^2
\bigr\},
\label{eq:app-4pn-local-scaffold-pivots-U}
\end{equation}
with COM images
\begin{equation}
\Pi_U(\mathcal P_U)=
\bigl\{
\nu^2 v^2,
\qquad
\nu^2\dot r^2
\bigr\}.
\label{eq:app-4pn-local-scaffold-pivot-images-U}
\end{equation}
Finally, for the static \(G^5/r^5\) block one may take
\begin{equation}
\mathcal P_W=\{p^2q^2\},
\qquad
\Pi_W(\mathcal P_W)=\{\nu^2\}.
\label{eq:app-4pn-local-scaffold-pivots-W}
\end{equation}
These pivot families are not unique, but they are especially convenient because
they already display the full five-block COM image with the smallest possible
number of representatives:
\begin{equation}
5+4+3+2+1=15.
\label{eq:app-4pn-local-scaffold-pivot-total}
\end{equation}

\subsection{Structural consequence}
\label{app:local-generic-frame-scaffold-consequence}

The structural interpretation is immediate.  Already before any fixed-chart
local target is imported, the raw constant-coefficient exchange-symmetric local
$4$PN scaffold is large enough to span the full COM interaction image block by
block.  So the first local $4$PN bottleneck is not undercompleteness of the raw
generic-frame basis.  It is the opposite: a very large COM-blind null family.
That is exactly why the next mathematically honest step is the one taken in
Section~\ref{sec:local-target-import}: import the fixed local Hamiltonian
chart, read off the true blockwise degree ceilings, and quotient the
$124$-dimensional COM-null sector by those fixed-chart conditions rather than by
COM data alone.

The local scaffold audit therefore establishes the first local $4$PN theorem
front end:
\begin{equation}
\boxed{
\text{there is no early kinematic shortage at local }4\text{PN.}
}
\label{eq:app-4pn-local-scaffold-no-shortage}
\end{equation}
Everything that remains in the later local theorem chain is a chart, quotient,
and seed-alignment problem, not a failure of the raw exchange-symmetric
scaffold.

% Appendix: Imported fixed local 4PN target and the Hamiltonian-first translation lesson

\section{Imported fixed local \texorpdfstring{$4$}{4}PN target and the Hamiltonian-first translation lesson}
\label{app:imported-local-target-hamiltonian-first}

Section~\ref{sec:local-target-import} used only the structural conclusions of the
fixed-chart local target import.  The purpose of the present appendix is to
freeze that import itself, together with the exact quartic translation lesson
that determines the correct order of the local generic-frame lift.  The result
is conceptually simple but technically decisive.  The local instantaneous target
is most naturally imported in the reduced center-of-mass (COM) \emph{Hamiltonian}
chart, where its blockwise mass-ratio degree ceilings match the constant-
coefficient generic-frame scaffold exactly.  After quartic translation to the
ordinary chart, by contrast, the middle and upper local blocks acquire extra
\(\nu^4\) tails that are not present in the Hamiltonian target and are not
reachable from the raw constant-coefficient ordinary scaffold.  This is the
precise reason the local \(4\)PN lift has to be done Hamiltonian-first.

\subsection{Imported fixed local Hamiltonian target}
\label{app:imported-local-target-hamiltonian-chart}

The first point to freeze is the target logic itself.  At conservative \(4\)PN,
the exact source split is
\begin{equation}
L_{4\mathrm{PN}}^{\mathrm{cons}}
=
L_{4\mathrm{PN}}^{\mathrm{local}}
+
L_{4\mathrm{PN}}^{\mathrm{tail}}.
\label{eq:app-local-target-split}
\end{equation}
So the present appendix concerns only the \emph{local instantaneous} target.
Concretely, the reference local Hamiltonian is the fixed-chart ADM local
\(4\)PN Hamiltonian, imported consistently with the already-frozen
local/nonlocal split.  The harmonic/Fokker and EFT constructions play only the
role of cross-check families; they are not used to redefine the target.

In reduced COM variables, the imported local Hamiltonian organizes into one pure
free-kinetic slot and fifteen interaction slots,
\begin{equation}
1\;\oplus\;5\;\oplus\;4\;\oplus\;3\;\oplus\;2\;\oplus\;1,
\label{eq:app-local-target-slot-structure}
\end{equation}
corresponding respectively to the free block and the interaction blocks
\(G/r\), \(G^2/r^2\), \(G^3/r^3\), \(G^4/r^4\), and \(G^5/r^5\).
This matches exactly the slot structure already seen in the raw local scaffold
appendix.

The strict one-body check is equally exact.  Applying the quartic Legendre
compiler of Appendix~\ref{app:quartic-legendre-identity} to the frozen
one-body Schwarzschild Lagrangian gives the reduced one-body \(4\)PN
Hamiltonian gate
\begin{equation}
\frac{7}{256}p^{10},
\qquad
\frac{45}{128}\frac{p^8}{r},
\qquad
\frac{13}{8}\frac{p^6}{r^2},
\qquad
\frac{105}{32}\frac{p^4}{r^3},
\qquad
\frac{105}{32}\frac{p^2}{r^4},
\qquad
-\frac{1}{16}\frac{1}{r^5}.
\label{eq:app-local-target-onebody-hamiltonian-gate}
\end{equation}
The imported fixed-chart local Hamiltonian matches
\eqref{eq:app-local-target-onebody-hamiltonian-gate} exactly in the strict
test-mass limit \(\nu\to0\).

The first structural theorem gate answered by the import is the blockwise
\(\nu\)-degree ceiling of the local Hamiltonian target.  The outcome is
\begin{equation}
\deg_{\nu}^{\rm max}(\text{free},G/r,G^2/r^2,G^3/r^3,G^4/r^4,G^5/r^5)
=
(4,4,3,3,2,2).
\label{eq:app-local-target-hamiltonian-degree-ceilings}
\end{equation}
In particular,
\begin{equation}
\deg_{\nu}^{\rm max}(G^4/r^4)=2,
\qquad
\deg_{\nu}^{\rm max}(G^5/r^5)=2,
\label{eq:app-local-target-upper-block-cutoff}
\end{equation}
so the imported fixed-chart local target does \emph{not} force any \(\nu^3\) or
\(\nu^4\) tails in the upper local interaction blocks.

This already removes one possible obstruction.  At the level of the imported
local Hamiltonian target, the upper interaction blocks are no more singular than
the raw constant-coefficient scaffold had suggested.

\subsection{Exact quartic COM translation and the ordinary local target}
\label{app:imported-local-target-ordinary-translation}

To make the chart lesson exact rather than heuristic, one translates the
imported local Hamiltonian target through the full reduced COM quartic compiler.
Use the full twenty-one-slot reduced COM ordinary basis,
\begin{equation}
6\;\oplus\;5\;\oplus\;4\;\oplus\;3\;\oplus\;2\;\oplus\;1,
\label{eq:app-local-target-ordinary-basis-structure}
\end{equation}
corresponding to the free block and the five local interaction blocks.  If the
ordinary \(4\)PN coefficients are denoted by \(L_1,\dots,L_{21}\), then the
quartic reduced COM compiler is diagonal slot by slot:
\begin{equation}
h_i=F_i(\nu)-L_i,
\qquad
\frac{\partial h_i}{\partial L_j}=-\delta_{ij}.
\label{eq:app-local-target-diagonal-com-map}
\end{equation}
So the local COM Hamiltonian-to-ordinary translation introduces no cross-slot
mixing.  The chart issue is not slot mixing; it is the new mass-ratio degree
content created inside individual slots.

The translated exact ordinary local target is as follows.

\paragraph{Free block.}
\begin{equation}
L_1^{\rm target}
=
\frac{7+71\nu-1135\nu^2+4082\nu^3-4497\nu^4}{256},
\qquad
L_2^{\rm target}=L_3^{\rm target}=L_4^{\rm target}=L_5^{\rm target}=L_6^{\rm target}=0.
\label{eq:app-local-target-ordinary-free}
\end{equation}
So the free block remains a one-slot block after quartic translation.

\paragraph{\texorpdfstring{$G/r$}{G/r} block.}
\begin{align}
L_7^{\rm target}
&=
\frac{150+512\nu-7292\nu^2+9285\nu^3+10070\nu^4}{256},
\label{eq:app-local-target-ordinary-L7}
\\
L_8^{\rm target}
&=
\frac{\nu(2414\nu^3-1723\nu^2+346\nu-16)}{64},
\label{eq:app-local-target-ordinary-L8}
\\
L_9^{\rm target}
&=
\frac{3\nu^2(270\nu^2-191\nu+30)}{128},
\label{eq:app-local-target-ordinary-L9}
\\
L_{10}^{\rm target}
&=
-\frac{5\nu^3(34\nu-13)}{64},
\label{eq:app-local-target-ordinary-L10}
\\
L_{11}^{\rm target}
&=
\frac{35\nu^3(2\nu-1)}{256}.
\label{eq:app-local-target-ordinary-L11}
\end{align}

\paragraph{\texorpdfstring{$G^2/r^2$}{G	extasciicircum2/r	extasciicircum2} block.}
\begin{align}
L_{12}^{\rm target}
&=
-\frac{7504\nu^4+22415\nu^3+3297\nu^2+340\nu-944}{256},
\label{eq:app-local-target-ordinary-L12}
\\
L_{13}^{\rm target}
&=
-\frac{\nu(17648\nu^3+17503\nu^2-12188\nu+1872)}{256},
\label{eq:app-local-target-ordinary-L13}
\\
L_{14}^{\rm target}
&=
-\frac{\nu(21888\nu^3-5105\nu^2+6915\nu-2916)}{768},
\label{eq:app-local-target-ordinary-L14}
\\
L_{15}^{\rm target}
&=
\frac{\nu(5760\nu^3-2673\nu^2+11510\nu-2952)}{1280}.
\label{eq:app-local-target-ordinary-L15}
\end{align}

\paragraph{\texorpdfstring{$G^3/r^3$}{G	extasciicircum3/r	extasciicircum3} block.}
\begin{align}
L_{16}^{\rm target}
&=
\frac{30412800\nu^4+128822400\nu^3-3987675\pi^2\nu^2+258968192\nu^2-76341312\nu-1294650\pi^2\nu+23385600}{3686400},
\label{eq:app-local-target-ordinary-L16}
\\
L_{17}^{\rm target}
&=
\frac{\nu(6144000\nu^3+15650400\nu^2+281025\pi^2\nu+5897600\nu-7914912+166050\pi^2)}{153600},
\label{eq:app-local-target-ordinary-L17}
\\
L_{18}^{\rm target}
&=
\frac{\nu(5836800\nu^3+4899200\nu^2-13348224\nu+579825\pi^2\nu-11250\pi^2+3289536)}{245760}.
\label{eq:app-local-target-ordinary-L18}
\end{align}

\paragraph{\texorpdfstring{$G^4/r^4$}{G	extasciicircum4/r	extasciicircum4} block.}
\begin{align}
L_{19}^{\rm target}
&=
-\frac{614400\nu^4+1382400\nu^3-3916025\pi^2\nu^2+40000704\nu^2-3160350\pi^2\nu+22640704\nu-1190400}{1228800},
\label{eq:app-local-target-ordinary-L19}
\\
L_{20}^{\rm target}
&=
-\frac{\nu(27648000\nu^3+89856000\nu^2+21708480\nu+7681425\pi^2\nu-5340450\pi^2+114494656)}{3686400}.
\label{eq:app-local-target-ordinary-L20}
\end{align}

\paragraph{\texorpdfstring{$G^5/r^5$}{G	extasciicircum5/r	extasciicircum5} block.}
\begin{equation}
L_{21}^{\rm target}
=
-\frac{-6430720\nu^2+555225\pi^2\nu^2-16243104\nu+1403325\pi^2\nu-14400}{230400}.
\label{eq:app-local-target-ordinary-L21}
\end{equation}

The translated ordinary target still satisfies the strict one-body gate exactly:
\begin{equation}
L_1\to\frac{7}{256},
\qquad
L_7\to\frac{75}{128},
\qquad
L_{12}\to\frac{59}{16},
\qquad
L_{16}\to\frac{203}{32},
\qquad
L_{19}\to\frac{31}{32},
\qquad
L_{21}\to\frac{1}{16},
\qquad (\nu\to0),
\label{eq:app-local-target-ordinary-onebody-gate}
\end{equation}
with every subleading slot vanishing in that same limit.  So the translation
preserves the strict one-body content exactly.  The problem it introduces is not
a one-body mismatch but new comparable-mass degree content in the ordinary
interaction blocks.

\subsection{The Hamiltonian-first lesson}
\label{app:imported-local-target-hamiltonian-first-lesson}

The comparison that matters is now exact.  From the raw local generic-frame
scaffold one already knows that the constant-coefficient interaction basis has
blockwise COM degree ceilings
\begin{equation}
\deg_{\nu}^{\rm max}(G/r,G^2/r^2,G^3/r^3,G^4/r^4,G^5/r^5)
=
(4,3,3,2,2).
\label{eq:app-local-target-scaffold-degree-ceilings}
\end{equation}
The imported local Hamiltonian target has exactly the same interaction ceilings,
so in the Hamiltonian chart the target and scaffold are degree-matched block by
block.

The translated ordinary target, however, has
\begin{equation}
\deg_{\nu}^{\rm max}(G/r,G^2/r^2,G^3/r^3,G^4/r^4,G^5/r^5)
=
(4,4,4,4,2).
\label{eq:app-local-target-ordinary-degree-ceilings}
\end{equation}
So the quartic compiler itself creates extra \(\nu^4\) tails in the ordinary
chart precisely in the middle and upper local interaction blocks,
\begin{equation}
G^2/r^2,
\qquad
G^3/r^3,
\qquad
G^4/r^4,
\label{eq:app-local-target-ordinary-extra-tail-blocks}
\end{equation}
while leaving the top static \(G^5/r^5\) block at the same ceiling as before.

This gives the exact chart lesson:
\begin{equation}
\boxed{
\text{The fixed local generic-frame lift at }4\text{PN must be performed in the Hamiltonian chart first.}
}
\label{eq:app-local-target-hamiltonian-first-lesson}
\end{equation}
The reason is not stylistic.  It is structural.  In the Hamiltonian chart, the
imported local target matches the constant-coefficient generic-frame scaffold
blockwise and can therefore be lifted directly before any further translation.
In the ordinary chart, the translated target contains new \(\nu^4\) tails that
are absent from the fixed Hamiltonian target and are not reachable from the raw
constant-coefficient ordinary scaffold.

So the correct local-first program is
\begin{equation}
\boxed{
\text{fixed local Hamiltonian target}
\;\longrightarrow\;
\text{generic-frame Hamiltonian lift}
\;\longrightarrow\;
\text{ordinary-chart translation only afterward.}
}
\label{eq:app-local-target-correct-program-order}
\end{equation}
The tail bridge remains separate throughout.  Nothing in the Hamiltonian-first
lesson relaxes the local/nonlocal split, and nothing here changes the fact that
the hereditary coefficient must be handled independently through the quadrupole
bridge of Section~\ref{sec:tail-bridge}.

At the same time, the lesson should not be overstated.  It does \emph{not} yet
settle the full local ordinary generic-frame representative.  It settles the
order in which the local problem must be attacked.  The actual Hamiltonian-chart
generic-frame residual, its canonical slice, the ordinary residual sign-flip
translation, and the aligned-seed correction are subsequent steps.  What this
appendix establishes is the chart logic that makes all of those later steps both
necessary and mathematically clean.

% Appendix: Hamiltonian-chart generic-frame lift

\section{Hamiltonian-chart generic-frame lift}
\label{app:hamiltonian-lift}

Section~\ref{sec:local-hamiltonian-lift} used only the headline conclusion of
this stage: once the fixed local target is imported in the Hamiltonian chart,
the generic-frame local $4$PN comparable-mass lift has no remaining span
obstruction.  The purpose of the present appendix is to freeze that result in
full.  Relative to the raw scaffold of
Appendix~\ref{app:local-generic-frame-scaffold-com} and the imported target of
Appendix~\ref{app:imported-local-target-hamiltonian-first}, the exact question
is the following: after subtracting the natural one-body/self-static seed, does
the constant-coefficient exchange-symmetric generic-frame scaffold span the full
fixed-chart local center-of-mass (COM) Hamiltonian residual, or does a further
local obstruction survive?

The answer is completely sharp.  In the Hamiltonian chart, the answer is yes.
The free slot is absorbed exactly by the natural Hamiltonian seed, the
interaction residual is polynomially reachable block by block, and the only
remaining ambiguity is a $92$-direction COM-blind null family.  In particular,
there is no Hamiltonian-chart analogue of the ordinary-chart $\nu^4$-tail
obstruction.  The only nontrivial task left after the present appendix is to
choose a canonical quotient representative inside the null family; that is the
role of Appendix~\ref{app:canonical-slice}.

\subsection{Natural Hamiltonian seed and exact free-slot absorption}
\label{app:hamiltonian-lift-seed}

The Hamiltonian-chart lift begins from the exact one-body $4$PN Hamiltonian
gate already frozen in
Appendix~\ref{app:imported-local-target-hamiltonian-first}.  In reduced COM
variables, write
\begin{equation}
X_A=\frac{1+\Delta}{2},
\qquad
X_B=\frac{1-\Delta}{2},
\qquad
X_A+X_B=1,
\qquad
X_AX_B=\nu.
\label{eq:app-4pn-hamiltonian-lift-XAXB}
\end{equation}
Then the natural reduced COM Hamiltonian self/static seed is
\begin{equation}
K^{\rm seed}=\frac{7}{256}\bigl(X_A^9+X_B^9\bigr),
\label{eq:app-4pn-hamiltonian-lift-Kseed}
\end{equation}
\begin{equation}
Q_1^{\rm seed}=\frac{45}{128}\bigl(X_A^8+X_B^8\bigr),
\qquad
T_1^{\rm seed}=\frac{13}{8}\bigl(X_A^7+X_B^7\bigr),
\qquad
S_1^{\rm seed}=\frac{105}{32}\bigl(X_A^6+X_B^6\bigr),
\label{eq:app-4pn-hamiltonian-lift-seed-QTS}
\end{equation}
\begin{equation}
U_1^{\rm seed}=\frac{105}{32}\bigl(X_A^5+X_B^5\bigr),
\qquad
W_1^{\rm seed}=-\frac{1}{16}\bigl(X_A^4+X_B^4\bigr),
\label{eq:app-4pn-hamiltonian-lift-seed-UW}
\end{equation}
with all subleading Hamiltonian slots initially set to zero.

This seed has one immediate exact consequence: it already absorbs the entire
free $4$PN COM Hamiltonian slot.  Writing the fixed local COM Hamiltonian target
as
\begin{equation}
H_{4,\mathrm{loc}}^{\mathrm{target,COM}}
=
K^{\mathrm{target}}
+
H_{4,\mathrm{loc}}^{\mathrm{int,target,COM}},
\label{eq:app-4pn-hamiltonian-lift-target-split}
\end{equation}
the free residual vanishes identically,
\begin{equation}
\Delta K
\equiv
K^{\mathrm{target}}-K^{\rm seed}
=
0.
\label{eq:app-4pn-hamiltonian-lift-DeltaK-zero}
\end{equation}
So the Hamiltonian-chart generic-frame lift problem is purely an interaction
problem.  No extra Hamiltonian free-slot repair survives beyond the natural
one-body seed.

\subsection{Exact polynomial coefficient-space formulation}
\label{app:hamiltonian-lift-coefficient-space}

The generic-frame Hamiltonian scaffold uses the same formal block families as the
raw local scaffold,
\begin{equation}
\frac{Gpq}{r}\,\mathcal Q,
\qquad
\frac{G^2pq}{r^2}\,\mathcal T,
\qquad
\frac{G^3pq}{r^3}\,\mathcal S,
\qquad
\frac{G^4}{r^4}\,\mathcal U,
\qquad
\frac{G^5}{r^5}\,\mathcal W,
\label{eq:app-4pn-hamiltonian-lift-block-families}
\end{equation}
with the same raw block dimensions
\begin{equation}
\dim\mathcal Q=52,
\qquad
\dim\mathcal T=46,
\qquad
\dim\mathcal S=29,
\qquad
\dim\mathcal U=10,
\qquad
\dim\mathcal W=2.
\label{eq:app-4pn-hamiltonian-lift-raw-dims}
\end{equation}
The only change is interpretive: in the present appendix these are read as
Hamiltonian-side monomials in the Newtonian-order momentum invariants rather
than as ordinary-chart velocity monomials.

Because the imported local Hamiltonian target has degree ceilings
\begin{equation}
\deg_{\nu}^{\rm max}
\bigl(G/r,\,G^2/r^2,\,G^3/r^3,\,G^4/r^4,\,G^5/r^5\bigr)
=
(4,3,3,2,2),
\label{eq:app-4pn-hamiltonian-lift-degree-ceilings}
\end{equation}
exactly as recorded in
Appendix~\ref{app:imported-local-target-hamiltonian-first}, the strict lift
question is stronger than slot counting alone.  One must ask whether the raw
Hamiltonian scaffold spans the full \emph{polynomial coefficient space} of the
imported COM target.  The blockwise coefficient-space dimensions are therefore
\begin{equation}
\begin{aligned}
\frac{G}{r}:&\quad 5\text{ slots}\times(\nu,\nu^2,\nu^3,\nu^4)
&&\Rightarrow 20\text{ coefficients},\\
\frac{G^2}{r^2}:&\quad 4\text{ slots}\times(\nu,\nu^2,\nu^3)
&&\Rightarrow 12,\\
\frac{G^3}{r^3}:&\quad 3\text{ slots}\times(\nu,\nu^2,\nu^3)
&&\Rightarrow 9,\\
\frac{G^4}{r^4}:&\quad 2\text{ slots}\times(\nu,\nu^2)
&&\Rightarrow 4,\\
\frac{G^5}{r^5}:&\quad 1\text{ slot}\times(\nu,\nu^2)
&&\Rightarrow 2.
\end{aligned}
\label{eq:app-4pn-hamiltonian-lift-coefficient-space}
\end{equation}
Hence the full interaction coefficient space has total dimension
\begin{equation}
20+12+9+4+2=47.
\label{eq:app-4pn-hamiltonian-lift-total-rank-target}
\end{equation}

This formulation is the correct exact lift test.  The issue is not whether the
scaffold spans fifteen reduced interaction slots after COM projection.  The
issue is whether it spans all forty-seven independent mass-ratio coefficients of
the fixed local COM Hamiltonian residual.

\subsection{Blockwise exact rank test and the Hamiltonian lift theorem}
\label{app:hamiltonian-lift-rank-test}

With the coefficient-space formulation fixed, the blockwise Hamiltonian-chart
rank test is exact:
\begin{equation}
\begin{aligned}
\frac{G}{r}:&\quad 52\to 20, &&\text{nullity }32,\\
\frac{G^2}{r^2}:&\quad 46\to 12, &&\text{nullity }34,\\
\frac{G^3}{r^3}:&\quad 29\to 9, &&\text{nullity }20,\\
\frac{G^4}{r^4}:&\quad 10\to 4, &&\text{nullity }6,\\
\frac{G^5}{r^5}:&\quad 2\to 2, &&\text{nullity }0.
\end{aligned}
\label{eq:app-4pn-hamiltonian-lift-ranks}
\end{equation}
So the interaction coefficient-space rank is maximal in every block, and the
combined totals are
\begin{equation}
\operatorname{rank}=47,
\qquad
\operatorname{nullity}=32+34+20+6+0=92.
\label{eq:app-4pn-hamiltonian-lift-total-rank-nullity}
\end{equation}

The local Hamiltonian lift theorem is then immediate.
\begin{theorem}[Hamiltonian-chart generic-frame lift theorem]
\label{thm:app-4pn-hamiltonian-chart-lift}
Let the fixed-chart local COM Hamiltonian target be imported with the natural
Hamiltonian seed
\eqref{eq:app-4pn-hamiltonian-lift-Kseed}--\eqref{eq:app-4pn-hamiltonian-lift-seed-UW}
and let the interaction residual be defined by subtracting that seed.  Then the
constant-coefficient exchange-symmetric generic-frame Hamiltonian scaffold
spans the full local comparable-mass interaction residual.  Equivalently, there
exists at least one generic-frame Hamiltonian representative
\begin{equation}
\Delta H_{4,\mathrm{loc}}^{\mathrm{gen}}
=
\frac{Gpq}{r}Q_{\mathrm{gen}}
+
\frac{G^2pq}{r^2}T_{\mathrm{gen}}
+
\frac{G^3pq}{r^3}S_{\mathrm{gen}}
+
\frac{G^4}{r^4}U_{\mathrm{gen}}
+
\frac{G^5}{r^5}W_{\mathrm{gen}}
\label{eq:app-4pn-hamiltonian-lift-generic-residual}
\end{equation}
such that
\begin{equation}
\Pi_{\mathrm{COM}}\!\left[\Delta H_{4,\mathrm{loc}}^{\mathrm{gen}}\right]
=
\Delta H_{4,\mathrm{loc}}^{\mathrm{target,COM}}
\label{eq:app-4pn-hamiltonian-lift-referee-identity}
\end{equation}
slot by slot and coefficient by coefficient.  No further Hamiltonian seed
refinement is required before the generic-frame comparable-mass lift is
performed.
\end{theorem}

\begin{proof}
Equation~\eqref{eq:app-4pn-hamiltonian-lift-DeltaK-zero} removes the free slot,
so only the interaction residual remains.  The interaction coefficient space has
dimension $47$ by
\eqref{eq:app-4pn-hamiltonian-lift-total-rank-target}.  The exact blockwise rank
test \eqref{eq:app-4pn-hamiltonian-lift-ranks} shows that the generic-frame
Hamiltonian scaffold attains the full coefficient-space rank in every block.
Hence the image of the scaffold coincides with the entire imported local COM
Hamiltonian residual, which proves existence of at least one solution of
\eqref{eq:app-4pn-hamiltonian-lift-referee-identity}.  Because maximal rank is
already achieved with the natural Hamiltonian seed fixed, no additional
Hamiltonian-side seed dressing can be necessary at the level of comparable-mass
existence.  The only remaining freedom is the $92$-direction kernel recorded in
\eqref{eq:app-4pn-hamiltonian-lift-total-rank-nullity}.
\end{proof}

Two consequences are worth stating explicitly.
First, the previous ordinary-chart obstruction is not a fundamental local span
failure.  It is a chart artifact created by quartic translation.  Second, the
Hamiltonian-chart lift theorem is stronger than the earlier COM slot-counting
statement: the lift closes not merely at the level of fifteen reduced
interaction slots, but at the level of all forty-seven independent mass-ratio
coefficients carried by those slots.

\subsection{Why span is solved but quotienting is not}
\label{app:hamiltonian-lift-transition}

The present appendix closes the Hamiltonian-chart existence question and nothing
more.  The rank theorem shows that the generic-frame scaffold is large enough,
but it also shows that the lift is highly nonunique.  The exact residual family
contains a $92$-direction COM-blind null sector, so after
\eqref{eq:app-4pn-hamiltonian-lift-referee-identity} is solved one still has to
choose a distinguished representative.  That is not a span problem; it is a
quotient problem.

This is precisely why the next appendix is logically separate.  The task of
Appendix~\ref{app:canonical-slice} is to identify the COM-null ideals, show
that the null family is purely algebraic rather than physical, and select one
exact sparse canonical quotient representative.  The present appendix does not
need that construction.  Its theorem is already complete once the rank test is
finished.

Finally, the chart lesson can now be stated in fully exact form.  The generic-
frame local lift has to be done in the Hamiltonian chart because only there do
the imported degree ceilings match the constant-coefficient scaffold exactly.
The translated ordinary target carries extra $\nu^4$ tails in the middle and
upper blocks, whereas the Hamiltonian target does not.  So the correct local-
first route is
\begin{equation}
\boxed{\begin{gathered}
\text{fixed local Hamiltonian target}
\;\longrightarrow\;
\text{generic-frame Hamiltonian lift}
\\
\text{canonical quotient slice}
\;\longrightarrow\;
\text{ordinary translation only afterward}
\end{gathered}}
\label{eq:app-4pn-hamiltonian-lift-program}
\end{equation}
That is the precise carry-forward conclusion of the present appendix.


% Appendix: Canonical comparable-mass local residual slice

\section{Canonical comparable-mass local residual slice}
\label{app:canonical-slice}

Appendix~\ref{app:hamiltonian-lift} established the exact Hamiltonian-chart lift
theorem: after subtracting the natural one-body/self-static seed, the full local
$4$PN comparable-mass Hamiltonian residual lies in the image of the
constant-coefficient exchange-symmetric generic-frame scaffold, block by block.
What remained open was no longer span, but quotient choice.  The COM projection
forgets a large family of generic-frame directions, so one must decide how to
choose a clean representative inside that COM-blind null family.

The present appendix freezes that choice.  It identifies the relevant COM-null
ideals, records the exact $92$-direction ambiguity, and then fixes one
canonical representative by the simplest exact quotient rule available: in the
Gauss--Jordan lift of each block, set every null coordinate equal to zero.
This produces a sparse generic-frame Hamiltonian representative of the full
local comparable-mass residual.  The representative is exact, reduces to the
imported fixed-chart COM target in every slot, and is the canonical starting
point for the subsequent ordinary-chart translation.

\subsection{COM-null ideals and the 92-direction ambiguity}
\label{app:canonical-slice-null-ideals}

We use the same formal generic-frame invariant symbols
\(
a,b,c,d,e,p,q
\)
as in Appendix~\ref{app:hamiltonian-lift}.  The COM identities are
\begin{equation}
\mathcal C_1 = p a + q c,
\qquad
\mathcal C_2 = p c + q b,
\qquad
\mathcal C_3 = p d + q e,
\label{eq:app-canonical-slice-C123}
\end{equation}
\begin{equation}
\mathcal C_4 = ab-c^2,
\qquad
\mathcal C_5 = ae-cd,
\qquad
\mathcal C_6 = bd-ce.
\label{eq:app-canonical-slice-C456}
\end{equation}
These are the same algebraic COM-blind generators already familiar from the
$3$PN analysis.  At $4$PN, however, the null family is substantially larger.

From Appendix~\ref{app:hamiltonian-lift}, the blockwise Hamiltonian-chart
nullities are
\begin{equation}
Q:32,
\qquad
T:34,
\qquad
S:20,
\qquad
U:6,
\qquad
W:0,
\label{eq:app-canonical-slice-block-nullities}
\end{equation}
so the total COM-blind ambiguity is
\begin{equation}
32+34+20+6+0 = 92.
\label{eq:app-canonical-slice-total-nullity}
\end{equation}
The exact ideal-membership statement is sharper than the raw nullity count:
\begin{equation}
\mathcal I_{\rm COM}^{\rm full}
=
\langle \mathcal C_1,\mathcal C_2,\mathcal C_3,\mathcal C_4,\mathcal C_5,\mathcal C_6\rangle,
\qquad
\mathcal I_{\rm COM}^{\rm lin}
=
\langle \mathcal C_1,\mathcal C_2,\mathcal C_3\rangle.
\label{eq:app-canonical-slice-ideals}
\end{equation}

\begin{theorem}[Exact COM-null ideal theorem at local $4$PN]
\label{thm:app-canonical-slice-null-ideal}
For the Hamiltonian-chart local $4$PN generic-frame scaffold, every $Q$-block
null vector lies in the full COM ideal \(\mathcal I_{\rm COM}^{\rm full}\),
while every \(T\)-, \(S\)-, and \(U\)-block null vector already lies in the
linear-momentum ideal \(\mathcal I_{\rm COM}^{\rm lin}\).  The \(W\)-block has
zero nullity and is therefore fixed completely once the COM target is fixed.
\end{theorem}

This theorem is the precise algebraic meaning of the statement that the
Hamiltonian lift ambiguity is \emph{purely} COM-blind.  No further span defect
or hidden upper-block freedom survives.  The only remaining task is to pick one
representative in the quotient by these ideals.

\subsection{Exact Gauss--Jordan quotient slice}
\label{app:canonical-slice-gauss-jordan}

Let \(B\in\{Q,T,S,U,W\}\) denote one interaction block.  After choosing a raw
generic-frame scaffold basis for that block, the exact coefficient-space lift is
an affine linear system
\begin{equation}
M_B x_B = y_B,
\label{eq:app-canonical-slice-linear-system}
\end{equation}
where \(M_B\) is the blockwise COM image matrix and \(y_B\) is the imported COM
Hamiltonian residual written in coefficient space.  Because the ranks are
maximal in every block, the solution set has the exact form
\begin{equation}
x_B
=
x_B^{\rm part}
+
\sum_{\alpha=1}^{n_B}\lambda_{B,\alpha}\,n_{B,\alpha},
\label{eq:app-canonical-slice-affine-family}
\end{equation}
where \(x_B^{\rm part}\) is one particular solution, \(\{n_{B,\alpha}\}\) is a
basis of the null space, and
\begin{equation}
n_Q=32,
\qquad
n_T=34,
\qquad
n_S=20,
\qquad
n_U=6,
\qquad
n_W=0.
\label{eq:app-canonical-slice-null-dims}
\end{equation}

The canonical quotient slice used throughout the local $4$PN chain is the
simplest exact choice:

\begin{equation}
\lambda_{B,\alpha}=0
\qquad
\text{for every block \(B\) and every null coordinate \(\alpha\).}
\label{eq:app-canonical-slice-choice}
\end{equation}

Equivalently, one takes the exact Gauss--Jordan lift and discards every purely
COM-blind null direction.  This is not claimed as the only possible quotient
choice; it is claimed to be the cleanest canonical one available once the
Hamiltonian chart and the fixed COM target are frozen.

A useful structural consequence is sparsity.  The resulting canonical
representative uses only
\begin{equation}
Q:11,
\qquad
T:12,
\qquad
S:9,
\qquad
U:4,
\qquad
W:2
\label{eq:app-canonical-slice-sparse-counts}
\end{equation}
nonzero scaffold directions.  So the quotient choice removes the entire
\(92\)-direction COM-blind family while preserving one exact sparse
representative of the same COM target.

\subsection{Explicit canonical representative}
\label{app:canonical-slice-explicit}

Write the local comparable-mass Hamiltonian residual as
\begin{equation}
\Delta H_{4,\mathrm{loc}}^{\mathrm{can}}
=
\frac{Gpq}{r}Q_{\mathrm{can}}
+
\frac{G^2pq}{r^2}T_{\mathrm{can}}
+
\frac{G^3pq}{r^3}S_{\mathrm{can}}
+
\frac{G^4}{r^4}U_{\mathrm{can}}
+
\frac{G^5}{r^5}W_{\mathrm{can}}.
\label{eq:app-canonical-slice-representative}
\end{equation}
Then one exact canonical quotient-slice representative is the following.

\paragraph{$G/r$ block.}
\begin{equation}
\begin{aligned}
Q_{\mathrm{can}}={}&
-\frac{11}{64}a^2b^2
+\frac{5}{256}(a^2bc+ab^2c)
-\frac{3}{32}(a^2bd^2+ab^2e^2)
+\frac{1}{64}(a^2bde+ab^2de)
\\
&-\frac{27}{64}(a^2c^2+b^2c^2)
-\frac{9}{64}(a^2d^2e^2+b^2d^2e^2)
+\frac{3}{128}(a^2de^3+b^2d^3e)
+\frac{3}{64}(a^2e^4+b^2d^4)
\\
&-\frac{5}{32}(ad^2e^4+bd^4e^2)
+\frac{5}{64}(ad^3e^3+bd^3e^3)
-\frac{35}{256}(d^5e^3+d^3e^5).
\end{aligned}
\label{eq:app-canonical-slice-Qcan}
\end{equation}

\paragraph{$G^2/r^2$ block.}
\begin{equation}
\begin{aligned}
T_{\mathrm{can}}={}&
-\frac{87}{128}(a^2bp+ab^2q)
-\frac{2227}{256}(a^2bq+ab^2p)
+\frac{63}{64}(a^2cq+b^2cp)
\\
&+\frac{49}{16}(a^2d^2p+b^2e^2q)
-\frac{435}{64}(a^2dep+b^2deq)
+\frac{2435}{256}(a^2e^2p+b^2d^2q)
\\
&-\frac{183}{16}(ad^2e^2p+bd^2e^2q)
-\frac{8305}{768}(ad^2e^2q+bd^2e^2p)
+\frac{889}{192}(ad^3eq+bde^3p)
\\
&-\frac{5343}{1280}(d^3e^3p+d^3e^3q)
+\frac{325}{128}(d^4e^2q+d^2e^4p)
-\frac{369}{160}(d^5eq+de^5p).
\end{aligned}
\label{eq:app-canonical-slice-Tcan}
\end{equation}

\paragraph{$G^3/r^3$ block.}
\begin{equation}
\begin{aligned}
S_{\mathrm{can}}={}&
\left(-\frac{3307423}{28800}+\frac{40483\pi^2}{16384}\right)(a^2p^2+b^2q^2)
+\left(-\frac{211189}{19200}+\frac{2749\pi^2}{8192}\right)(a^2pq+b^2pq)
\\
&+\left(-\frac{184897}{900}+\frac{34985\pi^2}{8192}\right)abpq
+\left(\frac{139241}{1200}-\frac{12507\pi^2}{2048}\right)(ad^2p^2+be^2q^2)
\\
&+\left(\frac{63347}{1600}-\frac{1059\pi^2}{1024}\right)(ad^2pq+be^2pq)
+\left(-\frac{716971}{9600}+\frac{10389\pi^2}{2048}\right)(adep^2+bdeq^2)
\\
&+\left(-\frac{16727}{384}-\frac{35655\pi^2}{16384}\right)(d^2e^2p^2+d^2e^2q^2)
+\left(-\frac{51193}{960}-\frac{36405\pi^2}{8192}\right)d^2e^2pq
\\
&+\left(\frac{23533}{1280}-\frac{375\pi^2}{8192}\right)(d^3eq^2+de^3p^2).
\end{aligned}
\label{eq:app-canonical-slice-Scan}
\end{equation}

\paragraph{$G^4/r^4$ block.}
\begin{equation}
\begin{aligned}
U_{\mathrm{can}}={}&
\left(\frac{930047}{9600}-\frac{551243\pi^2}{49152}\right)(ap^2q+bpq^2)
+\left(\frac{500761}{19200}-\frac{21837\pi^2}{8192}\right)(apq^2+bp^2q)
\\
&+\left(\frac{274387}{1600}-\frac{62047\pi^2}{49152}\right)(d^2p^2q+e^2pq^2)
+\left(\frac{3401779}{57600}-\frac{28691\pi^2}{24576}\right)(d^2pq^2+e^2p^2q).
\end{aligned}
\label{eq:app-canonical-slice-Ucan}
\end{equation}

\paragraph{$G^5/r^5$ block.}
\begin{equation}
W_{\mathrm{can}}=
\left(-\frac{169799}{2400}+\frac{6237\pi^2}{1024}\right)(p^3q+pq^3)
+\left(-\frac{609427}{3600}+\frac{44825\pi^2}{3072}\right)p^2q^2.
\label{eq:app-canonical-slice-Wcan}
\end{equation}

These formulas are best read as a theorem witness.  They are not a heuristic
guess and they are not an approximate sparse fit.  They are the exact output of
the canonical quotient rule \eqref{eq:app-canonical-slice-choice} applied to the
full Hamiltonian-chart lift family.

\subsection{Exact verification and status boundary}
\label{app:canonical-slice-verification}

The decisive verification is immediate: after COM reduction, the representative
\eqref{eq:app-canonical-slice-representative} reproduces the imported fixed-chart
local COM Hamiltonian residual in every block and every slot:
\begin{equation}
\Pi_{\mathrm{COM}}\!\left[\Delta H_{4,\mathrm{loc}}^{\mathrm{can}}\right]
=
\Delta H_{4,\mathrm{loc}}^{\mathrm{target,COM}}.
\label{eq:app-canonical-slice-referee-identity}
\end{equation}
More explicitly, all five \(Q\)-slots, all four \(T\)-slots, all three
\(S\)-slots, both \(U\)-slots, and the single \(W\)-slot match exactly.

\begin{theorem}[Canonical comparable-mass local residual slice]
\label{thm:app-canonical-slice}
Within the Hamiltonian chart, the local $4$PN comparable-mass residual admits an
exact sparse generic-frame representative obtained by setting all COM-null
coordinates to zero in the blockwise Gauss--Jordan lift.  The resulting
representative is given by
\eqref{eq:app-canonical-slice-representative}--\eqref{eq:app-canonical-slice-Wcan},
contains only the nonzero scaffold counts
\eqref{eq:app-canonical-slice-sparse-counts}, and reduces exactly to the
imported fixed-chart local COM Hamiltonian residual.
\end{theorem}

What this appendix settles is therefore precise.  The Hamiltonian-chart local
existence problem is no longer open; it has been quotiented, canonically sliced,
and frozen in explicit form.  What it does \emph{not} settle is the stronger
question of fixed-chart uniqueness against the true full generic-frame ADM
Hamiltonian target.  For the purposes of the local-first program, however, that
stronger issue is no longer the next bottleneck.  The next clean move is the
ordinary-chart translation of the canonical representative, with the hereditary
tail bridge kept separate.

% Appendix: Ordinary translation and aligned-seed correction

\section{Ordinary translation and aligned-seed correction}
\label{app:ordinary-translation-aligned-seed}

Appendices~\ref{app:imported-local-target-hamiltonian-first},
\ref{app:hamiltonian-lift}, and \ref{app:canonical-slice} already settled the
local $4$PN problem in the chart in which it is naturally posed.  The fixed
instantaneous target was imported in the reduced COM Hamiltonian chart, the
constant-coefficient exchange-symmetric generic-frame scaffold was shown to span
its full comparable-mass interaction residual, and one sparse canonical
Hamiltonian representative was selected by quotienting the $92$-direction
COM-blind null family.  What remained open after those steps was not the
Hamiltonian-side existence problem, but the exact relation between that
canonical Hamiltonian representative and the ordinary local chart.

This appendix freezes that relation.  There are two logically separate claims.
First, once the lower-order ledger is held fixed and the ordinary $4$PN seed is
chosen in the \emph{Hamiltonian-aligned} convention, the quartic comparable-mass
residual translates by an exact sign flip,
\begin{equation}
\boxed{
\Delta L_{4,\mathrm{loc}}^{\mathrm{can}}
=
-\Delta H_{4,\mathrm{loc}}^{\mathrm{can}}.
}
\label{eq:app-ordinary-translation-sign-flip-preview}
\end{equation}
Second, the earlier ordinary-chart obstruction is therefore not a residual-lift
obstruction at all.  It is the separate chart problem carried by the
\emph{seed-alignment correction}
\begin{equation}
\delta L_{4,\mathrm{seed}}
=
L_{4,\mathrm{seed}}^{(\mathrm{aligned})}
-
L_{4,\mathrm{seed}}^{(\mathrm{natural})}.
\label{eq:app-ordinary-translation-delta-seed-preview}
\end{equation}
In particular, the true local ordinary obstruction is seed-sector in origin and
remains completely separate from the hereditary/tail bridge.

\subsection{Exact quartic residual compiler theorem}
\label{app:ordinary-translation-residual-compiler}

The quartic perturbative Legendre transform of
Appendix~\ref{app:quartic-legendre-identity} has the schematic structure
\begin{equation}
H_4 = F[L_1,L_2,L_3](v_0) - L_4(v_0),
\qquad
v_0=M^{-1}p,
\label{eq:app-ordinary-translation-general-form}
\end{equation}
where the feedback functional $F$ depends only on the already frozen lower-order
ordinary data.  Therefore if the ordinary $4$PN block is split as
\begin{equation}
L_4 = L_4^{\mathrm{seed}} + \Delta L_4,
\label{eq:app-ordinary-translation-seed-residual-split}
\end{equation}
and the seed is chosen to be the one obtained by translating the fixed local
Hamiltonian seed through the exact quartic compiler, then the residual is the
only genuinely new quartic object and it obeys an exact sign flip.

\begin{theorem}[Quartic ordinary residual sign-flip theorem]
\label{thm:app-ordinary-translation-sign-flip}
Let the local $4$PN ordinary chart be defined relative to the
Hamiltonian-aligned seed.  Then the exact comparable-mass local residual
satisfies
\begin{equation}
\boxed{
\Delta H_4 = -\Delta L_4.
}
\label{eq:app-ordinary-translation-sign-flip-theorem}
\end{equation}
In particular, for the canonical sparse representative selected in
Appendix~\ref{app:canonical-slice},
\begin{equation}
\boxed{
\Delta L_{4,\mathrm{loc}}^{\mathrm{can}}
=
-\Delta H_{4,\mathrm{loc}}^{\mathrm{can}}.
}
\label{eq:app-ordinary-translation-canonical-sign-flip}
\end{equation}
\end{theorem}

\begin{proof}
Equation~\eqref{eq:app-ordinary-translation-general-form} is affine in the new
quartic ordinary block $L_4$ and all lower-order feedback is frozen.  If the
quartic block is decomposed as in
\eqref{eq:app-ordinary-translation-seed-residual-split}, then the seed is
chosen precisely so that its quartic counterimage reproduces the fixed local
Hamiltonian self/static package.  Subtracting the seeded relation from the full
quartic relation leaves only the residual piece, and the dependence on
$\Delta L_4$ enters linearly with coefficient $-1$.  Hence
\eqref{eq:app-ordinary-translation-sign-flip-theorem} follows exactly, and
\eqref{eq:app-ordinary-translation-canonical-sign-flip} is its canonical local
specialization.
\end{proof}

The practical consequence is immediate.  The canonical sparse Hamiltonian
representative
\begin{equation}
\Delta H_{4,\mathrm{loc}}^{\mathrm{can}}
=
\frac{Gpq}{r}Q_{\mathrm{can}}
+
\frac{G^2pq}{r^2}T_{\mathrm{can}}
+
\frac{G^3pq}{r^3}S_{\mathrm{can}}
+
\frac{G^4}{r^4}U_{\mathrm{can}}
+
\frac{G^5}{r^5}W_{\mathrm{can}}
\label{eq:app-ordinary-translation-canonical-hamiltonian-residual-repeat}
\end{equation}
constructed in Appendix~\ref{app:canonical-slice} already determines a
canonical sparse ordinary representative,
\begin{equation}
\Delta L_{4,\mathrm{loc}}^{\mathrm{can}}
=
-
\frac{Gpq}{r}Q_{\mathrm{can}}
-
\frac{G^2pq}{r^2}T_{\mathrm{can}}
-
\frac{G^3pq}{r^3}S_{\mathrm{can}}
-
\frac{G^4}{r^4}U_{\mathrm{can}}
-
\frac{G^5}{r^5}W_{\mathrm{can}}.
\label{eq:app-ordinary-translation-canonical-ordinary-residual}
\end{equation}
So the comparable-mass local residual itself does not generate any new ordinary
uncertainty once the aligned seed convention has been fixed.

\subsection{Exact COM verification of the translated residual}
\label{app:ordinary-translation-com-verification}

The translated canonical ordinary residual must be compared with the exact fixed
local COM ordinary target reconstructed in
Appendix~\ref{app:imported-local-target-hamiltonian-first}.  The essential
point is that the comparison is not made relative to the earlier natural
one-body ordinary seed, but relative to the \emph{ordinary seed aligned with the
chosen Hamiltonian self/static seed}.  With that aligned seed, the COM
comparison is exact block by block:
\begin{equation}
\frac{G}{r}:\ 5\ \text{slots match},
\qquad
\frac{G^2}{r^2}:\ 4\ \text{slots match},
\qquad
\frac{G^3}{r^3}:\ 3\ \text{slots match},
\label{eq:app-ordinary-translation-block-match-1}
\end{equation}
\begin{equation}
\frac{G^4}{r^4}:\ 2\ \text{slots match},
\qquad
\frac{G^5}{r^5}:\ 1\ \text{slot matches}.
\label{eq:app-ordinary-translation-block-match-2}
\end{equation}
So the comparable-mass residual survives quartic translation intact.  The exact
ordinary residual is already COM-correct before any further generic-frame rescue
is attempted.

It is useful to record the degree ceilings of the translated residual in the
same block language.  Writing the five local interaction blocks in the order
$Q,T,S,U,W$, the ordinary residual has ceilings
\begin{equation}
Q:(4,4,4,4,3),
\qquad
T:(3,3,3,3),
\qquad
S:(3,3,3),
\qquad
U:(2,2),
\qquad
W:(2).
\label{eq:app-ordinary-translation-residual-degree-ceilings}
\end{equation}
Equivalently, the quartic ordinary translation introduces no new comparable-mass
polynomial degree beyond what the Hamiltonian-chart lift already spans.  This is
another way of saying that the residual itself is not the source of the earlier
ordinary-chart confusion.

\subsection{The exact seed decomposition of the local ordinary target}
\label{app:ordinary-translation-seed-decomposition}

Once the residual is isolated, the exact local ordinary target decomposes as
\begin{equation}
\boxed{
L_{4,\mathrm{target}}^{\mathrm{loc}}
=
L_{4,\mathrm{seed}}^{\mathrm{natural}}
+
\delta L_{4,\mathrm{seed}}
+
\Delta L_{4,\mathrm{loc}}^{\mathrm{can}},
}
\label{eq:app-ordinary-translation-target-decomposition}
\end{equation}
with
\begin{equation}
\delta L_{4,\mathrm{seed}}
\equiv
L_{4,\mathrm{seed}}^{(\mathrm{aligned})}
-
L_{4,\mathrm{seed}}^{(\mathrm{natural})}.
\label{eq:app-ordinary-translation-seed-alignment-definition}
\end{equation}
Here
$L_{4,\mathrm{seed}}^{(\mathrm{natural})}$ means the direct symmetric promotion
of the exact one-body ordinary gate, while
$L_{4,\mathrm{seed}}^{(\mathrm{aligned})}$ means the ordinary self/static seed
selected by Hamiltonian-first quartic translation of the fixed local seed.
Equation~\eqref{eq:app-ordinary-translation-target-decomposition} is the clean
post-translation statement of the local ordinary problem.

The free block is already fixed exactly by the aligned seed.  In the twenty-one
slot reduced COM ordinary basis of
Appendix~\ref{app:imported-local-target-hamiltonian-first}, one has
\begin{equation}
L_1^{(\mathrm{aligned})}=L_1^{\mathrm{target}},
\qquad
L_2^{(\mathrm{aligned})}=L_3^{(\mathrm{aligned})}=L_4^{(\mathrm{aligned})}=L_5^{(\mathrm{aligned})}=L_6^{(\mathrm{aligned})}=0.
\label{eq:app-ordinary-translation-aligned-free-block}
\end{equation}
So the free local ordinary sector does not need a separate generic-frame repair.
The only local ordinary mismatch not already absorbed by the translated residual
is the seed-alignment correction itself.

At COM level, the \emph{interaction} part of $\delta L_{4,\mathrm{seed}}$ is
nonzero in the four lower interaction blocks and vanishes in the top static
block:
\begin{equation}
Q_\delta\neq 0,
\qquad
T_\delta\neq 0,
\qquad
S_\delta\neq 0,
\qquad
U_\delta\neq 0,
\qquad
W_\delta=0.
\label{eq:app-ordinary-translation-delta-block-support}
\end{equation}
The pure-kinetic correction is separate and is not part of the local interaction
lift.  For the interaction part, the decisive blockwise mass-ratio ceilings are
\begin{equation}
\deg_\nu^{\max}(Q_\delta,T_\delta,S_\delta,U_\delta,W_\delta)
=
(4,4,4,4,0).
\label{eq:app-ordinary-translation-delta-degree-pattern}
\end{equation}
This statement should be read together with the fixed-target import of
Appendix~\ref{app:imported-local-target-hamiltonian-first}.  The point is not
that the $G/r$ block reaches degree $4$---that was already allowed.  The point
is that the aligned-seed correction newly injects \emph{ordinary-chart}
$\nu^4$ content into the middle and upper local blocks
\begin{equation}
\frac{G^2}{r^2},
\qquad
\frac{G^3}{r^3},
\qquad
\frac{G^4}{r^4},
\label{eq:app-ordinary-translation-middle-upper-blocks}
\end{equation}
which is precisely the pattern that looked like an ordinary-chart obstruction in
the earlier direct translation audit.

\begin{corollary}[The ordinary obstruction is seed-sector only]
\label{cor:app-ordinary-translation-seed-sector-only}
The quartic ordinary-chart obstruction that appears after direct translation of
the local target is not a failure of the comparable-mass generic-frame lift.
It is exactly the seed-alignment correction
$\delta L_{4,\mathrm{seed}}$ of
\eqref{eq:app-ordinary-translation-seed-alignment-definition}.  Equivalently,
after the Hamiltonian-chart comparable-mass residual has been translated by the
sign-flip theorem, the only remaining local ordinary chart problem is the
generic-frame realization of the aligned ordinary seed.
\end{corollary}

\begin{proof}
The translated residual already matches the exact fixed-chart ordinary target in
all interaction blocks when measured relative to the aligned seed, by
\eqref{eq:app-ordinary-translation-block-match-1}--\eqref{eq:app-ordinary-translation-block-match-2}.
The aligned seed already fixes the entire free block by
\eqref{eq:app-ordinary-translation-aligned-free-block}.  Therefore the exact
difference between the natural-seed decomposition and the aligned-seed
decomposition is carried solely by
$\delta L_{4,\mathrm{seed}}$.  Its support and degree pattern are given by
\eqref{eq:app-ordinary-translation-delta-block-support} and
\eqref{eq:app-ordinary-translation-delta-degree-pattern}, so there is no second
local ordinary obstruction beyond that seed-sector object.
\end{proof}

\subsection{What this leaves for the next appendix}
\label{app:ordinary-translation-next-step}

The local-first $4$PN program is therefore split sharply into three parts:
\begin{equation}
\boxed{\text{(A) Hamiltonian-derived generic-frame residual: solved},}
\label{eq:app-ordinary-translation-split-A}
\end{equation}
\begin{equation}
\boxed{\text{(B) ordinary seed-alignment correction: the remaining local chart problem},}
\label{eq:app-ordinary-translation-split-B}
\end{equation}
\begin{equation}
\boxed{\text{(C) tail / hereditary quadrupole bridge: still completely separate}.}
\label{eq:app-ordinary-translation-split-C}
\end{equation}
Part~(A) is finished by the Hamiltonian-chart lift and the exact sign-flip
translation theorem.  Part~(C) belongs to the hereditary bridge and is untouched
by the present appendix.  The only local chart work that remains is therefore
Part~(B): realizing
$\delta L_{4,\mathrm{seed}}$ as an exact generic-frame object.

That is the content of the next local appendix.  Once the aligned-seed
correction is lifted in generic-frame form and added back to the natural seed,
the already translated canonical residual can be inserted without further local
ambiguity.  The ordinary-chart story then closes, and the remaining $4$PN
conditionality is pushed entirely onto the separate passive/outgoing quadrupole
normalization problem.

% Appendix: Structured generic-frame lift of the aligned-seed correction

\section{Structured generic-frame lift of the aligned-seed correction}
\label{app:structured-aligned-seed-lift}

Appendix~\ref{app:ordinary-translation-aligned-seed} already isolated the
exact local ordinary-chart problem.  Once the lower-order ledger is frozen and
one works relative to the Hamiltonian-aligned seed, the canonical comparable-mass
local residual translates by exact sign flip and is therefore no longer the
source of the ordinary-chart obstruction.  What remains is the separate
\emph{seed-alignment correction}
\begin{equation}
\delta L_{4,\mathrm{seed}}
=
L_{4,\mathrm{seed}}^{(\mathrm{aligned})}
-
L_{4,\mathrm{seed}}^{(\mathrm{natural})}.
\label{eq:app-structured-aligned-seed-delta-seed-def}
\end{equation}
The purpose of the present appendix is to freeze the exact generic-frame lift of
that correction.  The conclusion is completely sharp: the earlier ordinary-chart
obstruction is real, but it is narrowly seed-sector in origin, and it is removed
exactly by a small structured enlargement of the seed spaces.  Once that lift is
performed, the natural generic-frame local seed, the lifted aligned-seed
correction, and the already solved canonical ordinary residual combine to give
one exact generic-frame ordinary candidate for the entire local instantaneous
$4$PN sector.

\subsection{The actual local ordinary problem}
\label{app:structured-aligned-seed-problem}

The local interaction part of
\eqref{eq:app-structured-aligned-seed-delta-seed-def} lives in the same block
language used throughout the local $4$PN chain:
\begin{equation}
\delta L_{4,\mathrm{seed}}^{\mathrm{loc}}
=
\frac{G}{r}Q_\delta
+
\frac{G^2}{r^2}T_\delta
+
\frac{G^3}{r^3}S_\delta
+
\frac{G^4}{r^4}U_\delta
+
\frac{G^5}{r^5}W_\delta.
\label{eq:app-structured-aligned-seed-block-decomposition}
\end{equation}
At COM level this correction is nonzero in the lower four interaction blocks and
vanishes in the top static block:
\begin{equation}
Q_\delta\neq 0,
\qquad
T_\delta\neq 0,
\qquad
S_\delta\neq 0,
\qquad
U_\delta\neq 0,
\qquad
W_\delta=0.
\label{eq:app-structured-aligned-seed-which-blocks}
\end{equation}
Its decisive mass-ratio degree pattern is equally simple.  The induced COM
interaction correction has degree $4$ in $\nu$ in each of the nonvanishing
blocks,
\begin{equation}
Q_\delta:\deg_\nu=4,
\qquad
T_\delta:\deg_\nu=4,
\qquad
S_\delta:\deg_\nu=4,
\qquad
U_\delta:\deg_\nu=4,
\qquad
W_\delta=0,
\label{eq:app-structured-aligned-seed-degree-pattern}
\end{equation}
while the pure-kinetic COM correction remains separate and plays no role in the
local interaction lift.

This is the precise sense in which the remaining local ordinary issue was never
another comparable-mass residual-lift problem.  It was the additional ordinary
chart correction induced by aligning the local seed with the Hamiltonian-first
translation.

\subsection{Why the old ordinary scaffold failed}
\label{app:structured-aligned-seed-old-scaffold-failure}

The ordinary comparable-mass interaction scaffold inherited from the earlier
local span audit consists of the constant-coefficient exchange-symmetric block
families
\begin{equation}
\mathcal Q,
\qquad
\mathcal T,
\qquad
\mathcal S,
\qquad
\mathcal U,
\qquad
\mathcal W,
\label{eq:app-structured-aligned-seed-old-families}
\end{equation}
with mass degrees $0,1,2,3,4$ respectively in the invariant variables
$(a,b,c,d,e,p,q)$.  That scaffold is sufficient for the Hamiltonian-chart
comparable-mass lift and, after Appendix~\ref{app:ordinary-translation-aligned-seed},
also for the translated canonical ordinary residual.  It does \emph{not},
however, lift the seed-alignment correction.

The exact blockwise verdict is
\begin{equation}
Q_\delta\in \mathcal Q,
\qquad
W_\delta=0,
\label{eq:app-structured-aligned-seed-old-success-blocks}
\end{equation}
but
\begin{equation}
T_\delta\notin \mathrm{Im}(\mathcal T),
\qquad
S_\delta\notin \mathrm{Im}(\mathcal S),
\qquad
U_\delta\notin \mathrm{Im}(\mathcal U).
\label{eq:app-structured-aligned-seed-old-failure-blocks}
\end{equation}
So the earlier ordinary-chart obstruction was genuine, but it was already much
narrower than a new local existence problem: it sat entirely in the seed sector.

\subsection{Exact structured seed-sector lift theorem}
\label{app:structured-aligned-seed-lift-theorem}

The exact cure is to enlarge only the seed spaces, and only in a very specific
way.

\begin{theorem}[Structured generic-frame lift theorem for the aligned-seed correction]
\label{thm:app-structured-aligned-seed-lift}
Let $\delta L_{4,\mathrm{seed}}^{\mathrm{loc}}$ be the local ordinary
seed-alignment correction defined by
\eqref{eq:app-structured-aligned-seed-delta-seed-def}.  Then the correction is
exactly generic-frame liftable in the structured seed spaces
\begin{equation}
Q_\delta\in \mathcal Q,
\qquad
T_\delta\in \mathcal T\oplus(pq)\,\mathcal T,
\qquad
S_\delta\in \mathcal S\oplus(pq)\,\mathcal S,
\label{eq:app-structured-aligned-seed-lift-theorem-a}
\end{equation}
\begin{equation}
U_\delta\in \mathcal U\oplus(p^2q^2)\,\mathcal U,
\qquad
W_\delta=0.
\label{eq:app-structured-aligned-seed-lift-theorem-b}
\end{equation}
Moreover, the enlarged blockwise coefficient-space ranks are full:
\begin{equation}
Q:\mathrm{rank}=20/20,
\qquad
T:\mathrm{rank}=16/16,
\qquad
S:\mathrm{rank}=12/12,
\qquad
U:\mathrm{rank}=8/8.
\label{eq:app-structured-aligned-seed-full-ranks}
\end{equation}
\end{theorem}

\begin{proof}
The starting point is the failure statement
\eqref{eq:app-structured-aligned-seed-old-failure-blocks}.  The ordinary
seed-alignment correction has the same exchange symmetry and strict one-body
separation properties as the earlier local scaffold, so the correct repair must
preserve the block families rather than replacing them by unrelated new
monomials.  The exact audit shows that adjoining one $pq$-dressed copy of the
old $\mathcal T$ family and one $pq$-dressed copy of the old $\mathcal S$
family, together with one $(p^2q^2)$-dressed copy of the old $\mathcal U$
family, is already sufficient to make the image ranks maximal.  No enlargement
is required in $\mathcal Q$, and the top block vanishes identically.  Hence the
correction is exactly liftable in the structured seed spaces listed in
\eqref{eq:app-structured-aligned-seed-lift-theorem-a}--
\eqref{eq:app-structured-aligned-seed-lift-theorem-b}, with the full-rank
statement of \eqref{eq:app-structured-aligned-seed-full-ranks} following
blockwise.
\end{proof}

The conceptual consequence is the important one.  The seed-alignment correction
is not an inexplicable extra obstruction.  It is itself an exact generic-frame
object, but only after the seed sector is dressed in the structured way above.
This is why the ordinary-chart obstruction is sharp, narrow, and local rather
than diffuse.

\subsection{Explicit canonical representative of the aligned-seed correction}
\label{app:structured-aligned-seed-explicit}

Using the exact Gauss--Jordan lift with all null coordinates set to zero yields
one sparse canonical representative of the aligned-seed correction.  In the
structured spaces of Theorem~\ref{thm:app-structured-aligned-seed-lift}, the
nonzero-direction counts are
\begin{equation}
Q_\delta:13,
\qquad
T_\delta:14,
\qquad
S_\delta:12,
\qquad
U_\delta:8,
\qquad
W_\delta:0.
\label{eq:app-structured-aligned-seed-nonzero-counts}
\end{equation}
One exact canonical representative is the following.

\paragraph{$G/r$ block.}
\begin{equation}
\begin{aligned}
Q_\delta={}&-\frac{1}{32}\Big(
214 a^3 c - 698 a^2 b^2 - 476 a^2 b c - 122 a^2 b d^2 - 290 a^2 b d e - 181 a^2 b e^2
+ 16 a^2 c^2 - 8 a^2 c d^2 - 18 a^2 d^2 e^2 - 54 a^2 d e^3 - 27 a^2 e^4
\\
&\hspace{2em}- 476 a b^2 c - 181 a b^2 d^2 - 290 a b^2 d e - 122 a b^2 e^2
+ 30 a d^3 e^3 + 15 a d^2 e^4
+ 214 b^3 c + 16 b^2 c^2 - 8 b^2 c e^2 - 27 b^2 d^4 - 54 b^2 d^3 e - 18 b^2 d^2 e^2
\\
&\hspace{2em}+ 15 b d^4 e^2 + 30 b d^3 e^3
\Big).
\end{aligned}
\label{eq:app-structured-aligned-seed-Qdelta}
\end{equation}

\paragraph{$G^2/r^2$ block.}
\begin{equation}
\begin{aligned}
T_\delta={}&-\frac{1}{96}\Big(
2814 a^2 b p^2 q + 1065 a^2 b p - 4725 a^2 b q + 2256 a^2 c q
- 3495 a^2 d^2 p^2 q + 408 a^2 d^2 p + 8031 a^2 d e p^2 q - 1782 a^2 d e p
\\
&\hspace{2em}+ 2814 a b^2 p q^2 - 4725 a b^2 p + 1065 a b^2 q
- 80 a d^3 e q + 2736 a d^2 e^2 p^2 q - 592 a d^2 e^2 p + 80 a d^2 e^2 q
+ 2256 b^2 c p
\\
&\hspace{2em}+ 8031 b^2 d e p q^2 - 1782 b^2 d e q
- 3495 b^2 e^2 p q^2 + 408 b^2 e^2 q
+ 2736 b d^2 e^2 p q^2 + 80 b d^2 e^2 p - 592 b d^2 e^2 q - 80 b d e^3 p
\\
&\hspace{2em}+ 432 d^3 e^3 p^2 q + 432 d^3 e^3 p q^2 + 576 d^3 e^3 p + 576 d^3 e^3 q
\Big).
\end{aligned}
\label{eq:app-structured-aligned-seed-Tdelta}
\end{equation}

\paragraph{$G^3/r^3$ block.}
\begin{equation}
\begin{aligned}
S_\delta={}&\frac{1}{192}\Big(
3672 a^2 p^3 q + 4464 a^2 p^2 q^2 - (10660+3\pi^2)a^2 p^2 + (1220-3\pi^2)a^2 p q
- 11304 a d^2 p^3 q - 7464 a d^2 p^2 q^2
\\
&\hspace{2em}- (2460-9\pi^2)a d^2 p^2 + (9\pi^2-2292)a d^2 p q
+ 4464 b^2 p^2 q^2 + 3672 b^2 p q^3 + (1220-3\pi^2)b^2 p q - (10660+3\pi^2)b^2 q^2
\\
&\hspace{2em}- 7464 b e^2 p^2 q^2 - 11304 b e^2 p q^3 + (9\pi^2-2292)b e^2 p q - (2460-9\pi^2)b e^2 q^2
+ 960 d^3 e q^2
\\
&\hspace{2em}- 11848 d^2 e^2 p^3 q - 19136 d^2 e^2 p^2 q^2 - 8512 d^2 e^2 p^2
- 11848 d^2 e^2 p q^3 - 8512 d^2 e^2 q^2
+ 960 d e^3 p^2
\Big).
\end{aligned}
\label{eq:app-structured-aligned-seed-Sdelta}
\end{equation}

\paragraph{$G^4/r^4$ block.}
\begin{equation}
\begin{aligned}
U_\delta={}&-\frac{pq}{96}\Big(
372 a p^3 q^2 + 108 a p^2 q^3 + (30\pi^2-1799)a p + (9\pi^2-1200)a q
+ 108 b p^3 q^2 + 372 b p^2 q^3 + (9\pi^2-1200)b p + (30\pi^2-1799)b q
\\
&\hspace{2em}+ 7740 d^2 p^3 q^2 + 2340 d^2 p^2 q^3 - (6953+96\pi^2)d^2 p - (2688+27\pi^2)d^2 q
+ 2340 e^2 p^3 q^2 + 7740 e^2 p^2 q^3 - (2688+27\pi^2)e^2 p - (6953+96\pi^2)e^2 q
\Big).
\end{aligned}
\label{eq:app-structured-aligned-seed-Udelta}
\end{equation}

\paragraph{$G^5/r^5$ block.}
\begin{equation}
W_\delta=0.
\label{eq:app-structured-aligned-seed-Wdelta}
\end{equation}

These formulas are not meant to be unique.  They are one exact canonical choice
inside the enlarged seed spaces.  What matters for the theorem chain is that a
representative exists and that its COM reduction reproduces the exact aligned
seed correction slot by slot.

\subsection{Natural seed, aligned seed, and one full local generic-frame candidate}
\label{app:structured-aligned-seed-full-candidate}

The natural generic-frame local ordinary seed is the direct symmetric promotion
of the exact one-body $4$PN gate:
\begin{equation}
Q_{\mathrm{nat}}=\frac{75}{128}(a^4+b^4),
\qquad
T_{\mathrm{nat}}=\frac{59}{16}(qa^3+pb^3),
\label{eq:app-structured-aligned-seed-natural-seed-a}
\end{equation}
\begin{equation}
S_{\mathrm{nat}}=\frac{203}{32}(q^2a^2+p^2b^2),
\qquad
U_{\mathrm{nat}}=\frac{31}{32}(q^3a+p^3b),
\qquad
W_{\mathrm{nat}}=\frac{1}{16}(q^4+p^4).
\label{eq:app-structured-aligned-seed-natural-seed-b}
\end{equation}
Thus the aligned local seed is simply
\begin{equation}
L_{4,\mathrm{seed}}^{\mathrm{aligned,loc}}
=
L_{4,\mathrm{seed}}^{\mathrm{natural,loc}}
+
\delta L_{4,\mathrm{seed}}^{\mathrm{loc}},
\label{eq:app-structured-aligned-seed-aligned-seed}
\end{equation}
with $\delta L_{4,\mathrm{seed}}^{\mathrm{loc}}$ built from
\eqref{eq:app-structured-aligned-seed-block-decomposition} and
\eqref{eq:app-structured-aligned-seed-Qdelta}--
\eqref{eq:app-structured-aligned-seed-Wdelta}.

Combining that aligned seed with the already solved canonical ordinary residual
of Appendix~\ref{app:ordinary-translation-aligned-seed} yields one exact local
ordinary generic-frame candidate,
\begin{equation}
\boxed{
L_{4,\mathrm{loc}}^{(\mathrm{gen})}
=
L_{4,\mathrm{seed}}^{(\mathrm{natural,gen})}
+
\delta L_{4,\mathrm{seed}}^{(\mathrm{gen})}
+
\Delta L_{4,\mathrm{loc}}^{(\mathrm{can,gen})}.
}
\label{eq:app-structured-aligned-seed-full-local-candidate}
\end{equation}
Its COM reduction reproduces the full fixed-chart local ordinary $4$PN target
block by block.  Therefore the local-first $4$PN program no longer has a local
ordinary existence bottleneck.

More precisely, the logical chain is now complete:
\begin{enumerate}
\item the Hamiltonian-chart comparable-mass residual lift was solved in
Appendices~\ref{app:hamiltonian-lift} and \ref{app:canonical-slice};
\item the translated canonical ordinary residual was fixed exactly in
Appendix~\ref{app:ordinary-translation-aligned-seed};
\item the remaining seed-alignment correction is solved by the structured
seed-sector lift of the present appendix; and
\item the sum \eqref{eq:app-structured-aligned-seed-full-local-candidate}
reproduces the full fixed-chart local ordinary target.
\end{enumerate}
So the only genuinely open local question left is no longer existence, but the
sharper uniqueness/interpretation question for the aligned-seed sector itself.

% Appendix: Exact local referee reconstruction

\section{Exact local referee reconstruction}
\label{app:local-referee-reconstruction}

The previous appendices already froze all of the ingredients needed for the
local $4$PN theorem chain, but they did so in the logically correct order rather
than as one end-to-end replay.  Appendix~\ref{app:imported-local-target-hamiltonian-first}
imported the fixed-chart local ADM Hamiltonian target and showed that the local
problem must be handled \emph{Hamiltonian-first}.  Appendix~\ref{app:hamiltonian-lift}
then solved the generic-frame Hamiltonian comparable-mass lift, and
Appendix~\ref{app:canonical-slice} froze one exact canonical representative of
that Hamiltonian residual.  Appendix~\ref{app:ordinary-translation-aligned-seed}
showed that the translated canonical ordinary residual obeys the exact quartic
sign-flip rule, while Appendix~\ref{app:structured-aligned-seed-lift} proved
that the remaining ordinary-chart seed-alignment correction is itself exactly
liftable in a minimal structured enlargement of the seed sector.

The purpose of the present appendix is to assemble those already-frozen pieces
into one exact referee reconstruction of the \emph{entire local instantaneous}
ordinary $4$PN sector.  The conclusion is sharper than a consistency check.
There is now one explicit generic-frame ordinary representative whose center-of-mass
reduction reproduces the full fixed-chart local ordinary target slot by slot.
So the local $4$PN program is no longer blocked by local existence.  What remains
open after the present appendix is the separate aligned-seed uniqueness /
interpretation question and the tail / hereditary bridge.

\subsection{Exact input data reused by the referee reconstruction}
\label{app:local-referee-reconstruction-inputs}

The referee reconstruction uses four exact inputs and nothing else.

First, let
\begin{equation}
L_{4,\mathrm{loc}}^{\mathrm{ord,target}}
\label{eq:app-local-referee-target-def}
\end{equation}
be the fixed-chart local ordinary $4$PN target obtained from the imported local
ADM Hamiltonian target by the exact quartic Hamiltonian-to-ordinary compiler of
Appendix~\ref{app:quartic-compiler}.  By construction, this is the object whose
local COM slots must be reproduced.

Second, write the natural and Hamiltonian-aligned local ordinary seeds as
\begin{equation}
L_{4,\mathrm{seed}}^{(\mathrm{natural})},
\qquad
L_{4,\mathrm{seed}}^{(\mathrm{aligned})},
\label{eq:app-local-referee-natural-aligned-seeds}
\end{equation}
and define their difference by
\begin{equation}
\delta L_{4,\mathrm{seed}}
=
L_{4,\mathrm{seed}}^{(\mathrm{aligned})}
-
L_{4,\mathrm{seed}}^{(\mathrm{natural})}.
\label{eq:app-local-referee-delta-seed}
\end{equation}
The aligned seed already fixes the entire free block of the translated local
ordinary target: the unique nonzero free slot matches exactly, while the
remaining free slots vanish identically.  So the free local sector does not
require any additional generic-frame repair.

Third, let
\begin{equation}
\Delta H_{4,\mathrm{loc}}^{\mathrm{can}}
\label{eq:app-local-referee-hcan}
\end{equation}
be the canonical generic-frame Hamiltonian comparable-mass residual frozen in
Appendix~\ref{app:canonical-slice}.  Appendix~\ref{app:ordinary-translation-aligned-seed}
then gives the exact translated ordinary residual
\begin{equation}
\Delta L_{4,\mathrm{loc}}^{\mathrm{can}}
=
-\Delta H_{4,\mathrm{loc}}^{\mathrm{can}}.
\label{eq:app-local-referee-sign-flip}
\end{equation}
So the true comparable-mass local residual is already settled in the ordinary
chart: there is no remaining comparable-mass lift obstruction after the
Hamiltonian-first step.

Fourth, Appendix~\ref{app:structured-aligned-seed-lift} proved that the local
seed-alignment correction \eqref{eq:app-local-referee-delta-seed} is exactly
liftable in the structured seed spaces
\begin{equation}
Q,
\qquad
T\oplus(pq)T,
\qquad
S\oplus(pq)S,
\qquad
U\oplus(p^2q^2)U,
\qquad
W=0.
\label{eq:app-local-referee-structured-spaces}
\end{equation}
Write one such exact generic-frame lift as
\begin{equation}
\delta L_{4,\mathrm{seed}}^{(\mathrm{gen})}.
\label{eq:app-local-referee-delta-seed-gen}
\end{equation}
This is the only place where the old ordinary-chart obstruction ever lived.

\subsection{The referee reconstruction formula}
\label{app:local-referee-reconstruction-formula}

With the exact pieces above in hand, define the generic-frame local ordinary
candidate by
\begin{equation}
L_{4,\mathrm{loc}}^{(\mathrm{gen})}
=
L_{4,\mathrm{seed}}^{(\mathrm{natural,gen})}
+
\delta L_{4,\mathrm{seed}}^{(\mathrm{gen})}
+
\Delta L_{4,\mathrm{loc}}^{(\mathrm{can,gen})}.
\label{eq:app-local-referee-master-assembly}
\end{equation}
This is the exact local referee reconstruction formula.

Its interpretation is transparent.

\begin{enumerate}
\item
The natural generic-frame seed carries the exact one-body / self / static local
$4$PN content in the ordinary chart.
\item
The lifted seed-alignment correction repairs the fact that the ordinary target
must be aligned with the Hamiltonian-first translation rather than with the raw
natural seed.
\item
The translated canonical residual inserts the already-solved comparable-mass
local content.
\end{enumerate}

Because the three pieces are exact and because the quartic compiler is linear in
the $4$PN slot once the lower-order ledger is frozen, the only remaining question
is direct slot-by-slot verification after COM reduction.

\subsection{Exact local referee reconstruction theorem}
\label{app:local-referee-reconstruction-theorem}

The local ordinary target has a six-slot free block together with fifteen
interaction slots distributed by block as
\begin{equation}
Q:5,
\qquad
T:4,
\qquad
S:3,
\qquad
U:2,
\qquad
W:1.
\label{eq:app-local-referee-slot-counts}
\end{equation}
The referee reconstruction theorem is therefore most cleanly stated directly in
that slot language.

\begin{theorem}[Exact local referee reconstruction]
\label{thm:app-local-referee-reconstruction}
Let $L_{4,\mathrm{loc}}^{\mathrm{ord,target}}$ be the fixed-chart local ordinary
$4$PN target defined by \eqref{eq:app-local-referee-target-def}, and let
$L_{4,\mathrm{loc}}^{(\mathrm{gen})}$ be the generic-frame candidate assembled in
\eqref{eq:app-local-referee-master-assembly}.  Then the center-of-mass reduction
of $L_{4,\mathrm{loc}}^{(\mathrm{gen})}$ reproduces the full local ordinary target
exactly:
\begin{equation}
\bigl(L_{4,\mathrm{loc}}^{(\mathrm{gen})}\bigr)_{\mathrm{COM}}
=
L_{4,\mathrm{loc}}^{\mathrm{ord,target}}.
\label{eq:app-local-referee-reconstruction-identity}
\end{equation}
Equivalently,
\begin{equation}
\Delta Q=0,
\qquad
\Delta T=0,
\qquad
\Delta S=0,
\qquad
\Delta U=0,
\qquad
\Delta W=0,
\label{eq:app-local-referee-blockwise-zero}
\end{equation}
where each difference denotes the induced COM blockwise slot mismatch between
\eqref{eq:app-local-referee-master-assembly} and the fixed-chart local ordinary
target.  In particular, the entire local instantaneous ordinary $4$PN sector
admits an exact generic-frame representative inside the declared hierarchy.
\end{theorem}

\begin{proof}
The free block is already fixed by the aligned seed, so there is nothing further
to solve there.  For the interaction blocks, split the target as in
Appendix~\ref{app:ordinary-translation-aligned-seed} into natural seed,
seed-alignment correction, and canonical comparable-mass residual.  The natural
seed is carried unchanged by
$L_{4,\mathrm{seed}}^{(\mathrm{natural,gen})}$.  The seed-alignment correction is
carried exactly by $\delta L_{4,\mathrm{seed}}^{(\mathrm{gen})}$ because of the
structured-lift theorem of
Appendix~\ref{app:structured-aligned-seed-lift}.  The remaining
comparable-mass local residual is carried exactly by
$\Delta L_{4,\mathrm{loc}}^{(\mathrm{can,gen})}$ because the Hamiltonian-chart
canonical representative exists by Appendix~\ref{app:canonical-slice} and its
ordinary translation is fixed by the exact sign-flip rule
\eqref{eq:app-local-referee-sign-flip}.  By linearity of the COM slot map, the
sum of the three reconstructed pieces therefore reproduces the COM image of the
full target in each block.  This gives
\eqref{eq:app-local-referee-reconstruction-identity} and hence the blockwise
vanishing statement \eqref{eq:app-local-referee-blockwise-zero}.
\end{proof}

The theorem is stronger than a mere existence statement in coefficient space.
The exact referee audit checks the blockwise COM image directly.  The local
ordinary $4$PN target is reproduced in all five interaction blocks with the
exact slot counts of \eqref{eq:app-local-referee-slot-counts}; no unresolved
local coefficient remains after the assembly.

\subsection{Blockwise replay of the exact referee checks}
\label{app:local-referee-reconstruction-blockwise}

The end-to-end audit may be read as the following sequence of exact replay
checks.

\paragraph{Free block.}
The aligned ordinary seed reproduces the entire free local block: the unique
nonzero free coefficient matches the translated target exactly, and the
remaining five free slots are identically zero.  So the free ordinary local
sector is already closed before any interaction lift is attempted.

\paragraph{$Q$ block.}
The natural seed supplies the exact one-body $G/r$ ordinary content, the
structured seed lift contributes the required aligned-seed correction already
inside $Q$, and the translated canonical residual supplies the remaining
comparable-mass $G/r$ content.  The referee replay verifies that all five COM
slots of the reconstructed $Q$ block match the fixed-chart target exactly.

\paragraph{$T$ block.}
The natural seed supplies the one-body $G^2/r^2$ content.  The seed-alignment
correction is not contained in the old constant-coefficient $T$ family, but it
is exactly liftable in the enlarged space $T\oplus(pq)T$.
After inserting that lifted correction and the translated canonical residual,
all four COM $T$ slots match exactly.

\paragraph{$S$ block.}
The same pattern repeats in the $G^3/r^3$ sector.  The natural seed carries the
exact one-body block, the aligned-seed correction is lifted in
$S\oplus(pq)S$, and the translated canonical residual completes the
comparable-mass part.  The resulting $S$ block matches the three target COM
slots exactly.

\paragraph{$U$ block.}
At $G^4/r^4$ the seed-alignment correction is lifted in the minimal enlarged
space $U\oplus(p^2q^2)U$.  Once that correction and the translated canonical
residual are added to the natural seed, both COM $U$ slots match the fixed-chart
target exactly.

\paragraph{$W$ block.}
The aligned-seed correction vanishes identically in the top static block, so no
additional lift is required there.  The natural static seed together with the
translated canonical residual already reproduces the unique $W$-slot exactly.

So the referee reconstruction is genuinely end to end: every local interaction
block is checked directly after COM reduction, and every block passes.

\subsection{Canonical sparsity of the lifted representative}
\label{app:local-referee-reconstruction-sparsity}

The end-to-end theorem above is an existence statement for the full local
ordinary sector, but the referee archive also records one concrete sparse
realization of the lifted seed-alignment correction.  In the structured spaces
\eqref{eq:app-local-referee-structured-spaces}, the exact Gauss--Jordan slice
with all null coordinates set to zero yields the nonzero-direction counts
\begin{equation}
Q:13,
\qquad
T:14,
\qquad
S:12,
\qquad
U:8,
\qquad
W:0,
\label{eq:app-local-referee-sparse-counts}
\end{equation}
inside the basis sizes
\begin{equation}
Q:52,
\qquad
T:92,
\qquad
S:58,
\qquad
U:20,
\qquad
W:0.
\label{eq:app-local-referee-basis-sizes}
\end{equation}
These numbers do not by themselves prove uniqueness of the aligned-seed sector,
but they do show that the local theorem is not merely formal.  One explicit,
sparse, exact representative exists and has been checked directly against the
fixed-chart target.

\subsection{What this appendix settles and what it leaves open}
\label{app:local-referee-reconstruction-status}

The result of the present appendix should be read carefully.

It \emph{does} settle the local $4$PN existence problem in the ordinary chart.
The exact imported local target is reconstructed from the Hamiltonian-first
compiler chain; the generic-frame comparable-mass local residual is already
solved in the Hamiltonian chart and translates back by exact sign flip; the
remaining aligned-seed correction is exactly liftable in the minimal structured
seed spaces; and the combined generic-frame candidate reproduces the full local
ordinary target after COM reduction.

It \emph{does not} settle two sharper questions.
First, it does not yet provide a final uniqueness / interpretation theorem for
the aligned-seed sector beyond the exact sparse slice used by the referee
archive.  Second, it does not touch the separate hereditary bridge.  That bridge
remains isolated by construction and is controlled by the same passive / outgoing
STF quadrupole normalization that already survived the $2.5$PN analysis.

Accordingly, the correct post-referee status of the local chain is the compact
statement
\begin{equation}
\boxed{
L_{4,\mathrm{loc}}^{\mathrm{ord,target}}
\text{ is already reconstructed exactly in the generic frame.}}
\label{eq:app-local-referee-final-status}
\end{equation}
From this point onward, the mathematically clean open problem is no longer local
existence.  It is the separate tail / normalization interface.

% Appendix: Hereditary/tail bridge and coefficient relation

\section{Hereditary/tail bridge and coefficient relation}
\label{app:tail-hereditary-bridge-coefficient}

This appendix records the detailed coefficient arithmetic behind the hereditary
$4$PN bridge.  The local instantaneous $4$PN sector is already fixed by the
local referee chain, so the only remaining nonlocal issue is the coefficient of
the time-symmetric tail functional.  The point of the appendix is threefold.
First, we freeze the standard GR hereditary structure and its local logarithmic
shadow.  Second, we show that the hereditary coefficient is tied exactly to the
same canonically normalized STF quadrupole coefficient that already controls
our carried $2.5$PN Burke--Thorne channel.  Third, we isolate the only
remaining model-side ambiguity as a single scalar transport factor rather than a
new tensor channel, a new generic-frame existence problem, or a new independent
normalization datum.

\subsection{GR reference structure}
\label{app:tail-gr-reference-structure}

The conservative $4$PN dynamics in GR splits into a local instantaneous sector
plus a time-symmetric hereditary sector,
\begin{equation}
L_{4\mathrm{PN}}^{\mathrm{cons}}
=
L_{4\mathrm{PN}}^{\mathrm{local}}
+
L_{4\mathrm{PN}}^{\mathrm{tail}}.
\label{eq:app-tail-gr-split}
\end{equation}
On the action side the hereditary term may be written as
\begin{equation}
S_{\mathrm{tail}}
=
\frac{G^2 M}{5c^8}
\operatorname{Pf}_{2s/c}
\iint \frac{dt\,dt'}{|t-t'|}
I_{ij}^{(3)}(t)
I_{ij}^{(3)}(t'),
\label{eq:app-tail-action}
\end{equation}
while the corresponding time-symmetric nonlocal Hamiltonian term is
\begin{equation}
H_{\mathrm{tail}}(t)
=
-
\frac{G^2 M}{5c^8}
I_{ij}^{(3)}(t)
\operatorname{Pf}_{2s/c}
\int_{-\infty}^{+\infty}
\frac{dv}{|v|}
I_{ij}^{(3)}(t+v).
\label{eq:app-tail-hamiltonian}
\end{equation}
So the universal GR hereditary coefficient is
\begin{equation}
C_{\mathrm{tail}}^{\mathrm{GR}}
=
\frac{G^2 M}{5c^8}.
\label{eq:app-tail-gr-coefficient}
\end{equation}
At the level relevant for the local logarithmic shadow, the same hereditary
sector contributes the local functional
\begin{equation}
F_{\mathrm{tail}}(t)
=
\frac{2}{5}\frac{G^2 M}{c^8}
\bigl(I_{ij}^{(3)}(t)\bigr)^2.
\label{eq:app-tail-local-functional}
\end{equation}
This is already enough to see that the hereditary side is structurally much
smaller than the local lift: it is one fixed STF-quadrupole functional with one
scalar coefficient.

\subsection{Universal Newtonian quadrupole shadow}
\label{app:tail-newtonian-shadow}

Because the hereditary coefficient first enters at $4$PN, its local logarithmic
shadow is fixed by the Newtonian order-reduced STF quadrupole,
\begin{equation}
I_{ij}
=
\mu\,\mathrm{STF}(x_i x_j),
\qquad
M\equiv m_A+m_B,
\qquad
\mu\equiv \frac{m_A m_B}{M},
\qquad
\nu\equiv \frac{\mu}{M}.
\label{eq:app-tail-newtonian-quadrupole}
\end{equation}
The exact third derivative is
\begin{equation}
I_{ij}^{(3)}
=
-
\frac{2GM\mu}{r^3}
\left(
4x_{\langle i}v_{j\rangle}
-
3\frac{\dot r}{r}x_{\langle i}x_{j\rangle}
\right),
\label{eq:app-tail-third-derivative}
\end{equation}
so that
\begin{equation}
\bigl(I_{ij}^{(3)}\bigr)^2
=
\frac{8}{3}\frac{G^2M^2\mu^2}{r^4}
\left(12v^2-11\dot r^2\right).
\label{eq:app-tail-third-derivative-square}
\end{equation}
Substituting \eqref{eq:app-tail-third-derivative-square} into
\eqref{eq:app-tail-local-functional} gives the exact local logarithmic tail
shadow
\begin{equation}
\boxed{
\frac{F_{\mathrm{tail}}}{\mu}
=
\frac{16}{15}\,\nu\,\frac{U^4}{c^8}
\left(12v^2-11\dot r^2\right),
\qquad
U\equiv \frac{GM}{r}.
}
\label{eq:app-tail-local-shadow}
\end{equation}
Equation~\eqref{eq:app-tail-local-shadow} shows that the hereditary shadow lives
only in the degree-$2$ local $G^4/r^4$ block.  That is the precise reason the
tail problem is much smaller than the local generic-frame lift problem.

\subsection{Exact bridge to the carried \texorpdfstring{$2.5$}{2.5}PN quadrupole coefficient}
\label{app:tail-exact-bridge}

The $2.5$PN program already isolated the correct canonical language for the
surviving radiative branch.  Let
\begin{equation}
\gamma_{\rm quad}^{\rm eff}
=
\mathcal N_Q\,\frac{a_{\rm th}^5}{27c_s^5}
=
\overline\Gamma_5
=
9\frac{\overline K_2^{5/2}}{\overline K_0^{3/2}}
\label{eq:app-tail-carried-gamma}
\end{equation}
be the canonically normalized STF quadrupole coefficient on that branch.  In
GR the corresponding Burke--Thorne coefficient is
\begin{equation}
\gamma_{\rm GR}
=
\frac{2G}{5c^5}.
\label{eq:app-tail-burke-thorne}
\end{equation}
The key arithmetic observation is that the GR hereditary coefficient is related
exactly to the Burke--Thorne coefficient by
\begin{equation}
\boxed{
C_{\mathrm{tail}}^{\mathrm{GR}}
=
\frac{GM}{2c^3}\,\gamma_{\rm GR}.
}
\label{eq:app-tail-gr-bridge}
\end{equation}
Indeed,
\[
\frac{GM}{2c^3}\,\frac{2G}{5c^5}
=
\frac{G^2M}{5c^8}
=
C_{\mathrm{tail}}^{\mathrm{GR}}.
\]
So the conservative hereditary coefficient is not merely analogous to the
carried $2.5$PN odd quadrupole coefficient.  It is fixed by the same canonical
quadrupole coefficient through an exact bridge.

This upgrades immediately to the carried toy-model branch.  The unique
compatible conservative hereditary coefficient is
\begin{equation}
\boxed{
C_{\mathrm{tail}}
=
\frac{GM}{2c^3}\,\gamma_{\rm quad}^{\rm eff}.
}
\label{eq:app-tail-exact-coefficient-bridge}
\end{equation}
Equation~\eqref{eq:app-tail-exact-coefficient-bridge} is the central
coefficient relation of the appendix.

\subsection{Toy-model branch forms and the scalar transport factor}
\label{app:tail-toy-branch-form}

It is useful to isolate explicitly the only model-side ambiguity that can still
survive at this stage.  Let
\begin{equation}
\gamma_{\rm toy}
\equiv
\hat m_0^{\,2}\Gamma_5
\label{eq:app-tail-gamma-toy}
\end{equation}
be the canonically normalized STF quadrupole coefficient inherited from the
$2.5$PN branch analysis.  Then one may parameterize the reduced hereditary
transport by a single scalar factor $\Theta_{\rm tail}$ via
\begin{equation}
\boxed{
\alpha_{\rm tail}^{\rm toy}
=
\Theta_{\rm tail}
\frac{GM}{2c_s^3}
\gamma_{\rm toy}.
}
\label{eq:app-tail-theta-transport}
\end{equation}
On the canonical branch with
\begin{equation}
c_s=c,
\qquad
\gamma_{\rm toy}=\gamma_{\rm quad}^{\rm eff},
\label{eq:app-tail-canonical-branch}
\end{equation}
Eq.~\eqref{eq:app-tail-theta-transport} reduces to
\eqref{eq:app-tail-exact-coefficient-bridge} when
\begin{equation}
\Theta_{\rm tail}=1.
\label{eq:app-tail-theta-one}
\end{equation}
So the only possible remaining mismatch on the hereditary side is a single
scalar monopole-scattering transport coefficient.  No new tensor channel is
needed, and no new generic-frame existence solve is needed.

Using \eqref{eq:app-tail-carried-gamma}, the same bridge may be written in the
compact toy-model form
\begin{equation}
C_{\mathrm{tail}}^{\mathrm{toy}}
=
\frac{GM}{2c^3}\,\gamma_{\rm quad}^{\rm eff}
=
\frac{GM}{2c^3}\,\overline\Gamma_5
=
\frac{9GM}{2c^3}
\frac{\overline K_2^{5/2}}{\overline K_0^{3/2}}.
\label{eq:app-tail-toy-final-form}
\end{equation}
Thus the entire hereditary $4$PN bridge is controlled by the same normalized
quadrupole branch data as the $2.5$PN odd channel.

\subsection{Why no new independent \texorpdfstring{$4$}{4}PN normalization datum opens}
\label{app:tail-no-new-normalization-datum}

The appendix may now be summarized as the following coefficient theorem inside
the declared hierarchy.

\paragraph{Coefficient theorem.}
If the canonically normalized orbital/worldtube STF quadrupole branch satisfies
the passive/outgoing low-frequency relations already isolated by the $2.5$PN
audit, then
\begin{enumerate}
  \item the unique compatible conservative hereditary coefficient is
  \eqref{eq:app-tail-exact-coefficient-bridge},
  \item the hereditary $4$PN sector introduces no new independent quadrupole
  normalization datum, and
  \item full GR matching of the hereditary coefficient is equivalent to the
  already-known $2.5$PN normalization target
  \begin{equation}
  \gamma_{\rm quad}^{\rm eff}
  =
  \frac{2G}{5c^5}.
  \label{eq:app-tail-gr-recovery-condition}
  \end{equation}
\end{enumerate}
The moving-throat normalization language used later in the paper rewrites the
same condition as
\begin{equation}
\mathcal N_Q^{\rm target}
=
\frac{54Gc_s^5}{5a_{\rm th}^5c^5},
\label{eq:app-tail-normalization-target}
\end{equation}
and, on the source-map notation of the reduced PDE program, equivalently as
\begin{equation}
\hat m_0^{\,2}P_0
=
\frac{54Gc_s^5}{5a_{\rm th}^5c^5}.
\label{eq:app-tail-p0-target}
\end{equation}
Substituting \eqref{eq:app-tail-gr-recovery-condition} into
\eqref{eq:app-tail-exact-coefficient-bridge} gives
\begin{equation}
C_{\mathrm{tail}}
=
\frac{GM}{2c^3}\,\frac{2G}{5c^5}
=
\frac{G^2M}{5c^8}
=
C_{\mathrm{tail}}^{\mathrm{GR}}.
\label{eq:app-tail-gr-recovery}
\end{equation}
So once the already-isolated passive/outgoing STF quadrupole normalization is
fixed on the natural moving-throat branch, the $4$PN hereditary coefficient
closes automatically.

What remains open is therefore maximally narrow.  The present appendix does not
claim a first-principles moving-throat PDE derivation of the passive/outgoing
quadrupole normalization, nor the final insertion of the normalized tail module
into a single end-to-end referee/master script.  But it does show that the full
remaining $4$PN hereditary gap is exactly the same passive/outgoing quadrupole
normalization theorem already inherited from the $2.5$PN program.

% Appendix: Verification artifact ledger and reference run summaries

\section{Verification artifact ledger and reference run summaries}
\label{app:verification-artifact-ledger}

This appendix records the supplementary referee archive for the present $4$PN
paper.  Its role is narrower than the analytic derivation in the main text.  The
archive is not presented as a black-box generator of the theorem.  Instead, it
serves as an executable and symbolic backstop for the exact one-body gate, the
quartic compiler, the local Hamiltonian-first lift, the aligned-seed repair, the
local referee reconstruction, and the hereditary bridge formulas used in the
paper.  The archive therefore verifies the paper \emph{within the declared
closure hierarchy}; it does not replace any analytic step, and it does not by
itself upgrade the present work into an assumption-free theorem of a fully solved
moving-throat bulk PDE.

The verification logic at $4$PN is naturally split into three layers.  The
first layer is the inherited lower-order archive carried forward from the parent
$4$D, $1$PN, $2$PN, $2.5$PN, and $3$PN packages.  The second layer is the
dedicated dual-CAS $4$PN archive, whose job is to check the new theorem chain
specific to the present paper: the strict one-body gate, the quartic ordinary /
Hamiltonian compiler, the local scaffold and fixed-chart import, the
Hamiltonian-first generic-frame lift, the aligned-seed correction, the local
referee reconstruction, and the hereditary bridge.  The third layer is the
conditional end-to-end replay, which ties the already-frozen $2.5$PN quadrupole
narrowing to the exact local $4$PN closure and the hereditary bridge, and
thereby verifies that the only remaining $4$PN interface is the same
passive/outgoing STF quadrupole normalization already isolated on the $2.5$PN
side.

\subsection{Reading rule for the archive}
\label{app:verification-reading-rule}

The archive should be read with the same claim-status firewall used throughout
the paper.  Exact symbolic statements are checked as exact identities.  Reduced
and chart-level statements are checked inside the same Hamiltonian-first and
worldtube / small-body conventions declared in the main text.  Closure-level
statements are checked only inside that same hierarchy.  So the archive is not
being asked to prove a different theorem from the one stated in the paper.  It
is being asked to verify that the paper's own exact, reduced, and
closure-dependent claims close where they are claimed to close.

\subsection{Inherited reference masters}
\label{app:verification-inherited-masters}

The $4$PN archive is not standalone in the historical sense.  It sits on top of
four already-frozen referee layers.

\begin{enumerate}
  \item \textbf{Parent $4$D and full conservative $1$PN archive.}  This is the
  inherited lower-order freeze that carries the exact projection backbone,
  Newtonian mass-dressing coefficient $\kappa_\rho=1$, optical matching
  $n=5$, the added-inertia ledger $\kappa_{\rm add}=1/2$, the adiabatic
  $\kappa_{\rm PV}=3/2$ closure, the parity-even wake coefficients, and the
  exact conservative EIH equality.

  \item \textbf{Full conservative $2$PN archive.}  This is the inherited
  referee archive that freezes the one-body denominator closure, the self /
  static seed, the comparable-mass residual solve, and the exact equality to
  the standard generic-frame ADM Hamiltonian within the declared hierarchy.

  \item \textbf{$2.5$PN master audit.}  This is the inherited audit that
  narrowed the universal dissipative bridge to the canonically normalized
  orbital / worldtube STF quadrupole branch and reduced the remaining theorem
  gap to the passive/outgoing quadrupole normalization.

  \item \textbf{Full conservative $3$PN referee master.}  This is the carried
  archive that closes the grouped real $P_2$ middle block, the geometry-side
  static completion, and the pure-kinetic collapse beneath the same $2.5$PN
  outgoing-normalization bridge.
\end{enumerate}

These inherited archives matter because the present $4$PN package is not trying
to reopen lower-order viability conditions.  It is trying to prove that no new
independent normalization datum opens at $4$PN beyond the already-known
$2.5$PN quadrupole gate.

\subsection{Dedicated \texorpdfstring{$4$}{4}PN stage ledger}
\label{app:verification-4pn-stage-ledger}

The dedicated $4$PN publication bundle is organized around the following named
stage audits.  For each artifact we record both its theorem role and the short
reference summary against which reruns should be checked.

\begin{enumerate}
  \item \textbf{\texttt{4pn\_onebody\_audit.py}.}  \emph{Role:} exact
  Schwarzschild one-body gate through order $c^{-8}$.  \emph{Reference
  summary:} the strict one-body target forces the minimal repair ledger
  \[
  \mu_{\rho,4}=\frac18,
  \qquad
  d_4=\frac{205}{16},
  \qquad
  s_{34}=-\frac{15}{32},
  \qquad
  s_{26}=-\frac1{16}.
  \]

  \item \textbf{\texttt{4pn\_quartic\_legendre\_audit.py}.}  \emph{Role:}
  exact quartic perturbative Legendre-transform compiler.  \emph{Reference
  summary:} the audit verifies the exact coefficients $H_1$ through $H_4$ and
  freezes the quartic ordinary / Hamiltonian map used everywhere later in the
  local chain.

  \item \textbf{\texttt{4pn\_local\_scaffold\_audit.py}.}  \emph{Role:}
  raw local generic-frame interaction scaffold and COM completeness test.
  \emph{Reference summary:} the raw exchange-symmetric block sizes are
  \[
  52,\quad 46,\quad 29,\quad 10,\quad 2,
  \]
  for the $G/r$, $G^2/r^2$, $G^3/r^3$, $G^4/r^4$, and $G^5/r^5$ blocks,
  respectively, while COM image ranks are already maximal,
  \[
  5,\quad 4,\quad 3,\quad 2,\quad 1,
  \]
  with total COM-nullity $124$.  So there is no early local span shortage.

  \item \textbf{\texttt{4pn\_local\_target\_import\_audit.py}.}
  \emph{Role:} fixed-chart local ADM target import.  \emph{Reference
  summary:} the strict one-body Hamiltonian gate matches the imported local
  target exactly, the local COM Hamiltonian has one free slot plus interaction
  slot structure $5+4+3+2+1$, and the upper local Hamiltonian blocks stop at
  $\nu^2$ in the $G^4/r^4$ and $G^5/r^5$ sectors.

  \item \textbf{\texttt{4pn\_local\_hamiltonian\_to\_ordinary\_audit.py}.}
  \emph{Role:} exact reduced COM Hamiltonian-to-ordinary translation.
  \emph{Reference summary:} the quartic compiler translates the fixed local
  Hamiltonian target into the exact local ordinary target, but the ordinary
  chart develops extra $\nu^4$ tails in the $G^2/r^2$, $G^3/r^3$, and
  $G^4/r^4$ blocks.  This freezes the Hamiltonian-first lesson.

  \item \textbf{\texttt{4pn\_hamiltonian\_chart\_generic\_frame\_lift\_audit.py}.}
  \emph{Role:} Hamiltonian-chart generic-frame lift of the local comparable-mass
  residual.  \emph{Reference summary:} the exact blockwise coefficient-space
  ranks are maximal,
  \[
  20,\quad 12,\quad 9,\quad 4,\quad 2,
  \]
  so there is no Hamiltonian-side span obstruction.

  \item \textbf{\texttt{4pn\_hamiltonian\_chart\_canonical\_slice\_audit.py}.}
  \emph{Role:} canonical comparable-mass local residual slice.
  \emph{Reference summary:} the audit quotients the $92$-direction COM-blind
  null family and records one sparse canonical representative with blockwise
  nonzero counts
  \[
  Q:11,\qquad T:12,\qquad S:9,\qquad U:4,\qquad W:2,
  \]
  together with exact COM verification.

  \item \textbf{\texttt{4pn\_generic\_frame\_ordinary\_translation\_audit.py}.}
  \emph{Role:} translation of the canonical Hamiltonian residual back to the
  ordinary chart.  \emph{Reference summary:} the canonical comparable-mass
  local residual satisfies the exact sign-flip theorem
  \[
  \Delta H_4=-\Delta L_4,
  \]
  and the translated target is measured relative to the Hamiltonian-aligned
  ordinary seed.

  \item \textbf{\texttt{4pn\_generic\_frame\_aligned\_seed\_lift\_audit.py}.}
  \emph{Role:} generic-frame lift of the aligned-seed ordinary correction.
  \emph{Reference summary:} the old ordinary obstruction is shown to be a
  seed-sector problem only, and the aligned-seed correction lifts exactly in the
  structured spaces
  \[
  Q,
  \qquad
  T\oplus (pq)T,
  \qquad
  S\oplus (pq)S,
  \qquad
  U\oplus (p^2q^2)U,
  \qquad
  W=0.
  \]

  \item \textbf{\texttt{4pn\_local\_referee\_master\_sympy\_audit.py}.}
  \emph{Role:} end-to-end referee reconstruction of the full local ordinary
  target.  \emph{Reference summary:} the assembled generic-frame local ordinary
  candidate reproduces the fixed-chart local ordinary target slot by slot in the
  $Q$, $T$, $S$, $U$, and $W$ blocks, so the entire local instantaneous $4$PN
  sector is closed exactly inside the declared hierarchy.

  \item \textbf{\texttt{4pn\_tail\_bridge\_audit.py}.}  \emph{Role:}
  universal Newtonian quadrupole shadow of the conservative GR tail.
  \emph{Reference summary:} the order-reduced STF quadrupole shadow is confined
  to the degree-$2$ $G^4/r^4$ block, so the hereditary problem is structurally
  much smaller than the local generic-frame lift.

  \item \textbf{\texttt{4pn\_tail\_hereditary\_bridge\_audit.py}.}
  \emph{Role:} exact coefficient bridge from the normalized STF quadrupole to
  the hereditary $4$PN coefficient.  \emph{Reference summary:} the audit proves
  \[
  C_{\rm tail}=\frac{GM}{2c^3}\,\gamma_{\rm quad}^{\rm eff},
  \]
  so the hereditary coefficient introduces no new independent normalization
  datum.

  \item \textbf{\texttt{4pn\_conditional\_referee\_master\_audit.py}.}
  \emph{Role:} conditional full conservative $4$PN replay.  \emph{Reference
  summary:} the final replay ties the inherited $2.5$PN quadrupole narrowing to
  the exact local $4$PN closure and the hereditary bridge, and reduces the full
  remaining theorem gate to
  \[
  \gamma_{\rm quad}^{\rm eff}\stackrel{?}{=}\frac{2G}{5c^5},
  \qquad
  \hat m_0^{\,2}P_0\stackrel{?}{=}\frac{54Gc_s^5}{5a_{\rm th}^5c^5}.
  \]
\end{enumerate}

\subsection{Reference master reruns}
\label{app:verification-reference-runs}

For a referee-facing rerun, the shortest load-bearing replay is the following.

\begin{enumerate}
  \item Run the carried \texttt{2\_5pn\_master\_session\_sympy\_audit.py}
  archive and verify that the universal higher-order bridge is still narrowed to
  the canonically normalized orbital / worldtube STF quadrupole branch, with one
  passive/outgoing normalization gap left open.

  \item Run \texttt{4pn\_local\_referee\_master\_sympy\_audit.py} and verify
  that the full fixed-chart local ordinary target is reproduced exactly by one
  generic-frame local ordinary representative assembled from the natural seed,
  the aligned-seed lift, and the canonical comparable-mass residual.

  \item Run \texttt{4pn\_tail\_hereditary\_bridge\_audit.py} and verify the
  bridge relation
  \[
  C_{\rm tail}=\frac{GM}{2c^3}\,\gamma_{\rm quad}^{\rm eff}.
  \]

  \item Run \texttt{4pn\_conditional\_referee\_master\_audit.py} and verify
  that the only remaining global gate is the same passive/outgoing quadrupole
  normalization already inherited from the $2.5$PN package.
\end{enumerate}

These are the publication reference runs because they reproduce the paper's
sharpest theorem-status statement in the smallest number of reruns.  More local
reruns are still useful when one wants to inspect a specific stage in detail,
but the four-step replay above is the minimal end-to-end referee path.

For a second-engine replay, run the seven Mathematica scripts listed in
\nolinkurl{mathematica/wl_notes.txt} sequentially with \nolinkurl{math -script}.
Those grouped audits mirror the same theorem chain at coarser granularity and
carry embedded \texttt{Output:} blocks so that the publication archive contains
their live reference transcripts as well as the stagewise SymPy records.

\subsection{What the archive verifies}
\label{app:verification-what-verifies}

The archive verifies the load-bearing symbolic statements of the paper sector by
sector.

First, the front-end audits verify the exact one-body and compiler backbone.
They check the strict Schwarzschild $4$PN target, confirm that the carried $3$PN
self-sector packaging is insufficient by itself, solve the minimal repair
ledger, and freeze the exact quartic perturbative Legendre-transform identity.

Second, the local middle audits verify the entire Hamiltonian-first comparable-
mass closure chain.  They check the raw local scaffold and COM completeness, the
fixed local ADM target import, the chart-level degree-ceiling lesson, the
Hamiltonian-chart generic-frame lift, the canonical residual slice, the exact
ordinary sign flip, and the structured liftability of the aligned-seed
correction.  In the language of the main text, this is the stage chain that
justifies the statement that the whole local instantaneous $4$PN sector has an
exact generic-frame ordinary representative inside the declared hierarchy.

Third, the local referee master verifies the end-to-end local assembly itself.
It rebuilds the reduced COM local ordinary target from the imported local ADM
Hamiltonian target, verifies the one-body gate and natural seed, replays the
canonical residual translation and aligned-seed correction, and confirms that
the assembled generic-frame candidate reduces back to the full fixed-chart local
ordinary target slot by slot.

Fourth, the tail-side audits verify the narrow hereditary bridge.  They confirm
that the GR tail contributes a universal Newtonian quadrupole shadow in the
local logarithmic block and that, on the canonically normalized STF quadrupole
branch, the hereditary coefficient is exactly
\[
C_{\rm tail}=\frac{GM}{2c^3}\,\gamma_{\rm quad}^{\rm eff}.
\]
So the hereditary side is not an independent $4$PN coefficient hunt.  It is an
interface theorem tying the conservative hereditary coefficient to the same
quadrupole normalization that already controls the $2.5$PN odd sector.

Finally, the conditional master replay verifies the paper's sharpest global
status statement: once the inherited $2.5$PN narrowing, the exact local $4$PN
closure, and the hereditary bridge are replayed together, the only remaining
full-$4$PN gate is the already-known passive/outgoing quadrupole normalization.
In other words, the executable archive verifies the same conclusion as the
analytic paper: no new independent normalization datum opens at $4$PN.

\subsection{What the archive does not verify}
\label{app:verification-what-not-verify}

The archive does \emph{not} replace the analytic argument in the main text, and
it does not solve the full moving-throat PDE.  In particular, it does not by
itself derive the passive/outgoing STF quadrupole normalization on the actual
moving-throat branch; it does not prove from first principles that
\[
\gamma_{\rm quad}^{\rm eff}=\frac{2G}{5c^5}
\qquad\text{or}\qquad
\hat m_0^{\,2}P_0=\frac{54Gc_s^5}{5a_{\rm th}^5c^5}
\]
holds; and it does not upgrade the present work into an assumption-free theorem
of the full bulk dynamics.

It also does not attempt to settle every deeper microscopic interpretation issue
that remains outside the present $4$PN scope.  The archive is not a proof of
the full moving-throat origin of every grouped conservative datum, it is not a
final uniqueness theorem for every aligned-seed local representative, and it is
not a general radiative or higher-PN completion package.  Those questions remain
outside the stated theorem target of the paper.

So the correct reading rule is simple.  The referee archive verifies that the
paper's exact local theorem chain and hereditary bridge close where claimed, and
that the whole conservative $4$PN problem has been reduced to one already-known
quadrupole-normalization gate.  What it does not verify is the final PDE-side
realization of that gate.  That distinction is exactly the one on which the
conditional status of the present $4$PN theorem rests.

\end{document}
